% Reboisement et écotourisme — Allemand 1ère A — Cours signé OnBuch+
% Programme officiel MINESEC. Compiler avec : tectonic ALA10-reboisement-et-ecotourisme.tex
\newcommand{\DOCMATIERE}{Allemand}
\newcommand{\DOCNIVEAU}{1ère A}
\newcommand{\DOCMODULE}{Chapitre 3 : Environnement, santé et bien-être}
\newcommand{\DOCLECON}{ALA10}
\newcommand{\DOCTITRE}{Reboisement et écotourisme}
% OnBuch+ — préambule « cours signés » ENRICHI (compilé avec tectonic / XeLaTeX)
% Système visuel « Pop » d'OnBuch+ : fond crème (jamais blanc), bordures encre
% épaisses, ombres « dures » décalées, titres Archivo Black en capitales.
% Version MATHÉMATIQUES : graphes (pgfplots/tikz), boîtes pédagogiques
% (définition, propriété, méthode, attention, le savais-tu, exercice résolu,
% exercice, TP, bilan), page de garde et signature OnBuch+ ancrée en bas.
%
% Macros attendues dans main.tex AVANT \input{preamble} :
%   \DOCMATIERE  \DOCNIVEAU  \DOCMODULE  \DOCLECON (numéro)  \DOCTITRE
\documentclass[11pt]{article}

\usepackage{fontspec}
\usepackage[ngerman,french]{babel} % français = langue principale ; l'allemand via \all{..} / textebox
\usepackage[a4paper,margin=2cm,top=3.4cm,bottom=2.8cm,headheight=1.4cm,footskip=1.3cm]{geometry}
\usepackage{amsmath,amssymb,amsfonts}
\usepackage{mathtools} % \xmapsto, \coloneqq... souvent utilisés par les modèles
\usepackage{cancel} % simplifications barrées (bilans de forces)
% \abs{z} (module, valeur absolue) et \norm{u} (norme), utilisés par les modèles
\providecommand{\abs}{}\let\abs\relax\DeclarePairedDelimiter{\abs}{\lvert}{\rvert}
\providecommand{\norm}{}\let\norm\relax\DeclarePairedDelimiter{\norm}{\lVert}{\rVert}
\usepackage[table]{xcolor} % [table] : fournit aussi \rowcolors (alternance de lignes)
\usepackage{pagecolor}
\usepackage{tikz}
\usetikzlibrary{arrows.meta,positioning,calc,shapes.geometric,shapes.misc,shapes.arrows,decorations.pathmorphing,decorations.pathreplacing,decorations.markings,patterns,shadows,mindmap,trees,fit,backgrounds,matrix,intersections,angles,quotes}
\usepackage{pgfplots}
\usepackage[version=4]{mhchem} % équations nucléaires \ce{^{235}_{92}U}
\usepackage[european,siunitx]{circuitikz} % schémas électriques
\pgfplotsset{compat=1.18}
\usepgfplotslibrary{fillbetween,groupplots} % aires entre courbes (calcul intégral)
\usepackage{enumitem}
\usepackage{fancyhdr}
% [most] charge toutes les bibliothèques tcolorbox (listings, minted, raster,
% documentation...) et gonfle inutilement la mémoire TeX sur un document long
% (~300 boîtes) : on ne charge que ce qui est réellement utilisé.
\usepackage[skins,breakable]{tcolorbox}
\usepackage{graphicx}
\usepackage{colortbl}
\usepackage{booktabs}
\usepackage{tabularx}
\usepackage{diagbox} % case d'en-tête diagonale des tableaux à double entrée
% Colonnes extensibles centrée / gauche / droite (C, L, R), souvent utilisées par les modèles
\newcolumntype{C}{>{\centering\arraybackslash}X}
\newcolumntype{L}{>{\raggedright\arraybackslash}X}
\newcolumntype{R}{>{\raggedleft\arraybackslash}X}
\usepackage{array}
\usepackage{multicol}
\usepackage{multirow}
\usepackage{siunitx}
% parse-numbers=false : accepte aussi \SI{2,5 \times 10^{-3}}{...}
\sisetup{locale=FR,output-decimal-marker={,},group-minimum-digits=4,parse-numbers=false}
\DeclareSIUnit\ohm{\ensuremath{\symup{\Omega}}} % U+2126 absent de la police
\DeclareSIUnit\division{div} % réglages d'oscilloscope (ms/div, V/div)
\DeclareSIUnit\tr{tr} % tours (vitesse de rotation, ex. \SI{1500}{\tr\per\minute})
\DeclareSIUnit\molar{M} % concentration molaire (ex. \SI{3}{\molar}, \SI{1}{\micro\molar})
\DeclareSIUnit\an{an} % année (le modèle confond parfois avec \year, un compteur TeX) — ex. \SI{5}{\kilo\meter\per\an}
\DeclareSIUnit\FCFA{FCFA} % franc CFA (ex. \SI{500}{\FCFA\per\kilo\gram})
\DeclareSIUnit\degreeBrix{\textdegree Brix} % teneur en sucre d'un sirop/jus (réfractométrie)
\DeclareSIUnit\month{mois} % le modèle utilise parfois \month, un compteur TeX, comme unité de durée
\usepackage{qrcode}
% \degree : commande fréquemment utilisée par les modèles pour un angle ou une
% température en dehors de \SI{}{} (non fournie par les paquets chargés).
\providecommand{\degree}{^{\circ}}
% \qed (fin de démonstration), utilisé par les modèles sans amsthm
\providecommand{\qed}{\hfill$\square$}
% \barre : double barre des valeurs interdites dans un tableau de variations (tkz-tab)
\providecommand{\barre}{\ensuremath{\|}}
% environnement proof (amsthm non chargé) utilisé par les modèles
\ifdefined\proof\else\newenvironment{proof}[1][Démonstration]{\par\noindent\textit{#1.}\ }{\qed\par}\fi
\usepackage{lastpage}
\usepackage{hyperref}
\hypersetup{hidelinks}

% ============================================================
% Palette « Pop » OnBuch+ — mêmes hex que la classe P (plus_ui.dart)
% ============================================================
\definecolor{poporange}{HTML}{F59321}
\definecolor{popdark}{HTML}{E07A0C}
\definecolor{popcream}{HTML}{FFF6E8}
\definecolor{popink}{HTML}{1C1714}
\definecolor{poppurple}{HTML}{7A5AE0}
\definecolor{popgreen}{HTML}{2E9E6A}
\definecolor{poppink}{HTML}{E0457B}
\definecolor{popblue}{HTML}{2F7DE1}
\definecolor{popgold}{HTML}{C98A12}
\definecolor{popmuted}{HTML}{8A7F76}
\definecolor{popline}{HTML}{EADFD2}
\definecolor{popgray}{HTML}{8A7F76}
\colorlet{popgrey}{popgray}
% Alias tolérants (le modèle nomme parfois les couleurs différemment)
\colorlet{popwhite}{white}
\colorlet{popblack}{popink}
\colorlet{popred}{poppink}
\colorlet{popyellow}{popgold}
% Teintes claires (fonds de tableaux/encadrés) — le modèle nomme parfois
% les couleurs avec un suffixe L (ex. popblueL) sans les avoir définies.
\colorlet{poporangeL}{poporange!20}
\colorlet{popdarkL}{popdark!20}
\colorlet{poppurpleL}{poppurple!20}
\colorlet{popgreenL}{popgreen!20}
\colorlet{poppinkL}{poppink!20}
\colorlet{popblueL}{popblue!20}
\colorlet{popgoldL}{popgold!20}
\colorlet{popredL}{popred!20}
\colorlet{popgrayL}{popgray!20}
% Teintes claires pour fonds de tableaux / zones
\colorlet{popblueL}{popblue!10!white}
\colorlet{poporangeL}{poporange!14!white}
\colorlet{popgreenL}{popgreen!12!white}
\colorlet{poppurpleL}{poppurple!12!white}
\colorlet{poppinkL}{poppink!12!white}
\colorlet{popgoldL}{popgold!14!white}

\pagecolor{popcream}

% ============================================================
% Typographie
% ============================================================
\defaultfontfeatures{Ligatures=TeX}
\setmainfont{PlusJakartaSans-Medium.ttf}[Path=../../../fonts/,BoldFont=PlusJakartaSans-Bold.ttf]
% Maths en sans-serif (Fira Math) avec chiffres et lettres droites de Plus
% Jakarta, pour que les formules se fondent dans le texte.
\usepackage[math-style=ISO,bold-style=ISO]{unicode-math}
\setmathfont{FiraMath-Regular.otf}[Path=../../../fonts/]
\setmathfont{PlusJakartaSans-Medium.ttf}[Path=../../../fonts/,range={up/num,up/{Latin,latin},\mathup}]
\usepackage{icomma} % virgule décimale sans espace en mode math
% Caractères mathématiques Unicode absents des polices de texte (Plus Jakarta,
% Archivo Black) : rendus par leur équivalent mathématique, en texte comme en
% formule — sinon ils s'affichent en carré vide (ex. « Arithmétique dans ℤ »).
\usepackage{newunicodechar}
\newunicodechar{ℝ}{\ensuremath{\mathbb{R}}}
\newunicodechar{ℤ}{\ensuremath{\mathbb{Z}}}
\newunicodechar{ℂ}{\ensuremath{\mathbb{C}}}
\newunicodechar{ℕ}{\ensuremath{\mathbb{N}}}
\newunicodechar{ℚ}{\ensuremath{\mathbb{Q}}}
\newunicodechar{𝔻}{\ensuremath{\mathbb{D}}}
\newunicodechar{α}{\ensuremath{\alpha}}
\newunicodechar{β}{\ensuremath{\beta}}
\newunicodechar{γ}{\ensuremath{\gamma}}
\newunicodechar{δ}{\ensuremath{\delta}}
\newunicodechar{ε}{\ensuremath{\varepsilon}}
\newunicodechar{θ}{\ensuremath{\theta}}
\newunicodechar{λ}{\ensuremath{\lambda}}
\newunicodechar{μ}{\ensuremath{\mu}}
\newunicodechar{σ}{\ensuremath{\sigma}}
\newunicodechar{τ}{\ensuremath{\tau}}
\newunicodechar{φ}{\ensuremath{\varphi}}
\newunicodechar{ψ}{\ensuremath{\psi}}
\newunicodechar{ω}{\ensuremath{\omega}}
\newunicodechar{Σ}{\ensuremath{\Sigma}}
% majuscules produites par \MakeUppercase dans les titres (α-aminés -> Α)
\newunicodechar{Α}{\ensuremath{\alpha}}
\newunicodechar{Β}{\ensuremath{\beta}}
\newunicodechar{Γ}{\ensuremath{\Gamma}}
\newunicodechar{Δ}{\ensuremath{\Delta}}
\newunicodechar{Ε}{\ensuremath{\varepsilon}}
\newunicodechar{Θ}{\ensuremath{\Theta}}
\newunicodechar{Λ}{\ensuremath{\Lambda}}
\newunicodechar{Μ}{\ensuremath{\mu}}
\newunicodechar{Τ}{\ensuremath{\tau}}
\newunicodechar{Φ}{\ensuremath{\Phi}}
\newunicodechar{Ψ}{\ensuremath{\Psi}}
\newunicodechar{Ω}{\ensuremath{\Omega}}
\newunicodechar{↦}{\ensuremath{\mapsto}}
\newunicodechar{∈}{\ensuremath{\in}}
\newunicodechar{∉}{\ensuremath{\notin}}
\newunicodechar{∧}{\ensuremath{\wedge}}
\newunicodechar{⇔}{\ensuremath{\Leftrightarrow}}
\newunicodechar{⟺}{\ensuremath{\Longleftrightarrow}}
\newunicodechar{⇒}{\ensuremath{\Rightarrow}}
\newunicodechar{⟹}{\ensuremath{\Longrightarrow}}
\newunicodechar{∘}{\ensuremath{\circ}}
\newunicodechar{∪}{\ensuremath{\cup}}
\newunicodechar{∩}{\ensuremath{\cap}}
\newunicodechar{≡}{\ensuremath{\equiv}}
\newunicodechar{⊂}{\ensuremath{\subset}}
\newunicodechar{⊥}{\ensuremath{\perp}}
\newunicodechar{∃}{\ensuremath{\exists}}
\newunicodechar{∀}{\ensuremath{\forall}}
\newunicodechar{∛}{\ensuremath{\sqrt[3]{\,}}}
\newunicodechar{∜}{\ensuremath{\sqrt[4]{\,}}}
\newunicodechar{⃗}{\ensuremath{{}^{\to}}}
\newunicodechar{ⁿ}{\ifmmode^{n}\else\textsuperscript{\ensuremath{n}}\fi}
\newunicodechar{ˣ}{\ifmmode^{x}\else\textsuperscript{\ensuremath{x}}\fi}
\newunicodechar{ᵇ}{\ifmmode^{b}\else\textsuperscript{\ensuremath{b}}\fi}
\newunicodechar{ʳ}{\ifmmode^{r}\else\textsuperscript{\ensuremath{r}}\fi}
\newunicodechar{ᵘ}{\ifmmode^{u}\else\textsuperscript{\ensuremath{u}}\fi}
\newunicodechar{ʸ}{\ifmmode^{y}\else\textsuperscript{\ensuremath{y}}\fi}
\newunicodechar{ᵃ}{\ifmmode^{a}\else\textsuperscript{\ensuremath{a}}\fi}
\newunicodechar{ᵉ}{\ifmmode^{e}\else\textsuperscript{\ensuremath{e}}\fi}
\newunicodechar{ᵏ}{\ifmmode^{k}\else\textsuperscript{\ensuremath{k}}\fi}
\newunicodechar{ᵖ}{\ifmmode^{p}\else\textsuperscript{\ensuremath{p}}\fi}
\newunicodechar{ᴺ}{\ifmmode^{N}\else\textsuperscript{\ensuremath{N}}\fi}
\newunicodechar{ᶜ}{\ifmmode^{c}\else\textsuperscript{\ensuremath{c}}\fi}
\newunicodechar{ʰ}{\ifmmode^{h}\else\textsuperscript{\ensuremath{h}}\fi}
\newunicodechar{ᵠ}{\ifmmode^{q}\else\textsuperscript{\ensuremath{q}}\fi}
\newunicodechar{ₙ}{\ifmmode_{n}\else\textsubscript{\ensuremath{n}}\fi}
\newunicodechar{ₐ}{\ifmmode_{a}\else\textsubscript{\ensuremath{a}}\fi}
\newunicodechar{ᵢ}{\ifmmode_{i}\else\textsubscript{\ensuremath{i}}\fi}
\newunicodechar{ᵣ}{\ifmmode_{r}\else\textsubscript{\ensuremath{r}}\fi}
\newunicodechar{ₖ}{\ifmmode_{k}\else\textsubscript{\ensuremath{k}}\fi}
\newunicodechar{ᵦ}{\ifmmode_{\beta}\else\textsubscript{\ensuremath{\beta}}\fi}
\newunicodechar{ₓ}{\ifmmode_{x}\else\textsubscript{\ensuremath{x}}\fi}
\newfontfamily\popdisplay{ArchivoBlack-Regular.ttf}[Path=../../../fonts/]
\newfontfamily\poplogo{SpaceGrotesk-Bold.ttf}[Path=../../../fonts/]
\newfontfamily\popsemi{PlusJakartaSans-SemiBold.ttf}[Path=../../../fonts/]
\newfontfamily\popxbold{PlusJakartaSans-ExtraBold.ttf}[Path=../../../fonts/]
\newfontfamily\popxboldls{PlusJakartaSans-ExtraBold.ttf}[Path=../../../fonts/,LetterSpace=3.0]

\setlist[enumerate]{leftmargin=1.7em,itemsep=5pt,topsep=4pt}
\setlist[itemize]{leftmargin=1.7em,itemsep=4pt,topsep=4pt,label={\color{poporange}\textbullet}}
\setlength{\parindent}{0pt}
\setlength{\parskip}{5pt}
\renewcommand{\arraystretch}{1.25}

% pgfplots : style Pop par défaut
\pgfplotsset{
  every axis/.append style={axis line style={popink,line width=1pt},tick style={popink},
    grid style={popline,line width=0.5pt},label style={font=\small},tick label style={font=\footnotesize}},
  popcurve/.style={poporange,line width=1.8pt,no markers,smooth},
}
% Même style disponible dans un \draw TikZ simple (hors pgfplots)
\tikzset{popcurve/.style={poporange,line width=1.8pt}}
\tikzset{popgrid/.style={popline,very thin}}
\colorlet{popcurve}{poporange} % le nom du style est parfois utilisé comme couleur

% ============================================================
% Pill / squiggle
% ============================================================
\newcommand{\poppill}[3]{%
  \tikz[baseline=-0.55ex]\node[draw=popink,line width=1.2pt,rounded corners=2.6pt,%
    fill=#2,inner xsep=6.5pt,inner ysep=3pt]{%
    \popxboldls\fontsize{7}{7}\selectfont\color{#3}\MakeUppercase{#1}};%
}
\newcommand{\popsquiggle}[1]{%
  \tikz[baseline=-0.4ex]{
    \draw[#1,line width=2.6pt,line cap=round]
      plot [smooth] coordinates {(0,0.09) (0.34,0.01) (0.68,0.21) (1.02,0.01) (1.36,0.10)};
  }%
}
% Mot-clé surligné (marqueur orange sous le texte)
\newcommand{\cle}[1]{{\popxbold\color{popdark}#1}}

% ============================================================
% En-tête / pied
% ============================================================
\pagestyle{fancy}
\fancyhf{}
\renewcommand{\headrulewidth}{0pt}
\renewcommand{\footrulewidth}{0pt}
\lhead{\raisebox{-0.15ex}{\poplogo\fontsize{13}{13}\selectfont\color{popink}OnBuch\color{poporange}+}}
\rhead{\poppill{\DOCMATIERE\ \textbullet\ \DOCNIVEAU\ \textbullet\ Leçon \DOCLECON}{white}{popink}}
\lfoot{\popxbold\fontsize{7.6}{7.6}\selectfont\color{popmuted}\MakeUppercase{Cours signé OnBuch+}\ \textbullet\ {\popsemi\DOCTITRE}}
\rfoot{\tikz[baseline=-0.6ex]\node[rounded corners=8.5pt,fill=popink,minimum height=17pt,minimum width=17pt,inner xsep=5pt,inner sep=0]{\popxbold\fontsize{8}{8}\selectfont\color{popcream}\thepage\,/\,\pageref*{LastPage}};}

% ============================================================
% Titres
% ============================================================
\newcommand{\chaptitre}[2]{%
  \begin{tcolorbox}[enhanced,colback=poporange,colframe=popink,boxrule=2.6pt,arc=11pt,
      drop shadow=popink,left=15pt,right=15pt,top=12pt,bottom=11pt,before skip=3pt,after skip=18pt]
    \poppill{#2}{popink}{popcream}\par\vspace{9pt}
    {\raggedright\hyphenpenalty=10000\exhyphenpenalty=10000\popdisplay\fontsize{22}{24}\selectfont\color{popcream}\MakeUppercase{#1}\par}
  \end{tcolorbox}
}
% Section numérotée : pastille orange avec le numéro + titre Archivo
\newcounter{popsec}
\newcommand{\coursec}[1]{%
  \stepcounter{popsec}\setcounter{popsub}{0}%
  \par\vspace{16pt}\noindent
  \begin{minipage}{\linewidth}
  \tikz[baseline=-0.7ex]\node[draw=popink,line width=1.6pt,fill=poporange,rounded corners=4pt,minimum size=22pt,inner sep=0]{\popdisplay\fontsize{12}{12}\selectfont\color{popcream}\thepopsec};%
  \hspace{8pt}{\popdisplay\fontsize{15}{17}\selectfont\color{popink}\MakeUppercase{#1}}\par
  \vspace{4pt}\noindent{\color{poporange}\rule{36pt}{3.6pt}}
  \end{minipage}\par\nopagebreak\vspace{8pt}
}
\newcounter{popsub}
\newcommand{\courssub}[1]{%
  \stepcounter{popsub}%
  \par\vspace{9pt}\noindent{\popxbold\fontsize{11.5}{13}\selectfont\color{popink}{\color{poporange}\thepopsec.\thepopsub}\hspace{6pt}#1}\par\nopagebreak\vspace{4pt}
}

% ============================================================
% Boîtes pédagogiques — même recette : fond blanc, bordure épaisse colorée,
% ombre dure encre, badge plein. Titre optionnel [#1].
% ============================================================
\tcbset{popbase/.style={breakable,enhanced,colback=white,boxrule=2.3pt,arc=9pt,
  shadow={3pt}{-3pt}{0pt}{popink},left=14pt,right=14pt,top=11pt,bottom=11pt,before skip=11pt,after skip=15pt,
  fontupper=\popsemi\fontsize{10.6}{14.5}\selectfont\color{popink},
  fontlower=\popsemi\fontsize{10.6}{14.5}\selectfont\color{popink}}}
% Coche dessinée (le glyphe ✓ n'existe pas dans les polices du document)
\newcommand{\popcheck}[1]{\tikz[baseline=-0.5ex]\draw[#1,line width=1.6pt,line cap=round,line join=round](-0.13,0.0)--(-0.04,-0.09)--(0.13,0.1);}
\providecommand{\coche}{\popcheck{popgreen}}
\providecommand{\resultat}[1]{\par\smallskip\noindent\fcolorbox{popgreen}{popgreenL}{\parbox{\dimexpr\linewidth-2\fboxsep-2\fboxrule\relax}{\textbf{Résultat.} #1}}\par\smallskip}
\newcommand{\popboxhead}[3]{\poppill{#1}{#2}{white}\ifx\relax#3\relax\else\hspace{6pt}{\popxbold\fontsize{10.6}{12}\selectfont\color{popink}#3}\fi\par\vspace{7pt}}

\newtcolorbox{aretenir}[1][]{popbase,colframe=popgreen,before upper={\popboxhead{À retenir}{popgreen}{#1}}}
% Texte en allemand (document de lecture, dialogue, extrait) : cadre
% d'étude. Un titre optionnel (source, auteur) s'affiche dans l'en-tête.
\newtcolorbox{textebox}[1][]{popbase,colframe=popblue,before upper={\popboxhead{Text}{popblue}{#1}\begin{otherlanguage}{ngerman}},after upper={\end{otherlanguage}}}
% Marques ✓ / ✗ (forme correcte / fautive) souvent utilisées par les modèles
\providecommand{\cross}{\textcolor{poppink}{\ensuremath{\times}}}
\providecommand{\xmark}{\cross}\providecommand{\crossmark}{\cross}\providecommand{\cmark}{\ensuremath{\checkmark}}
% Ligne de réponse / justification à compléter (inventée par les modèles)
\providecommand{\justification}[1]{\par\noindent\textit{Justification~:} \dotfill\par}
% \textipa (package tipa non chargé) : la police ne couvre pas l'API -> simple passage
\providecommand{\textipa}[1]{#1}
% Soulignement (ulem non chargé) : \uline, \ul
\providecommand{\uline}[1]{\underline{#1}}\providecommand{\ul}[1]{\underline{#1}}
% \alert (beamer) inventée par les modèles : texte mis en évidence
\providecommand{\alert}[1]{\textbf{\textcolor{poppink}{#1}}}
% variantes ASCII d'ordinaux écrites en macro
\providecommand{\iere}{\textsuperscript{re}}\providecommand{\ieme}{\textsuperscript{e}}\providecommand{\ere}{\textsuperscript{re}}\providecommand{\eme}{\textsuperscript{e}}\providecommand{\er}{\textsuperscript{er}}\providecommand{\ie}{\textsuperscript{e}}
% Ordinaux français écrits en macro par les modèles : 1\ière, 3\ième
\newcommand{\ière}{\textsuperscript{re}}\newcommand{\ième}{\textsuperscript{e}}
% \exemple (inventée par les modèles) : étiquette « Exemple : »
\providecommand{\exemple}{\textbf{Exemple~:}\ }
% \square utilisable aussi hors mode math (cases à cocher)
\AtBeginDocument{\let\squareMath\square\renewcommand{\square}{\ensuremath{\squareMath}}}
% Mot ou phrase en allemand à l'intérieur d'un paragraphe français
\newcommand{\all}[1]{\ifmmode\text{\textit{#1}}\else\begingroup\selectlanguage{ngerman}\itshape #1\endgroup\fi}
\let\esp\all % alias toléré
\newtcolorbox{exemplebox}[1][]{popbase,colframe=poporange,before upper={\popboxhead{Exemple}{poporange}{#1}}}
\newtcolorbox{definition}[1][]{popbase,colframe=popblue,before upper={\popboxhead{Définition}{popblue}{#1}}}
\newtcolorbox{propriete}[1][]{popbase,colframe=poppurple,before upper={\popboxhead{Propriété}{poppurple}{#1}}}
\newtcolorbox{methode}[1][]{popbase,colframe=popgold,before upper={\popboxhead{Méthode}{popgold}{#1}}}
\newtcolorbox{attention}[1][]{popbase,colframe=poppink,before upper={\popboxhead{Attention}{poppink}{#1}}}
\newtcolorbox{experience}[1][]{popbase,colframe=popblue,colback=popblueL,before upper={\popboxhead{Expérience}{popblue}{#1}}}
% « Le savais-tu ? » : carte violette pleine (contexte, culture, Cameroun)
\newtcolorbox{savaistu}[1][]{popbase,colback=poppurple,colframe=popink,
  fontupper=\popsemi\fontsize{10.4}{14.2}\selectfont\color{white},
  before upper={\poppill{Le savais-tu\,?}{white}{poppurple}\ifx\relax#1\relax\else\hspace{6pt}{\popxbold\color{white}#1}\fi\par\vspace{7pt}}}
% Exercice résolu : énoncé en haut, \tcblower puis la solution sur fond crème
\newcounter{popexo}\newcounter{popexr}
\newtcolorbox{exoresolu}[1][]{popbase,colframe=poporange,colbacklower=popcream,
  segmentation style={popink,line width=1.2pt,dashed},
  before upper={\stepcounter{popexr}\popboxhead{Exercice résolu \thepopexr}{poporange}{#1}},
  before lower={\poppill{Solution}{popink}{popcream}\par\vspace{6pt}}}
% Exercice (non résolu en place) — la correction est dans la partie Corrigés
\newtcolorbox{exercice}[1][]{popbase,colframe=popink,
  before upper={\stepcounter{popexo}\popboxhead{Exercice \thepopexo}{popink}{#1}}}
% Corrigé
\newtcolorbox{corrige}[1][]{popbase,colframe=popgreen,colback=popgreenL,
  before upper={\popboxhead{Corrigé}{popgreen}{#1}}}
% Objectifs / prérequis / bilan
\newtcolorbox{objectifs}{popbase,colframe=popink,colback=poporangeL,
  before upper={\popboxhead{Objectifs de la leçon}{popink}{}}}
\newtcolorbox{prerequis}{popbase,colframe=popink,colback=popblueL,
  before upper={\popboxhead{Prérequis}{popblue}{}}}
\newtcolorbox{bilan}{popbase,colframe=popink,colback=white,boxrule=2.8pt,
  before upper={\popboxhead{Fiche bilan}{poporange}{L'essentiel à connaître pour l'examen}}}
% Figure encadrée avec légende
\newtcolorbox{popfigure}[1][]{enhanced,breakable=false,colback=white,colframe=popline,boxrule=1.4pt,arc=8pt,
  left=10pt,right=10pt,top=10pt,bottom=8pt,before skip=10pt,after skip=12pt,halign=center,
  lower separated=false,halign lower=center,
  fontlower=\popsemi\fontsize{9}{11.5}\selectfont\color{popmuted},
  #1}
\newcommand{\legende}[1]{\par\vspace{6pt}{\popsemi\fontsize{9}{11.5}\selectfont\color{popmuted}#1}}

% ============================================================
% Signature OnBuch+
% ============================================================
\newcommand{\poplogoinline}[1]{{\poplogo\fontsize{#1}{#1}\selectfont\color{popink}OnBuch\color{poporange}+}}
\newcommand{\signatureonbuch}{%
  \begin{tcolorbox}[enhanced,colback=popink,colframe=popink,boxrule=2.2pt,arc=10pt,
      drop shadow=poporange,left=14pt,right=14pt,top=10pt,bottom=10pt,before skip=0pt,after skip=0pt]
    \tikz[baseline=-0.7ex]\node[circle,fill=popgreen,draw=popcream,line width=1.3pt,minimum size=22pt,inner sep=0]{\popcheck{white}};%
    \hspace{9pt}\begin{minipage}[c]{0.62\linewidth}
      {\popxbold\fontsize{12}{13}\selectfont\color{popcream}Signé\ \textbullet\ {\poplogo OnBuch\color{poporange}+}}\par\vspace{2pt}
      {\raggedright\popsemi\fontsize{8}{10.5}\selectfont\color{popline}\MakeUppercase{Équipe pédagogique OnBuch+}\\\MakeUppercase{Conforme au programme officiel MINESEC}\par}
    \end{minipage}\hfill
    {\poplogo\fontsize{22}{22}\selectfont\color{popcream}OnBuch\color{poporange}+}
  \end{tcolorbox}%
}

% ============================================================
% Page de garde
% #1 titre · #2 matière · #3 niveau · #4 module · #5 n° leçon · #6 durée ·
% #7 description / objectifs (texte ou itemize)
% ============================================================
\newcommand{\pagegarde}[7]{%
  \thispagestyle{empty}
  \newgeometry{a4paper,margin=1.7cm,top=1.6cm,bottom=1.4cm}%
  \begin{tikzpicture}[remember picture,overlay]
    \fill[poporange!18] ([xshift=-3.2cm,yshift=-3.6cm]current page.north east) circle (4.6cm);
    \fill[poppurple!14] ([xshift=2.4cm,yshift=7.6cm]current page.south west) circle (3.8cm);
  \end{tikzpicture}%
  {\poplogo\fontsize{30}{30}\selectfont\color{popink}OnBuch\color{poporange}+}\hfill
  \raisebox{0.6ex}{\poppill{Cours signé \textbullet\ Premium}{popink}{popcream}}\par
  \vspace{1pt}\popsquiggle{poppurple}\par
  \vspace{8pt}
  \begin{tcolorbox}[enhanced,colback=poporange,colframe=popink,boxrule=2.8pt,arc=13pt,
      drop shadow=popink,left=18pt,right=18pt,top=12pt,bottom=13pt,after skip=10pt]
    \poppill{#2 \textbullet\ #3}{popink}{popcream}\hspace{5pt}\poppill{#4}{white}{popink}\par\vspace{8pt}
    {\popdisplay\fontsize{14}{14}\selectfont\color{popink}LEÇON #5}\par\vspace{4pt}
    {\raggedright\hyphenpenalty=10000\exhyphenpenalty=10000\popdisplay\fontsize{25}{27}\selectfont\color{popcream}\MakeUppercase{#1}\par}
    \vspace{8pt}
    {\popxbold\fontsize{9.5}{9.5}\selectfont\color{popink}\MakeUppercase{Durée conseillée : #6}}
  \end{tcolorbox}
  \begin{tcolorbox}[enhanced,colback=white,colframe=popink,boxrule=2.2pt,arc=10pt,
      drop shadow=popink,left=16pt,right=16pt,top=10pt,bottom=10pt,after skip=8pt]
    \poppill{Dans cette leçon}{poppurple}{white}\par\vspace{5pt}
    \popsemi\fontsize{9.3}{12.6}\selectfont\color{popink}#7
  \end{tcolorbox}
  \noindent\begin{minipage}[t]{0.487\linewidth}
    \begin{tcolorbox}[enhanced,colback=popgreen,colframe=popink,boxrule=2.2pt,arc=10pt,
        drop shadow=popink,left=11pt,right=11pt,top=10pt,bottom=10pt,height=3.1cm,valign=center]
      \tcbox[colback=white,colframe=popink,boxrule=1.3pt,arc=5pt,boxsep=0pt,nobeforeafter,
        left=3pt,right=3pt,top=3pt,bottom=3pt,valign=center]{\qrcode[height=1.9cm]{https://play.google.com/store/apps/details?id=com.luvvix.onbuch&hl=fr}}%
      \hspace{8pt}\begin{minipage}[c]{0.5\linewidth}
        {\popdisplay\fontsize{11}{12.5}\selectfont\color{white}\MakeUppercase{Télécharge OnBuch}}\par\vspace{4pt}
        {\raggedright\popsemi\fontsize{8.4}{11}\selectfont\color{white}Gratuit \textbullet\ Google Play\\Tuteur IA, annales, concours, résultats\par}
      \end{minipage}
    \end{tcolorbox}
  \end{minipage}\hfill
  \begin{minipage}[t]{0.487\linewidth}
    \begin{tcolorbox}[enhanced,colback=poppurple,colframe=popink,boxrule=2.2pt,arc=10pt,
        drop shadow=popink,left=11pt,right=11pt,top=10pt,bottom=10pt,height=3.1cm,valign=center]
      \tikz[baseline=-0.7ex]\node[circle,fill=white,draw=popink,line width=1.3pt,minimum size=1.6cm,inner sep=0]{\popdisplay\fontsize{24}{24}\selectfont\color{poppurple}+};%
      \hspace{8pt}\begin{minipage}[c]{0.6\linewidth}
        {\popdisplay\fontsize{11}{12.5}\selectfont\color{white}\MakeUppercase{Passe à OnBuch+}}\par\vspace{4pt}
        {\raggedright\popsemi\fontsize{8.4}{11}\selectfont\color{white}Tous les cours signés, Tuteur IA illimité, fiches PDF — onglet OnBuch+ de l'app\par}
      \end{minipage}
    \end{tcolorbox}
  \end{minipage}
  \vspace{8pt}
  \popcontenu
  \vfill
  \signatureonbuch
  \restoregeometry
  \newpage
}

% Rangée « Ce cours contient » (page de garde)
\newcommand{\poptile}[3]{% #1 couleur #2 gros texte #3 légende
  \begin{tcolorbox}[enhanced,colback=#1,colframe=popink,boxrule=2pt,arc=9pt,shadow={2.5pt}{-2.5pt}{0pt}{popink},
      left=6pt,right=6pt,top=9pt,bottom=9pt,halign=center,valign=center,height=1.75cm,width=\linewidth,nobeforeafter]
    {\popdisplay\fontsize{12.5}{13}\selectfont\color{white}\MakeUppercase{#2}}\par\vspace{3pt}
    {\centering\popsemi\fontsize{7.8}{9.5}\selectfont\color{white}#3\par}
  \end{tcolorbox}}
\newcommand{\popcontenu}{%
  \par\noindent{\popxboldls\fontsize{7.5}{7.5}\selectfont\color{popmuted}CE COURS SIGNÉ CONTIENT}\par\vspace{7pt}
  \noindent\begin{minipage}[t]{0.235\linewidth}\poptile{popblue}{Cours}{complet, illustré, pas à pas}\end{minipage}\hfill
  \begin{minipage}[t]{0.235\linewidth}\poptile{poporange}{Méthodes}{et exercices résolus}\end{minipage}\hfill
  \begin{minipage}[t]{0.235\linewidth}\poptile{poppink}{Exercices}{type examen + corrigés}\end{minipage}\hfill
  \begin{minipage}[t]{0.235\linewidth}\poptile{popgreen}{Fiche bilan}{l'essentiel à retenir}\end{minipage}\par}

% Sommaire
\newtcolorbox{sommaire}{enhanced,colback=white,colframe=popink,boxrule=2.2pt,arc=10pt,
  drop shadow=popink,left=18pt,right=18pt,top=13pt,bottom=13pt,before skip=6pt,after skip=20pt,
  before upper={\poppill{Sommaire}{poppurple}{white}\par\vspace{9pt}\popsemi\fontsize{10.6}{14.5}\selectfont\color{popink}}}

% Signature de fin de cours — poussée tout en bas de la dernière page
\newcommand{\findecours}{%
  \par\vfill\nopagebreak
  \begin{center}\popsquiggle{poporange}\end{center}\vspace{4pt}
  \signatureonbuch
}
\AtBeginDocument{\def\llbracket{\mathopen{[\mkern-3.5mu[}}\def\rrbracket{\mathclose{]\mkern-3.5mu]}}} % intervalles d'entiers (glyphe ⟦ absent de la police)
% Glyphes absents de Fira Math : redéfinis à partir de glyphes disponibles
\AtBeginDocument{%
  \def\setminus{\mathbin{\backslash}}\let\smallsetminus\setminus
  \def\vdots{\mathord{\vbox{\baselineskip3.2pt\lineskiplimit0pt\kern4pt\hbox{.}\hbox{.}\hbox{.}}}}%
  \def\ddots{\mathinner{\mkern1mu\raise6.5pt\hbox{.}\mkern2mu\raise3.6pt\hbox{.}\mkern2mu\raise0.7pt\hbox{.}\mkern1mu}}%
  \def\triangle{\mathord{\scalebox{0.9}{$\symup{\Delta}$}}}%
  \def\nabla{\mathord{\raisebox{\depth}{\scalebox{1}[-1]{$\symup{\Delta}$}}}}%
  \def\checkmark{\popcheck{popdark}}%
  \def\star{\mathbin{\ast}}\def\bigstar{\mathord{\ast}}%
  \def\diamond{\mathbin{\scalebox{0.6}{$\Diamond$}}}%
  \def\bigcup{\mathop{\mathchoice{\raisebox{-0.3em}{\scalebox{1.7}{$\cup$}}}{\scalebox{1.25}{$\cup$}}{\cup}{\cup}}\displaylimits}%
  \def\bigcap{\mathop{\mathchoice{\raisebox{-0.3em}{\scalebox{1.7}{$\cap$}}}{\scalebox{1.25}{$\cap$}}{\cap}{\cap}}\displaylimits}%
  \def\lesseqqgtr{\mathrel{\lessgtr}}\def\bigoplus{\mathop{\oplus}\displaylimits}%
}
\providecommand{\ssi}{si et seulement si} % abréviation fréquente
\AtBeginDocument{\providecommand{\wideparen}{\overparen}} % arc de cercle

%% FIGFIX-BEGIN v5 — correctifs globaux des figures (docs/onbuch-generation/tools/figfix)
\usepackage{adjustbox}\usepackage{etoolbox}\tcbuselibrary{raster}
% toute figure TikZ/pgfplots plus large que la ligne est réduite à la largeur disponible
\BeforeBeginEnvironment{tikzpicture}{\begin{adjustbox}{max width=\linewidth}}
\AfterEndEnvironment{tikzpicture}{\end{adjustbox}}
% légendes pgfplots lisibles par-dessus les courbes
\pgfplotsset{every axis/.append style={legend style={fill=white,fill opacity=0.92,text opacity=1,draw=black!15,inner sep=3pt}}}
% étiquettes d'arêtes (syntaxe edge["..."]) avec fond blanc
\tikzset{every edge quotes/.append style={fill=white,inner sep=1.2pt}}
%% FIGFIX-END
\begin{document}

\pagegarde{Reboisement et écotourisme}{Allemand}{1ère A}{Chapitre 3 : Environnement, santé et bien-être}{ALA10}{2 h}{Cette leçon de type « texte » propose la lecture et le commentaire d'un texte allemand sur le reboisement et l'écotourisme. Elle développe le vocabulaire thématique de la forêt et de la protection de l'environnement, et réinvestit l'impératif ainsi que les prépositions suivies de l'accusatif dans une production guidée.
\vspace{4pt}
\begin{multicols}{2}\small
\begin{itemize}
\item À la fin de cette leçon, je sais comprendre globalement puis en détail un texte allemand sur le reboisement et l'écotourisme.
\item À la fin de cette leçon, je sais employer correctement le vocabulaire du thème (der Wald, der Baum, pflanzen, der Naturpark, schützen, die Umwelt) avec articles et pluriels.
\item À la fin de cette leçon, je sais former et utiliser l'impératif allemand (2e pers. sing., 2e pers. plur., 1re pers. plur., forme de politesse).
\item À la fin de cette leçon, je sais utiliser les prépositions suivies de l'accusatif (durch, für, gegen, ohne, um, bis, entlang).
\item À la fin de cette leçon, je sais répondre par écrit et à l'oral à des questions de compréhension sur un texte.
\item À la fin de cette leçon, je sais résumer un texte en quelques phrases avec mes propres mots.
\end{itemize}\end{multicols}}

\begin{sommaire}
\begin{multicols}{2}
\begin{enumerate}
\item Texte support : « Aufforstung und Ökotourismus in Kamerun »
\item Compréhension du texte
\item Lexique thématique : la forêt et la protection de l'environnement
\item Grammaire 1 : l'impératif
\item Grammaire 2 : les prépositions avec l'accusatif
\item Méthode du commentaire et production guidée
\item Activité d'intégration
\item Méthodes et exercices résolus
\item Exercices
\item Corrigés des exercices
\item Fiche bilan
\end{enumerate}
\end{multicols}
\end{sommaire}

\begin{objectifs}
\begin{itemize}
\item À la fin de cette leçon, je sais comprendre globalement puis en détail un texte allemand sur le reboisement et l'écotourisme.
\item À la fin de cette leçon, je sais employer correctement le vocabulaire du thème (der Wald, der Baum, pflanzen, der Naturpark, schützen, die Umwelt) avec articles et pluriels.
\item À la fin de cette leçon, je sais former et utiliser l'impératif allemand (2e pers. sing., 2e pers. plur., 1re pers. plur., forme de politesse).
\item À la fin de cette leçon, je sais utiliser les prépositions suivies de l'accusatif (durch, für, gegen, ohne, um, bis, entlang).
\item À la fin de cette leçon, je sais répondre par écrit et à l'oral à des questions de compréhension sur un texte.
\item À la fin de cette leçon, je sais résumer un texte en quelques phrases avec mes propres mots.
\item À la fin de cette leçon, je sais rédiger un court paragraphe ou une affiche de sensibilisation à la protection de l'environnement.
\item À la fin de cette leçon, je sais repérer et éviter les erreurs typiques des francophones (faux amis, ordre des mots, cas).
\end{itemize}
\end{objectifs}

\begin{prerequis}
\begin{itemize}
\item Le présent de l'indicatif des verbes réguliers et irréguliers (Seconde)
\item Les articles définis et indéfinis au nominatif et à l'accusatif (Seconde)
\item Le vocabulaire de base de la nature et des saisons (Seconde)
\item La place du verbe en position 2 dans la phrase déclarative (Seconde)
\item Les questions totales et partielles avec W-Fragewörter (début d'année de 1ère)
\end{itemize}
\end{prerequis}

\newpage

\coursec{Texte support : « Aufforstung und Ökotourismus in Kamerun »}

Au Cameroun, des élèves de lycées participent chaque année à des campagnes de reboisement autour de Yaoundé et de Kribi : on distribue des plants, on creuse des trous, on arrose, on protège les jeunes arbres. En Allemagne, en Autriche et en Suisse, les forêts couvrent près d'un tiers du territoire et les \all{Naturparke} attirent des millions de touristes soucieux de la nature. Comment parler en allemand de ces actions pour protéger l'environnement ? Comment donner des consignes à des volontaires qui plantent des arbres ? Cette leçon nous donne les mots et les structures pour lire, comprendre et commenter un texte sur le reboisement et l'écotourisme.

\textbf{Problématique :} comment un texte allemand simple raconte-t-il un projet de reboisement, et quels outils de langue (vocabulaire, impératif, prépositions) utilise-t-il pour informer et mobiliser le lecteur ?

\courssub{Présentation du texte}

Le texte que nous allons lire est un \cle{reportage} (\all{der Bericht}) : un journaliste décrit un projet de reboisement dans un village près de Yaoundé. Il donne la parole à un professeur et à des élèves. Ce type de texte mélange deux choses : des \textbf{faits} (qui ? où ? quoi ?) et des \textbf{citations} au discours direct (entre guillemets allemands „…“).

Avant de lire, observons le titre : \all{Aufforstung und Ökotourismus in Kamerun}. Deux mots composés : \all{die Aufforstung} (le reboisement, de \all{der Wald} → on « refait » la forêt) et \all{der Ökotourismus} (l'écotourisme). Le titre annonce donc le thème : la protection de la nature au Cameroun.

\begin{methode}[Lire un texte en trois étapes]
\begin{enumerate}
\item \textbf{Lecture globale} : lire le texte une première fois \emph{sans dictionnaire}. Repérer le thème, le lieu, les personnages. Ne pas s'arrêter sur les mots inconnus.
\item \textbf{Lecture détaillée} : relire avec le glossaire. Souligner les mots du champ lexical, chercher les structures grammaticales (ici : impératifs, prépositions + accusatif).
\item \textbf{Reformulation} : raconter le texte avec ses propres mots, d'abord en français, puis en allemand simple.
\end{enumerate}
\end{methode}

\courssub{Lecture globale du texte}

\begin{textebox}[Aufforstung und Ökotourismus in Kamerun]
In einem Dorf in der Nähe von Yaoundé pflanzen Schüler und Freiwillige jedes Jahr tausend Bäume. „Pflanzt Bäume und schützt die Umwelt!“, sagt der Lehrer. Die jungen Leute arbeiten für die Zukunft ihres Dorfes.

Am Morgen kommen die Schüler mit Sämlingen und Werkzeugen zum Hügel hinter dem Dorf. Sie graben Löcher, setzen die jungen Bäume in die Erde und geben ihnen Wasser. Die Arbeit ist schwer, aber alle sind stolz. „Unser Wald ist unser Leben“, sagt Amina, eine Schülerin aus der ersten Klasse. „Ohne Bäume gibt es keinen Regen und keine Tiere.“

Durch den neuen Wald kommen jetzt auch Touristen: Sie wandern ohne Lärm durch den Naturpark und beobachten Vögel und Affen. Ein Guide aus dem Dorf zeigt ihnen die alten Bäume und erklärt die Namen der Pflanzen. Die Touristen schlafen in kleinen Häusern bei Familien und essen im Dorf.

Das Dorf verdient Geld mit dem Ökotourismus und die Natur bleibt geschützt. Mit den Einnahmen kauft die Schule Bücher und Material für die Klassen. „Geht mit uns in den Wald und helft uns!“, rufen die Schüler. „Besucht unseren Naturpark und entdeckt die Schönheit Kameruns!“
\end{textebox}

\textbf{Glossaire (mots nouveaux) :}

\begin{center}
\begin{tabularx}{\linewidth}{|l|X|}
\hline
\rowcolor{popblueL} \textbf{Deutsch} & \textbf{Français} \\
\hline
die Aufforstung, -en & le reboisement \\
\hline
der Freiwillige, -n & le/la volontaire, bénévole (nom faible : \all{dem Freiwilligen}, \all{des Freiwilligen}) \\
\hline
der Sämling, -e & le plant, le semis (jeune arbre) \\
\hline
das Werkzeug, -e & l'outil \\
\hline
der Hügel, - & la colline (pl. \all{die Hügel}) \\
\hline
graben (gräbt, grub, gegraben) & creuser (verbe fort) \\
\hline
der Lebensraum, -räume & l'espace de vie, l'habitat \\
\hline
der Lärm & le bruit \\
\hline
nachhaltig & durable (respectueux de l'avenir) \\
\hline
die Einnahmen (Pl.) & les recettes, les revenus \\
\hline
schützen & protéger \\
\hline
die Schönheit, -en & la beauté \\
\hline
\end{tabularx}
\end{center}

\textbf{Repérage global (première lecture) :} réponds mentalement à ces questions avant toute traduction.

\begin{itemize}
\item \textbf{Ort (lieu) :} un village près de Yaoundé, au Cameroun ; une colline derrière le village ; un Naturpark.
\item \textbf{Personen (personnages) :} des élèves (\all{Schüler}), des volontaires (\all{Freiwillige}), un professeur (\all{der Lehrer}), Amina, un guide, des touristes.
\item \textbf{Thema (thème) :} le reboisement et l'écotourisme comme source de revenus.
\item \textbf{Aktionen (actions) :} planter, creuser, arroser, protéger, observer, gagner de l'argent.
\end{itemize}

\begin{popfigure}
\begin{center}
\begin{center}\tcbox[colback=popblue,colframe=popblue,coltext=white,arc=9pt,boxsep=4pt,left=6pt,right=6pt]{\bfseries\small der Text}\end{center}
\vspace{-2pt}
\begin{tcbraster}[raster columns=2,raster equal height=rows,raster column skip=6pt,raster row skip=6pt,enhanced,blankest]
\begin{tcolorbox}[enhanced,colback=popblue,colframe=popblue,coltext=white,boxrule=0.9pt,arc=3pt,left=4pt,right=4pt,top=3pt,bottom=3pt,boxsep=1pt,halign=left]\bfseries\scriptsize Personen: Schüler, Lehrer, Freiwillige, Touristen\end{tcolorbox}
\begin{tcolorbox}[enhanced,colback=poporange,colframe=poporange,coltext=white,boxrule=0.9pt,arc=3pt,left=4pt,right=4pt,top=3pt,bottom=3pt,boxsep=1pt,halign=left]\bfseries\scriptsize Ort: Dorf bei Yaoundé, Naturpark\end{tcolorbox}
\begin{tcolorbox}[enhanced,colback=popgreen,colframe=popgreen,coltext=white,boxrule=0.9pt,arc=3pt,left=4pt,right=4pt,top=3pt,bottom=3pt,boxsep=1pt,halign=left]\bfseries\scriptsize Thema: Aufforstung, Ökotourismus\end{tcolorbox}
\begin{tcolorbox}[enhanced,colback=poppurple,colframe=poppurple,coltext=white,boxrule=0.9pt,arc=3pt,left=4pt,right=4pt,top=3pt,bottom=3pt,boxsep=1pt,halign=left]\bfseries\scriptsize Aktionen: pflanzen, graben, schützen, wandern\end{tcolorbox}
\end{tcbraster}

\end{center}
\legende{Carte mentale : les quatre repères de la lecture globale.}
\end{popfigure}

\courssub{Lecture détaillée : le champ lexical de la forêt et de l'environnement}

Soulignons dans le texte tous les mots qui appartiennent au champ lexical de la \textbf{forêt} et de la \textbf{protection de l'environnement} :

\begin{itemize}
\item \textbf{La nature :} \all{der Wald} (la forêt), \all{der Baum, die Bäume} (l'arbre), \all{die Pflanzen} (les plantes), \all{die Erde} (la terre), \all{der Regen} (la pluie), \all{die Tiere} (les animaux), \all{die Vögel} (les oiseaux), \all{die Affen} (les singes), \all{der Naturpark} (le parc naturel).
\item \textbf{Les actions de protection :} \all{pflanzen} (planter), \all{schützen} (protéger), \all{graben} (creuser), \all{helfen} (aider), \all{Wasser geben} (arroser, litt. « donner de l'eau »).
\item \textbf{Le tourisme :} \all{die Touristen}, \all{wandern} (randonner), \all{beobachten} (observer), \all{besuchen} (visiter), \all{entdecken} (découvrir), \all{der Guide} (le guide).
\end{itemize}

\begin{definition}[L'écotourisme]
\all{Der Ökotourismus} (l'écotourisme) désigne un tourisme \textbf{respectueux de la nature} et \textbf{profitable aux populations locales} : les visiteurs découvrent la nature sans la détruire (\all{ohne Lärm}), et l'argent qu'ils dépensent reste dans le village (\all{die Einnahmen}). Mot composé : \all{Öko-} (écologie) + \all{Tourismus}.
\end{definition}

\begin{popfigure}
\begin{center}
\begin{tikzpicture}[node distance=0.6cm and 1.4cm]
\node[draw=popink, fill=poppinkL, rounded corners, align=center, font=\small\bfseries] (p) {Problem\\ Abholzung\\ (déforestation)};
\node[draw=popink, fill=popblueL, rounded corners, align=center, font=\small\bfseries, right=of p] (pr) {Projekt\\ Bäume pflanzen\\ (reboisement)};
\node[draw=popink, fill=popgreenL, rounded corners, align=center, font=\small\bfseries, right=of pr] (e) {Ergebnis\\ Ökotourismus\\ Einnahmen};
\draw[-{Stealth}, very thick, poporange] (p) -- (pr);
\draw[-{Stealth}, very thick, popgreen] (pr) -- (e);
\end{tikzpicture}
\end{center}
\legende{Frise logique du texte : du problème au résultat.}
\end{popfigure}

\courssub{Traduction guidée des phrases clés}

\begin{enumerate}
\item \all{In einem Dorf in der Nähe von Yaoundé pflanzen Schüler und Freiwillige jedes Jahr tausend Bäume.} → « Dans un village près de Yaoundé, des élèves et des volontaires plantent chaque année mille arbres. » (Note : \all{in der Nähe von} + Datif = « à proximité de ».)
\item \all{Die jungen Leute arbeiten für die Zukunft ihres Dorfes.} → « Les jeunes travaillent pour l'avenir de leur village. » (\all{für} + Akkusatif ; \all{ihres Dorfes} : Génitif).
\item \all{Unser Wald ist unser Leben.} → « Notre forêt est notre vie. »
\item \all{Ohne Bäume gibt es keinen Regen und keine Tiere.} → « Sans arbres, il n'y a pas de pluie ni d'animaux. » (\all{es gibt} + Akkusatif = « il y a » ; \all{keinen/keine} = négation de l'article indéfini).
\item \all{Das Dorf verdient Geld mit dem Ökotourismus und die Natur bleibt geschützt.} → « Le village gagne de l'argent grâce à l'écotourisme et la nature reste protégée. » (\all{mit} + Datif ; \all{bleibt geschützt} = passif d'état).
\end{enumerate}

\begin{exoresolu}[Questions de compréhension]
Réponds en allemand, par des phrases complètes : 1. \all{Wer pflanzt die Bäume?} 2. \all{Wo arbeiten die Schüler am Morgen?} 3. \all{Warum kommen Touristen ins Dorf?}
\tcblower
1. \all{Schüler und Freiwillige pflanzen die Bäume.} (Des élèves et des volontaires plantent les arbres.) 2. \all{Am Morgen arbeiten die Schüler auf dem Hügel hinter dem Dorf.} (Le matin, les élèves travaillent sur la colline derrière le village.) 3. \all{Die Touristen kommen ins Dorf, weil sie durch den Naturpark wandern und Vögel und Affen beobachten wollen.} (Les touristes viennent parce qu'ils veulent randonner dans le parc et observer les oiseaux et les singes.)
\end{exoresolu}

\begin{savaistu}[Forêts et parcs naturels en Allemagne]
L'Allemagne compte plus de 100 \all{Naturparke} (parcs naturels) et la forêt (\all{der Wald}) couvre environ 32 \% du pays. Le mot \all{Wald} est un des plus anciens mots de la langue allemande, et le « bain de forêt » (\all{das Waldbaden}) y est devenu une activité de bien-être très à la mode. Au Cameroun, la forêt du bassin du Congo est la deuxième plus grande forêt tropicale du monde : un trésor à protéger, comme le font les élèves de notre texte !
\end{savaistu}

% ============================================================
% GRAMMAIRE 1 : L'IMPÉRATIF
% ============================================================
\coursec{Grammaire 1 : L'impératif (Der Imperativ)}

L'impératif exprime un \textbf{ordre}, un \textbf{conseil}, une \textbf{invitation} ou une \textbf{interdiction}. En allemand, il existe trois formes selon la personne à qui l'on s'adresse (du, ihr, Sie) et une forme pour la 1\textsuperscript{re} personne du pluriel (wir = « faisons... »).

\begin{propriete}[Formation de l'impératif]
\begin{center}
\begin{tabularx}{\linewidth}{|l|X|X|X|X|}
\hline
\rowcolor{popblueL} \textbf{Personne} & \textbf{Règle de formation} & \textbf{Exemple : \all{pflanzen}} & \textbf{Exemple : \all{graben}} & \textbf{Exemple : \all{sein}} \\
\hline
\all{du} (2\textsuperscript{e} sg) & Radical du présent \textbf{sans -e} (souvent) & \all{Pflanz!} / \all{Pflanze!} & \all{Grab!} & \all{Sei!} \\
\hline
\all{ihr} (2\textsuperscript{e} pl) & \textbf{Identique au présent indicatif} (\all{ihr pflanzt}) & \all{Pflanzt!} & \all{Grabt!} & \all{Seid!} \\
\hline
\all{Sie} (politesse) & \textbf{Infinitif + \all{Sie}} (verbe en 1\textsuperscript{re} position) & \all{Pflanzen Sie!} & \all{Graben Sie!} & \all{Seien Sie!} \\
\hline
\all{wir} (1\textsuperscript{re} pl) & \textbf{Infinitif + \all{wir}} (verbe en 1\textsuperscript{re} position) & \all{Pflanzen wir!} & \all{Graben wir!} & \all{Seien wir!} \\
\hline
\end{tabularx}
\end{center}
\end{propriete}

\begin{attention}[Particularités importantes pour les francophones]
\begin{enumerate}
\item \textbf{Position du verbe :} L'impératif est \textbf{toujours en première position} dans la phrase (sauf si un mot de liaison comme \all{und} ou \all{aber} précède). \all{Pflanzt Bäume!} — Pas de sujet exprimé pour \all{du/ihr}.
\item \textbf{Verbes à voyelle modifiée (du-form) :} Si le présent change de voyelle (\all{e} → \all{i/ie}, \all{a} → \all{ä}), l'impératif \all{du} garde la voyelle modifiée \textbf{sans -e} final. \all{geben} → \all{du gibst} → \all{Gib!} ; \all{graben} → \all{du gräbst} → \all{Grab!} ; \all{lesen} → \all{du liest} → \all{Lies!} ; \all{fahren} → \all{du fährst} → \all{Fahr!}.
\item \textbf{Verbes en -eln / -ern :} Garde souvent le \all{-e} à la forme \all{du} : \all{handeln} → \all{Handle!} ; \all{lächeln} → \all{Lächele!}.
\item \textbf{Verbes à particule séparable :} La particule va à la \textbf{fin} de la phrase. \all{anfangen} → \all{Fangt an!} ; \all{mitmachen} → \all{Macht mit!} ; \all{aufräumen} → \all{Räumt auf!}
\item \textbf{Verbes pronominaux (réfléchis) :} Le pronom réfléchi suit le verbe. \all{sich setzen} → \all{Setz dich!} / \all{Setzt euch!} / \all{Setzen Sie sich!}
\item \textbf{\all{sein} et \all{haben} :} Irréguliers. \all{Sei! / Seid! / Seien Sie!} — \all{Hab! / Habt! / Haben Sie!}
\end{enumerate}
\end{attention}

\begin{exemplebox}[Impératifs dans le texte support]
\begin{itemize}
\item \all{„Pflanzt Bäume und schützt die Umwelt!“} (\all{ihr}-form, s'adresse au groupe).
\item \all{„Geht mit uns in den Wald und helft uns!“} (\all{geht}, \all{helft} = \all{ihr}-form de \all{gehen}, \all{helfen}).
\item \all{„Besucht unseren Naturpark und entdeckt die Schönheit Kameruns!“} (\all{besucht}, \all{entdeckt} = \all{ihr}-form).
\end{itemize}
\textit{Remarque : Le texte utilise exclusivement la forme \all{ihr} (pluriel), car le professeur et les élèves s'adressent à un groupe (la communauté, les touristes).}
\end{exemplebox}

\begin{exercice}[Entraînement : L'impératif]
Mets les verbes entre parenthèses à la forme d'impératif demandée.
\begin{enumerate}
\item \all{(pflanzen, du)} \_\_\_\_\_\_\_\_\_\_\_\_ einen Baum!
\item \all{(graben, ihr)} \_\_\_\_\_\_\_\_\_\_\_\_ das Loch!
\item \all{(schützen, Sie)} \_\_\_\_\_\_\_\_\_\_\_\_ die Natur!
\item \all{(helfen, wir)} \_\_\_\_\_\_\_\_\_\_\_\_ den Schülern!
\item \all{(aufpassen, du)} \_\_\_\_\_\_\_\_\_\_\_\_ auf die Sämlinge! (verbe séparable)
\item \all{(sich setzen, ihr)} \_\_\_\_\_\_\_\_\_\_\_\_ hin! (verbe pronominal)
\end{enumerate}
\end{exercice}

\begin{corrige}[Exercice : L'impératif]
\begin{enumerate}
\item \all{Pflanz!} (ou \all{Pflanze!}) \all{einen Baum!}
\item \all{Grabt!} \all{das Loch!}
\item \all{Schützen Sie!} \all{die Natur!}
\item \all{Helfen wir!} \all{den Schülern!} (Datif après \all{helfen})
\item \all{Pass auf!} \all{auf die Sämlinge!} (particule \all{auf} à la fin)
\item \all{Setzt euch!} \all{hin!} (pronom \all{euch} + particule \all{hin} à la fin)
\end{enumerate}
\end{corrige}

% ============================================================
% GRAMMAIRE 2 : PRÉPOSITIONS AVEC L'ACCUSATIF
% ============================================================
\coursec{Grammaire 2 : Les prépositions avec l'accusatif (Präpositionen mit Akkusativ)}

Certaines prépositions allemandes \textbf{exigent toujours l'accusatif}, quel que soit le contexte (mouvement ou position). Il faut les connaître par cœur.

\begin{propriete}[Les prépositions « accusatif seul » (FUDGOBUW)]
\begin{center}
\begin{tabularx}{\linewidth}{|l|l|X|}
\hline
\rowcolor{popblueL} \textbf{Préposition} & \textbf{Sens principal} & \textbf{Exemple (article masc. \all{der Wald})} \\
\hline
\all{durch} & à travers, par (moyen), grâce à & \all{durch den Wald} (à travers la forêt) \\
\hline
\all{für} & pour, en faveur de & \all{für den Wald} (pour la forêt) \\
\hline
\all{gegen} & contre, vers (direction vague, temps) & \all{gegen den Wald} (contre la forêt) ; \all{gegen Abend} \\
\hline
\all{ohne} & sans & \all{ohne den Wald} (sans la forêt) \\
\hline
\all{um} & autour de, à (heure précise) & \all{um den Wald} (autour de la forêt) ; \all{um 10 Uhr} \\
\hline
\all{bis} & jusqu'à (souvent sans article) & \all{bis zum Wald} (\all{bis zu dem}) ; \all{bis nächstes Jahr} \\
\hline
\all{entlang} & le long de (souvent \textbf{après} le nom) & \all{den Wald entlang} (le long de la forêt) \\
\hline
\all{wider} & contre (formel, littéraire) & \all{wider den Wald} \\
\hline
\end{tabularx}
\end{center}
\textbf{Mnémotechnique :} \cle{FUDGOBUW} (\all{Für, Um, Durch, Gegen, Ohne, Bis, Entlang, Wider}).
\end{propriete}

\begin{attention}[Pièges fréquents pour francophones]
\begin{enumerate}
\item \textbf{\all{durch} vs \all{wegen} :} \all{durch} + Akk = « par le biais de, grâce à » (cause efficiente). \all{wegen} + Génitif (langage soutenu) / Datif (langage courant) = « à cause de ». Dans le texte : \all{Durch den neuen Wald kommen Touristen} = « Grâce à la nouvelle forêt... » (cause positive).
\item \textbf{\all{für} = « pour » (but, destinataire) :} Ne pas confondre avec \all{vor} (devant, avant) + Datif/Akk.
\item \textbf{\all{ohne} + Akk :} Toujours sans article indéfini au pluriel ? \all{Ohne Bäume} (sans arbres). Correct.
\item \textbf{\all{entlang} :} Place \textbf{après} le nom (\all{den Fluss entlang}). Si placé avant, il prend le Datif (\all{entlang des Flusses}) — niveau B1/B2.
\item \textbf{Prépositions « Wechselpräpositionen » (in, an, auf, über, unter, vor, hinter, neben, zwischen) :} Elles prennent l'\textbf{Accusatif} s'il y a \textbf{mouvement vers un lieu} (Wo hin?), le \textbf{Datif} s'il y a \textbf{position statique} (Wo?).
\end{enumerate}
\end{attention}

\begin{exemplebox}[Prépositions + Akk dans le texte]
\begin{itemize}
\item \all{arbeiten \textbf{für die} Zukunft} (\all{für} + Akk féminin \all{die Zukunft}).
\item \all{\textbf{Durch den} neuen Wald kommen...} (\all{durch} + Akk masc \all{den Wald}).
\item \all{wandern \textbf{ohne} Lärm \textbf{durch den} Naturpark} (\all{ohne} + Akk \all{Lärm} ; \all{durch} + Akk \all{den Naturpark}).
\item \all{Geht \textbf{in den} Wald!} (\all{in} + Akk masc \all{den Wald} → mouvement \all{Wo hin?}).
\item \all{für die Klassen} (\all{für} + Akk pluriel \all{die Klassen}).
\end{itemize}
\end{exemplebox}

\begin{propriete}[Déclinaison de l'article défini après préposition accusative]
\begin{center}
\begin{tabular}{|l|c|c|c|c|}
\hline
\rowcolor{popblueL} \textbf{Cas} & \textbf{Masculin} & \textbf{Féminin} & \textbf{Neutre} & \textbf{Pluriel} \\
\hline
\textbf{Accusatif} & \all{den} Wald & \all{die} Zukunft & \all{das} Kind & \all{die} Bäume \\
\hline
\end{tabular}
\end{center}
\end{propriete}

\begin{exercice}[Entraînement : Prépositions + Accusatif]
Complète avec l'article correct (défini ou indéfini) à l'accusatif.
\begin{enumerate}
\item Die Schüler pflanzen Bäume \_\_\_\_\_\_\_\_\_\_\_\_ (für) \_\_\_\_\_\_\_\_\_\_\_\_ Zukunft.
\item Wir wandern \_\_\_\_\_\_\_\_\_\_\_\_ (durch) \_\_\_\_\_\_\_\_\_\_\_\_ Wald.
\item \_\_\_\_\_\_\_\_\_\_\_\_ (ohne) \_\_\_\_\_\_\_\_\_\_\_\_ Lärm ist es schöner.
\item Der Guide zeigt \_\_\_\_\_\_\_\_\_\_\_\_ (um) \_\_\_\_\_\_\_\_\_\_\_\_ Naturpark herum.
\item Geht \_\_\_\_\_\_\_\_\_\_\_\_ (in) \_\_\_\_\_\_\_\_\_\_\_\_ Park! (mouvement)
\end{enumerate}
\end{exercice}

\begin{corrige}[Exercice : Prépositions + Accusatif]
\begin{enumerate}
\item \all{für die} Zukunft (féminin : \all{die}).
\item \all{durch den} Wald (masculin : \all{den}).
\item \all{Ohne} \all{Lärm} (masculin, pas d'article indéfini au singulier négatif souvent, ou \all{den} Lärm si défini. Ici sens général : \all{Ohne Lärm}).
\item \all{um den} Naturpark (masculin : \all{den}).
\item \all{in den} Park (masculin, mouvement → Akk : \all{den}).
\end{enumerate}
\end{corrige}

% ============================================================
% MÉTHODE DU COMMENTAIRE ET PRODUCTION GUIDÉE
% ============================================================
\coursec{Méthode du commentaire et production guidée}

\courssub{Méthode : Commenter un texte court (Niveau A2)}

\begin{methode}[Commenter un texte en 4 étapes]
\begin{enumerate}
\item \textbf{Introduction (1 phrase) :} Type de texte, titre, auteur (si connu), thème, lieu.
\all{Der Text „Aufforstung und Ökotourismus in Kamerun“ ist ein Bericht über ein Projekt in einem Dorf bei Yaoundé.}
\item \textbf{Contenu (2-3 phrases) :} Résumer les faits principaux (Qui ? Quoi ? Où ? Résultat ?). Utiliser \all{Präteritum} ou \all{Perfekt}.
\all{Schüler und Freiwillige pflanzen jedes Jahr Bäume. Sie arbeiten für die Zukunft. Durch den Wald kommen Touristen. Das Dorf verdient Geld.}
\item \textbf{Analyse linguistique (1-2 phrases) :} Relever 2-3 structures étudiées (impératifs, prépositions, vocabulaire).
\all{Der Text benutzt viele Imperative: „Pflanzt!“, „Geht!“, „Besucht!“. Es gibt Präpositionen mit Akkusativ: für, durch, ohne, in.}
\item \textbf{Opinion personnelle / Conclusion (1 phrase) :} \all{Ich finde das Projekt gut, weil es die Natur schützt und dem Dorf hilft.}
\end{enumerate}
\end{methode}

\courssub{Production écrite guidée : « Un appel pour la forêt »}

\begin{experience}[Jeu de rôle / Production orale : « Engagement pour l'environnement »]
\textbf{Situation :} Vous êtes délégué(e) écolo de votre lycée à Douala ou Yaoundé. Vous organisez une journée « Plantons pour l'avenir ».
\textbf{Consigne :} Rédigez une courte affiche (5-6 phrases) en allemand pour motiver vos camarades. Utilisez :
\begin{itemize}
\item Au moins \textbf{3 impératifs} (formes \all{ihr} ou \all{Sie}).
\item Au moins \textbf{3 prépositions avec accusatif} (\all{für, durch, ohne, in, um, gegen}).
\item Du vocabulaire de la leçon (\all{der Wald, der Baum, pflanzen, schützen, die Umwelt, der Sämling, die Zukunft}).
\end{itemize}
\textbf{Aide au démarrage :}
\all{„Liebe Mitschüler, ... Pflanzt ... Schützt ... Geht ... Ohne ... Für ...“}
\end{experience}

\begin{exercice}[Expression écrite : L'affiche pour le reboisement]
Rédigez votre affiche (environ 60-80 mots) sur votre cahier. Relisez-vous en vérifiant :
\begin{enumerate}
\item Majuscules à \textbf{tous les noms}.
\item Verbe en \textbf{1\textsuperscript{re} position} pour les impératifs.
\item \textbf{Accusatif} correct après \all{für, durch, ohne, in} (mouvement).
\item Orthographe des \textbf{Umlaute} (\all{ä, ö, ü}) et \all{ß}.
\end{enumerate}
\end{exercice}

\begin{corrige}[Exercice : Expression écrite (Proposition de correction)]
\all{„Liebe Mitschüler, pflanzt Bäume für die Zukunft! Schützt unsere Umwelt! Geht mit uns in den Wald! Ohne Bäume gibt es kein Leben. Durch eure Hilfe wächst ein neuer Wald. Besucht unsere Aktion am Samstag!“}
\textit{Vérification : Impératifs (pflanzt, schützt, geht, besucht) ; Prépositions + Akk (für die Zukunft, in den Wald, ohne Bäume, durch eure Hilfe). Vocabulaire thématique présent.}
\end{corrige}

\begin{savaistu}[Le mot « Wald » dans la culture germanophone]
\all{Der Wald} n'est pas seulement un écosystème, c'est un \textbf{mythe culturel} en Allemagne. Les contes des frères Grimm (\all{Hänsel und Gretel, Rotkäppchen}) se déroulent dans la forêt sombre et dangereuse. Au XIXe siècle, les romantiques (Caspar David Friedrich) en font le symbole de l'âme allemande. Aujourd'hui, \all{Waldsterben} (mort des forêts) et \all{Waldbaden} (bain de forêt) montrent la relation ambivalente : menace écologique et recherche de bien-être. Au Cameroun, la forêt est aussi source de vie (cacao, bois, pharmacopée, climat) : la protéger, c'est protéger l'avenir commun.
\end{savaistu}

\coursec{Compréhension du texte}

Après avoir découvert le texte support \emph{« Aufforstung und Ökotourismus in Kamerun »} dans la section précédente, nous allons maintenant l'exploiter pas à pas : d'abord une compréhension globale pour saisir l'essentiel, puis une analyse détaillée pour maîtriser le vocabulaire et les structures, enfin une reformulation et un échange oral pour ancrer les acquis. Cette démarche progressive — du global au précis, du réceptif au productif — est la clé d'une lecture efficace en langue étrangère.

\courssub{Rappel du texte support}

\begin{textebox}[Aufforstung und Ökotourismus in Kamerun]
In Kamerun sind die Wälder sehr wichtig für das Klima und die Artenvielfalt. Doch die Abholzung durch die Holzindustrie und die Landwirtschaft gefährdet die Natur. Deshalb starten Schulen, Vereine und das Ministerium für Forstwirtschaft gemeinsame Aufforstungsaktionen. Jedes Jahr pflanzen Tausende von Schülern und Freiwilligen junge Bäume in den degradierten Zonen.

Ein Beispiel ist der Naturpark „Mont Cameroun“. Dort schützen Ranger die Tiere und führen Touristen auf Wanderwegen durch den Regenwald. Die Besucher lernen die heimischen Baumarten kennen und sehen, wie wichtig der Wald für das Wasser und den Boden ist. Der Eintritt in den Park finanziert die Aufforstung und die Arbeit der Ranger.

Ökotourismus bedeutet: Reisen ohne die Natur zu zerstören. Touristen schlafen in einfachen Öko-Lodges, essen lokale Produkte und respektieren die Regeln des Parks. So schaffen sie Einkommen für die Dörfer und schützen gleichzeitig die Umwelt. „Wir pflanzen heute für die Kinder von morgen“, sagt ein Ranger. Das ist nachhaltige Entwicklung.
\end{textebox}

\begin{center}
\begin{tabularx}{\linewidth}{|X|X|X|}
\hline
\rowcolor{popblueL}
\textbf{Deutsch} & \textbf{Français} & \textbf{Remarque} \\
\hline
der Wald, die Wälder & la forêt & nom masculin, pluriel en \textit{-ä-} + \textit{-er} \\
die Abholzung, die -en & la déforestation & nom féminin, dérivé du verbe \textit{abholzen} \\
die Holzindustrie, die -n & l'industrie du bois & composé : \textit{das Holz} + \textit{die Industrie} \\
die Artenvielfalt (kein Pl.) & la biodiversité & \textit{die Art} (espèce) + \textit{die Vielfalt} (diversité) \\
gefährden & menacer / mettre en danger & verbe faible régulier : \textit{gefährdet, gefährdete, hat gefährdet} \\
die Aufforstung, die -en & le reboisement & nom féminin, dérivé du verbe \textit{aufforsten} \\
degradiert & dégradé & adjectif (participe passé de \textit{degradieren}) \\
der Naturpark, die -s & le parc naturel & composé : \textit{die Natur} + \textit{der Park} \\
der Ranger, die Ranger & le garde forestier / l'éco-garde & emprunt à l'anglais, pluriel identique au singulier \\
der Wanderweg, die -e & le sentier de randonnée & \textit{wandern} (randonner) + \textit{der Weg} (chemin) \\
heimisch & indigène / local & adjectif : \textit{heimische Baumarten} \\
die Öko-Lodge, die -n & l'éco-lodge & féminin, emprunt récent à l'anglais \\
nachhaltig & durable / soutenable & adjectif clé : \textit{nachhaltige Entwicklung} \\
\hline
\end{tabularx}
\end{center}

\courssub{Compréhension globale : Wer? Wo? Was? Warum?}

Avant d'entrer dans le détail, répondons aux quatre questions fondamentales qui donnent le squelette du texte. En classe, on les pose d'abord à l'oral, puis on écrit les réponses au tableau.

\begin{exemplebox}[Questions globales et réponses modèles]
\textbf{Wer handelt im Text?} \\
\all{Schüler, Freiwillige, das Ministerium für Forstwirtschaft, Ranger und Touristen handeln im Text.} \\
« Ce sont des élèves, des volontaires, le ministère des Forêts, des éco-gardes et des touristes qui agissent dans le texte. »

\textbf{Wo spielt die Handlung?} \\
\all{Die Handlung spielt in Kamerun, besonders im Naturpark „Mont Cameroun“.} \\
« L'action se passe au Cameroun, surtout dans le parc naturel « Mont Cameroun ». »

\textbf{Was passiert?} \\
\all{Man pflanzt Bäume, schützt den Wald und betreibt Ökotourismus.} \\
« On plante des arbres, on protège la forêt et on pratique l'écotourisme. »

\textbf{Warum macht man das?} \\
\all{Weil die Abholzung die Natur gefährdet und man die Umwelt für die Zukunft schützen will.} \\
« Parce que la déforestation menace la nature et qu'on veut protéger l'environnement pour l'avenir. »
\end{exemplebox}

\begin{methode}[Répondre à une question de compréhension en allemand]
\textbf{Étape 1 :} Identifier le mot interrogatif de la question (\cle{Wer}, \cle{Wo}, \cle{Was}, \cle{Warum}, \cle{Wie}, \cle{Wann}). \\
\textbf{Étape 2 :} Relever l'information dans le texte (surligner ou noter le passage concerné). \\
\textbf{Étape 3 :} Formuler une \textbf{phrase complète} en reprenant le sujet et le verbe de la question. \\
\textbf{Étape 4 :} Respecter l'ordre des mots : \textbf{verbe conjugué en 2e position}. \\
\textbf{Étape 5 :} Vérifier l'orthographe des noms (majuscules !) et la conjugaison.
\end{methode}

\begin{attention}[Piège classique : l'ordre des mots dans la réponse]
En français, on dit souvent : « Les élèves plantent des arbres. » En allemand, le verbe conjugué \textbf{doit} occuper la 2e place :
\begin{itemize}
    \item Correct : \all{Die Schüler pflanzen Bäume.} (Sujet — Verbe — Objet)
    \item Correct : \all{Jedes Jahr pflanzen die Schüler Bäume.} (Complément — Verbe — Sujet — Objet)
    \item Incorrect : \all{Die Schüler Bäume pflanzen.} (ordre français, \textbf{faux} en phrase principale)
\end{itemize}
Même si la question commence par \all{Was pflanzen die Schüler?}, la réponse garde le verbe en 2e position : \all{Die Schüler pflanzen Bäume.}
\end{attention}

\courssub{Compréhension détaillée : six questions ciblées}

Voici six questions précises qui obligent à relire le texte avec attention. Les réponses attendues sont des phrases complètes, justifiées par le texte.

\begin{enumerate}
    \item \all{Was gefährdet die Natur in Kamerun?}
    \item \all{Wer pflanzt junge Bäume in den degradierten Zonen?}
    \item \all{Was machen die Ranger im Naturpark „Mont Cameroun“?}
    \item \all{Wie lernen die Besucher die Bedeutung des Waldes kennen?}
    \item \all{Wovon wird die Aufforstung finanziert?}
    \item \all{Was bedeutet „Ökotourismus“ laut Text?}
\end{enumerate}

\begin{exoresolu}[Réponses détaillées avec citations]
\textbf{1. Was gefährdet die Natur in Kamerun?} \\
\all{Die Abholzung durch die Holzindustrie und die Landwirtschaft gefährdet die Natur.} \\
\textit{Beleg :} „Doch die Abholzung durch die Holzindustrie und die Landwirtschaft gefährdet die Natur.“

\textbf{2. Wer pflanzt junge Bäume in den degradierten Zonen?} \\
\all{Tausende von Schülern und Freiwilligen pflanzen junge Bäume in den degradierten Zonen.} \\
\textit{Beleg :} „Jedes Jahr pflanzen Tausende von Schülern und Freiwilligen junge Bäume in den degradierten Zonen.“

\textbf{3. Was machen die Ranger im Naturpark „Mont Cameroun“?} \\
\all{Die Ranger schützen die Tiere und führen Touristen auf Wanderwegen durch den Regenwald.} \\
\textit{Beleg :} „Dort schützen Ranger die Tiere und führen Touristen auf Wanderwegen durch den Regenwald.“

\textbf{4. Wie lernen die Besucher die Bedeutung des Waldes kennen?} \\
\all{Die Besucher lernen die heimischen Baumarten kennen und sehen, wie wichtig der Wald für das Wasser und den Boden ist.} \\
\textit{Beleg :} „Die Besucher lernen die heimischen Baumarten kennen und sehen, wie wichtig der Wald für das Wasser und den Boden ist.“

\textbf{5. Wovon wird die Aufforstung finanziert?} \\
\all{Der Eintritt in den Park finanziert die Aufforstung und die Arbeit der Ranger.} \\
\textit{Beleg :} „Der Eintritt in den Park finanziert die Aufforstung und die Arbeit der Ranger.“

\textbf{6. Was bedeutet „Ökotourismus“ laut Text?} \\
\all{Ökotourismus bedeutet: Reisen ohne die Natur zu zerstören.} \\
\textit{Beleg :} „Ökotourismus bedeutet: Reisen ohne die Natur zu zerstören.“
\end{exoresolu}

\courssub{Vrai / Faux avec justification (Richtig / Falsch + Beleg)}

Cet exercice entraîne la lecture de près : il faut non seulement décider si l'affirmation est exacte, mais \textbf{la prouver par une citation littérale} du texte (recopier la phrase ou le segment concerné).

\begin{exercice}[Richtig oder falsch? Belegen Sie Ihre Antwort!]
Entscheiden Sie: \textbf{Richtig} oder \textbf{Falsch}? Zitieren Sie den Beleg aus dem Text.

\begin{enumerate}
    \item Die Wälder in Kamerun sind unwichtig für das Klima.
    \item Nur das Ministerium pflanzt Bäume.
    \item Im Naturpark „Mont Cameroun“ gibt es keine Touristen.
    \item Die Besucher sehen, dass der Wald für Wasser und Boden wichtig ist.
    \item Öko-Lodges sind große Luxushotels.
    \item Der Eintritt in den Park hilft, neue Bäume zu pflanzen.
    \item Touristen essen importierte Produkte in den Lodges.
    \item Der Ranger sagt: „Wir pflanzen heute für die Kinder von morgen.“
\end{enumerate}
\end{exercice}

\begin{corrige}[Exercice : Richtig oder falsch?]
\begin{enumerate}
    \item \textbf{Falsch.} Beleg : „In Kamerun sind die Wälder sehr wichtig für das Klima und die Artenvielfalt.“
    \item \textbf{Falsch.} Beleg : „Schulen, Vereine und das Ministerium für Forstwirtschaft starten gemeinsame Aufforstungsaktionen.“
    \item \textbf{Falsch.} Beleg : „Ranger … führen Touristen auf Wanderwegen durch den Regenwald.“
    \item \textbf{Richtig.} Beleg : „Die Besucher … sehen, wie wichtig der Wald für das Wasser und den Boden ist.“
    \item \textbf{Falsch.} Beleg : „Touristen schlafen in einfachen Öko-Lodges.“ (\textit{einfach} = simple, pas luxueux)
    \item \textbf{Richtig.} Beleg : „Der Eintritt in den Park finanziert die Aufforstung …“
    \item \textbf{Falsch.} Beleg : „… essen lokale Produkte …“ (\textit{lokal} = local, pas importé)
    \item \textbf{Richtig.} Beleg : „‚Wir pflanzen heute für die Kinder von morgen‘, sagt ein Ranger.“
\end{enumerate}
\end{corrige}

\courssub{Recherche lexicale : champ lexical de la nature et de la protection}

Repérez dans le texte tous les mots (noms, verbes, adjectifs) qui appartiennent au domaine sémantique de la \textbf{nature}, de la \textbf{forêt} et de la \textbf{protection de l'environnement}. Le tableau ci-dessous les classe avec l'article, le pluriel et la traduction : c'est un modèle de fiche de vocabulaire à reproduire dans votre cahier.

\begin{center}
\begin{tabularx}{\linewidth}{|X|X|X|X|}
\hline
\rowcolor{popgreenL}
\textbf{Mot (texte)} & \textbf{Article / Pluriel} & \textbf{Traduction} & \textbf{Catégorie} \\
\hline
Wälder & der Wald, die Wälder & la forêt & Nom (nature) \\
Abholzung & die Abholzung, die -en & la déforestation & Nom (action négative) \\
Holzindustrie & die Holzindustrie, die -n & l'industrie du bois & Nom (acteur) \\
Artenvielfalt & die Artenvielfalt (kein Pl.) & la biodiversité & Nom (nature) \\
Natur & die Natur, die -en & la nature & Nom (nature) \\
Aufforstung & die Aufforstung, die -en & le reboisement & Nom (action positive) \\
Bäume & der Baum, die Bäume & l'arbre & Nom (nature) \\
Naturpark & der Naturpark, die -s & le parc naturel & Nom (lieu protégé) \\
Ranger & der Ranger, die Ranger & l'éco-garde & Nom (acteur) \\
Tiere & das Tier, die Tiere & l'animal & Nom (nature) \\
Regenwald & der Regenwald, die Regenwälder & la forêt tropicale & Nom (nature) \\
Baumarten & die Baumart, die -en & l'essence d'arbre & Nom (nature) \\
Wasser & das Wasser (kein Pl.) & l'eau & Nom (ressource) \\
Boden & der Boden, die Böden & le sol & Nom (ressource) \\
Ökotourismus & der Ökotourismus (kein Pl.) & l'écotourisme & Nom (concept) \\
Öko-Lodges & die Öko-Lodge, die -n & l'éco-lodge & Nom (hébergement) \\
Umwelt & die Umwelt (kein Pl.) & l'environnement & Nom (milieu) \\
Dörfer & das Dorf, die Dörfer & le village & Nom (lieu) \\
nachhaltig & — & durable & Adjectif \\
schützen & — & protéger & Verbe \\
gefährden & — & menacer & Verbe \\
pflanzen & — & planter & Verbe \\
finanzieren & — & financer & Verbe \\
respektieren & — & respecter & Verbe \\
\hline
\end{tabularx}
\end{center}

\begin{aretenir}[Familles de mots : das Wortschatz-Netz]
Autour du verbe \cle{schützen} (protéger) se construit tout un réseau :
\begin{itemize}
    \item \all{der Schutz} (la protection) — \all{schützenswert} (qui mérite d'être protégé)
    \item \all{der Naturschutz} (la protection de la nature) — \all{der Umweltschutz} (la protection de l'environnement)
    \item \all{der Artenschutz} (la protection des espèces) — \all{der Klimaschutz} (la protection du climat)
\end{itemize}
Même chose pour \cle{pflanzen} : \all{die Pflanze, die -n} (la plante), \all{die Pflanzung, die -en} (la plantation), \all{anpflanzen} (planter, cultiver), \all{umpflanzen} (replanter, transplanter). \\
\textbf{Astuce :} en allemand, les mots composés se devinent : \all{der Natur} + \all{der Park} = \all{der Naturpark}. Le dernier élément donne toujours le genre et le sens principal.
\end{aretenir}

\begin{popfigure}
\begin{center}\tcbox[colback=popink,colframe=popink,coltext=white,arc=9pt,boxsep=4pt,left=6pt,right=6pt]{\bfseries\small Natur \& Umwelt}\end{center}
\vspace{-2pt}
\begin{tcbraster}[raster columns=2,raster equal height=rows,raster column skip=6pt,raster row skip=6pt,enhanced,blankest]
\begin{tcolorbox}[enhanced,colback=white,colframe=popblue!70,boxrule=0.9pt,arc=3pt,title={Wald (der Wald, die Wälder)},fonttitle=\bfseries\scriptsize,coltitle=white,colbacktitle=popblue!70,left=4pt,right=4pt,top=2pt,bottom=2pt,boxsep=1pt,toptitle=1.5pt,bottomtitle=1.5pt,halign=left]
\begin{itemize}[nosep,leftmargin=*,itemsep=1pt,font=\scriptsize]
\item Baum, die Bäume
\item Regenwald
\item Baumart, die -en
\item Aufforstung
\end{itemize}
\end{tcolorbox}
\begin{tcolorbox}[enhanced,colback=white,colframe=poporange!70,boxrule=0.9pt,arc=3pt,title={Schutz (der Schutz)},fonttitle=\bfseries\scriptsize,coltitle=white,colbacktitle=poporange!70,left=4pt,right=4pt,top=2pt,bottom=2pt,boxsep=1pt,toptitle=1.5pt,bottomtitle=1.5pt,halign=left]
\begin{itemize}[nosep,leftmargin=*,itemsep=1pt,font=\scriptsize]
\item Naturschutz
\item Artenschutz
\item Klimaschutz
\item Ranger
\end{itemize}
\end{tcolorbox}
\begin{tcolorbox}[enhanced,colback=white,colframe=poppurple!70,boxrule=0.9pt,arc=3pt,title={Tourismus (der Tourismus)},fonttitle=\bfseries\scriptsize,coltitle=white,colbacktitle=poppurple!70,left=4pt,right=4pt,top=2pt,bottom=2pt,boxsep=1pt,toptitle=1.5pt,bottomtitle=1.5pt,halign=left]
\begin{itemize}[nosep,leftmargin=*,itemsep=1pt,font=\scriptsize]
\item Ökotourismus
\item Naturpark
\item Wanderwege
\item Öko-Lodge
\end{itemize}
\end{tcolorbox}
\begin{tcolorbox}[enhanced,colback=white,colframe=poppink!70,boxrule=0.9pt,arc=3pt,title={Ressourcen},fonttitle=\bfseries\scriptsize,coltitle=white,colbacktitle=poppink!70,left=4pt,right=4pt,top=2pt,bottom=2pt,boxsep=1pt,toptitle=1.5pt,bottomtitle=1.5pt,halign=left]
\begin{itemize}[nosep,leftmargin=*,itemsep=1pt,font=\scriptsize]
\item Wasser
\item Boden, die Böden
\item Artenvielfalt
\item lokale Produkte
\end{itemize}
\end{tcolorbox}
\end{tcbraster}

\legende{Carte mentale du champ lexical « Nature, forêt, protection, écotourisme » issu du texte}
\end{popfigure}

\courssub{Réformulation : résumer le texte en trois phrases (Zusammenfassung)}

Résumer, c'est sélectionner l'essentiel, reformuler avec ses propres mots (ou ceux du texte) et lier les idées par des connecteurs logiques. Au niveau A1+/A2, on utilise le \textbf{Präsens} et des connecteurs simples.

\begin{propriete}[Règles d'un bon résumé (Zusammenfassung) à ce niveau]
\begin{itemize}
    \item \textbf{Temps :} Präsens — les faits sont actuels et généraux.
    \item \textbf{Longueur :} 3 phrases maximum (une par paragraphe logique du texte).
    \item \textbf{Connecteurs :} \all{und} (et), \all{aber} (mais), \all{weil} (parce que), \all{deshalb} (c'est pourquoi), \all{auch} (aussi).
    \item \textbf{Personne :} 3e personne ou forme impersonnelle \all{man}.
    \item \textbf{Pas de détails :} pas de chiffres précis, pas de citations, pas d'exemples.
    \item \textbf{Attention à l'ordre des mots :} après \all{deshalb} (1re position), le verbe vient en 2e position, puis le sujet : \all{Deshalb pflanzen die Schüler Bäume.}
\end{itemize}
\end{propriete}

\begin{exemplebox}[Modèle de résumé en 3 phrases]
\textbf{Satz 1 (Problème + action) :} \\
\all{In Kamerun gefährdet die Abholzung die Wälder, deshalb pflanzen Schüler, Freiwillige und das Ministerium gemeinsam neue Bäume.} \\
« Au Cameroun, la déforestation menace les forêts ; c'est pourquoi des élèves, des volontaires et le ministère plantent ensemble de nouveaux arbres. »

\textbf{Satz 2 (Solution concrète : parc + éco-gardes) :} \\
\all{Im Naturpark „Mont Cameroun“ schützen Ranger die Tiere und zeigen den Touristen die Bedeutung des Regenwaldes für Wasser und Boden.} \\
« Dans le parc naturel « Mont Cameroun », des éco-gardes protègent les animaux et montrent aux touristes l'importance de la forêt tropicale pour l'eau et le sol. »

\textbf{Satz 3 (Modèle durable : écotourisme) :} \\
\all{Der Ökotourismus finanziert durch den Parkeintritt die Aufforstung und schafft Einkommen für die Dörfer, ohne die Natur zu zerstören.} \\
« L'écotourisme finance le reboisement grâce au droit d'entrée du parc et crée des revenus pour les villages, sans détruire la nature. »
\end{exemplebox}

\begin{exoresolu}[Écrire son propre résumé — étape par étape]
\textbf{Consigne :} Rédigez votre propre résumé en 3 phrases en vous aidant du modèle ci-dessus.

\textbf{Étape 1 :} Identifiez les 3 blocs logiques du texte.
\begin{itemize}
    \item Bloc A (paragraphe 1) : problème (déforestation) + action collective (reboisement).
    \item Bloc B (paragraphe 2) : parc naturel, éco-gardes, sensibilisation des touristes.
    \item Bloc C (paragraphe 3) : définition de l'écotourisme, financement, développement durable.
\end{itemize}

\textbf{Étape 2 :} Notez les mots-clés de chaque bloc (en allemand).
\begin{itemize}
    \item A : \all{Abholzung, gefährden, pflanzen, Schüler, Freiwillige, Ministerium}
    \item B : \all{Naturpark, Ranger, schützen, führen, Besucher, Bedeutung, Wasser, Boden}
    \item C : \all{Ökotourismus, finanzieren, Eintritt, Aufforstung, Einkommen, Dörfer, nicht zerstören}
\end{itemize}

\textbf{Étape 3 :} Assemblez en 3 phrases avec des connecteurs.
\tcblower
\textbf{Solution possible :}
\begin{quote}
\all{In Kamerun gefährdet die Abholzung die Natur, aber Schulen und das Ministerium pflanzen neue Bäume. Im Naturpark schützen Ranger die Tiere und erklären den Touristen die Wichtigkeit des Waldes. Durch Ökotourismus wird die Aufforstung finanziert, und die Dörfer profitieren.}
\end{quote}
\textbf{Étape 4 (auto-contrôle) :} 3 phrases ? Präsens ? Connecteurs ? Majuscules aux noms ? Verbe conjugué en 2e position dans chaque phrase principale ?
\end{exoresolu}

\courssub{Expression orale guidée : « Und bei uns in Kamerun? »}

La dernière étape de la compréhension est l'appropriation personnelle : relier le texte à sa propre réalité. Ici, le texte parle \emph{du} Cameroun — l'élève est donc \emph{dans} le texte. L'activité vise à faire émerger un discours personnel, argumenté, en réutilisant le vocabulaire et les structures vus.

\begin{experience}[Débat guidé : Reboisement et écotourisme chez nous]
\textbf{Phase 1 — Préparation individuelle (5 min) :} chaque élève note 3 idées en allemand (mots-clés ou phrases courtes) sur une fiche.
\begin{itemize}
    \item \all{Was gibt es schon?} (Qu'est-ce qui existe déjà ?) — \textit{ex. : Parc national de Waza, reboisement scolaire, éco-gardes}
    \item \all{Was fehlt?} (Qu'est-ce qui manque ?) — \textit{ex. : plus de pépinières, meilleurs sentiers, formation de guides}
    \item \all{Was kann ich tun?} (Que puis-je faire ?) — \textit{ex. : planter un arbre, visiter un parc, sensibiliser}
\end{itemize}

\textbf{Phase 2 — Mise en commun en binôme (5 min) :} en paires, les élèves comparent leurs notes et forment 4 à 5 phrases complètes. Ils s'aident des \textbf{phrases-outils} au tableau :
\begin{center}
\begin{tabularx}{\linewidth}{|X|X|}
\hline
\rowcolor{popblueL}
\textbf{Deutsch} & \textbf{Français} \\
\hline
Bei uns in Kamerun gibt es … & Chez nous au Cameroun, il y a … \\
Das ist gut, weil … & C'est bien, parce que … \\
Aber es fehlt … & Mais il manque … \\
Wir brauchen mehr … & Nous avons besoin de plus de … \\
Ich kann … & Je peux … \\
Wir sollten … & Nous devrions … \\
\hline
\end{tabularx}
\end{center}

\textbf{Phase 3 — Présentation orale (tour de classe, 10 min) :} chaque binôme présente une idée. Le professeur note au tableau les structures récurrentes pour les corriger collectivement (ex. : \all{Wir brauchen mehr Bäume} — accusatif pluriel !).

\textbf{Phase 4 — Trace écrite (devoir) :} rédiger un court paragraphe (5 à 6 lignes) : \emph{« Meine Meinung zum Ökotourismus in Kamerun »}.
\end{experience}

\begin{savaistu}[Le Cameroun et l'Allemagne : une coopération verte]
Le Cameroun et l'Allemagne coopèrent depuis longtemps dans la protection des forêts du bassin du Congo : des programmes allemands d'aide au développement soutiennent la gestion durable des forêts et la formation des éco-gardes, en collaboration avec le MINFOF (Ministère camerounais des Forêts et de la Faune). La région du \textbf{Mont Cameroun}, près de Buea, est une zone pilote pour l'écotourisme communautaire : une partie des revenus générés par les visites revient aux villages riverains, ce qui incite les habitants à protéger la forêt plutôt qu'à la couper. C'est exactement le principe du texte : \all{Schutz durch nachhaltige Nutzung} — « la protection par l'utilisation durable ». En allemand, le mot \all{nachhaltig} (durable) vient d'ailleurs de la foresterie : on ne devait couper que ce que la forêt pouvait « faire durer » (\all{nachhalten}).
\end{savaistu}

\begin{popfigure}
\begin{tikzpicture}[
    node distance=1.2cm and 2.5cm,
    box/.style={draw=popblue, thick, rounded corners, fill=popblueL, text width=3.2cm, align=center, minimum height=1.6cm},
    arrow/.style={-Stealth, thick, popdark}
]
\node[box, align=center] (glob){\textbf{Compréhension globale}\\\textit{Wer? Wo? Was? Warum?}\\\textit{4 questions essentielles}};
\node[box, right=of glob, align=center] (det){\textbf{Compréhension détaillée}\\\textit{6 questions ciblées}\\\textit{Phrases complètes + Beleg}};
\node[box, right=of det, align=center] (vf){\textbf{Vrai / Faux}\\\textit{Richtig / Falsch}\\\textit{Justification par citation}};
\node[box, below=of det, align=center] (lex){\textbf{Recherche lexicale}\\\textit{Champ lexical nature}\\\textit{Tableau + familles de mots}};
\node[box, below=of vf, align=center] (sum){\textbf{Résumé (3 Sätze)}\\\textit{Präsens + connecteurs}\\\textit{und, aber, weil, deshalb}};
\node[box, below=of sum, align=center] (oral){\textbf{Expression orale}\\\textit{« Und bei uns? »}\\\textit{Binôme → Classe → Écrit}};

\draw[arrow] (glob) -- (det);
\draw[arrow] (det) -- (vf);
\draw[arrow] (det) -- (lex);
\draw[arrow] (lex) -- (sum);
\draw[arrow] (vf) -- (sum);
\draw[arrow] (sum) -- (oral);
\end{tikzpicture}
\legende{Organigramme de la progression didactique : du global à l'oral personnel}
\end{popfigure}

\courssub{Exercices de consolidation}

\begin{exercice}[Compréhension et expression — Entraînement]
\begin{enumerate}
    \item Répondez en allemand par une phrase complète : \all{Warum pflanzt man in Kamerun neue Bäume?}
    \item Citez du texte la phrase qui définit l'écotourisme.
    \item Transformez au pluriel : \all{Der Ranger schützt den Baum.} → \all{\_\_\_\_\_\_\_\_\_\_\_\_}
    \item Complétez avec la préposition correcte suivie de l'accusatif (\all{durch}, \all{für}, \all{ohne}, \all{gegen}) : \\
    \all{Die Abholzung \_\_\_\_\_\_\_\_\_ die Holzindustrie gefährdet die Natur.} \\
    \all{Wir pflanzen Bäume \_\_\_\_\_\_\_\_\_ die Zukunft.} \\
    \all{Ökotourismus heißt: reisen \_\_\_\_\_\_\_\_\_ die Natur zu zerstören.} \\
    \all{Der Park schützt die Tiere \_\_\_\_\_\_\_\_\_ die Wilderei.}
    \item Mini-production (3 phrases) : décrivez une action de reboisement dans votre école ou votre quartier. Utilisez : \all{pflanzen}, \all{Schüler}, \all{Bäume}, \all{gemeinsam}, \all{weil}.
\end{enumerate}
\end{exercice}

\begin{corrige}[Exercice : Compréhension et expression]
\begin{enumerate}
    \item \all{Man pflanzt neue Bäume, weil die Abholzung die Natur gefährdet.} \\
    « On plante de nouveaux arbres parce que la déforestation menace la nature. » (Attention : après \all{weil}, le verbe conjugué va à la fin.)
    \item „Ökotourismus bedeutet: Reisen ohne die Natur zu zerstören.“
    \item \all{Die Ranger schützen die Bäume.} \\
    (Rappel : \all{der Ranger} a un pluriel identique ; \all{der Baum} → \all{die Bäume} ; le verbe passe à la 3e personne du pluriel : \all{schützen}.)
    \item \all{durch} (la cause : \textit{à travers / par}) — \all{für} (le but : \textit{pour}) — \all{ohne} (l'absence : \textit{sans}) — \all{gegen} (l'opposition : \textit{contre}). \\
    Phrases complètes : \all{Die Abholzung durch die Holzindustrie gefährdet die Natur.} / \all{Wir pflanzen Bäume für die Zukunft.} / \all{Ökotourismus heißt: reisen ohne die Natur zu zerstören.} / \all{Der Park schützt die Tiere gegen die Wilderei.}
    \item \textit{Exemple de production :} \all{In unserer Schule pflanzen wir gemeinsam Bäume. Das machen wir, weil der Wald wichtig für das Klima ist. Jeder Schüler pflanzt einen Baum.} \\
    « Dans notre école, nous plantons des arbres ensemble. Nous le faisons parce que la forêt est importante pour le climat. Chaque élève plante un arbre. »
\end{enumerate}
\end{corrige}

\coursec{Lexique thématique : la forêt et la protection de l'environnement}

Après avoir lu et compris le texte „Aufforstung und Ökotourismus in Kamerun“, nous allons maintenant fixer le vocabulaire essentiel. Ce lexique thématique est structuré en quatre champs sémantiques pour faciliter la mémorisation : les noms (avec article et pluriel), les verbes d'action, les termes liés à l'environnement et au tourisme, ainsi que le vocabulaire des problèmes et des solutions. Nous observerons aussi la formation des mots composés, très productive en allemand, et les familles de mots qui permettent d'enrichir son vocabulaire de manière exponentielle.

\courssub{Champ 1 : La forêt et les arbres (Der Wald und die Bäume)}

Commençons par les noms concrets désignant les éléments de la forêt. En allemand, \cle{tout nom s'apprend avec son article (genre) et son pluriel}. Observez les mutations vocaliques (Umlaut) fréquentes au pluriel pour les mots d'une syllabe.

\begin{center}
\begin{tabularx}{\linewidth}{|>{\bfseries}l|X|X|l|}
\hline
\rowcolor{popblueL}
\textbf{Allemand (Article + Pluriel)} & \textbf{Français} & \textbf{Exemple de phrase} & \textbf{Traduction} \\
\hline
der Wald, die \cle{Wälder} & la forêt & Der Wald in Kamerun ist groß. & La forêt au Cameroun est grande. \\
\hline
der Baum, die \cle{Bäume} & l'arbre & Wir pflanzen neue Bäume. & Nous plantons de nouveaux arbres. \\
\hline
der Ast, die \cle{Äste} & la branche & Der Ast bricht im Wind. & La branche casse dans le vent. \\
\hline
das Blatt, die \cle{Blätter} & la feuille & Die Blätter sind grün. & Les feuilles sont vertes. \\
\hline
die Wurzel, die \cle{Wurzeln} & la racine & Die Wurzeln halten die Erde. & Les racines retiennent la terre. \\
\hline
der Sämling, die \cle{Sämlinge} & le semis, le plant & Der Sämling wächst schnell. & Le plant pousse vite. \\
\hline
die Pflanze, die \cle{Pflanzen} & la plante & Diese Pflanze ist selten. & Cette plante est rare. \\
\hline
\end{tabularx}
\end{center}

\begin{propriete}[Familles de mots : der Baum]
L'allemand construit beaucoup de mots composés (\cle{Komposita}). À partir de \all{der Baum}, on forme :
\begin{itemize}
    \item \all{der Obstbaum, die Obstbäume} (l'arbre fruitier) — \all{der Baumstamm, die Baumstämme} (le tronc d'arbre).
    \item \all{der Baumsaft} (la sève), \all{der Baumschatten} (l'ombre de l'arbre).
\end{itemize}
\emph{Règle d'or :} le genre et le pluriel du composé sont ceux du \cle{dernier élément} (\all{der Baum} $\rightarrow$ \all{der Obstbaum}).
\end{propriete}

\begin{exemplebox}[Exemples contextualisés]
\all{„Der tropische Wald schützt das Klima.“} (« La forêt tropicale protège le climat. ») \\
\all{„Ein alter Baum hat viele Äste und Blätter.“} (« Un vieil arbre a beaucoup de branches et de feuilles. ») \\
\all{„Die Wurzeln der Bäume verhindern die Erosion.“} (« Les racines des arbres empêchent l'érosion. »)
\end{exemplebox}

\courssub{Champ 2 : Les actions — Verbes de l'écologie (Handlungen)}

Voici les verbes clés du texte. Notez les verbes à particule séparable (\all{abholzen}, \all{aufforsten}) et les verbes forts avec modification de voyelle au présent (\all{wachsen} $\rightarrow$ \all{er wächst}, \all{gießen} $\rightarrow$ \all{er gießt}).

\begin{center}
\begin{tabularx}{\linewidth}{|>{\bfseries}l|l|X|l|}
\hline
\rowcolor{popgreenL}
\textbf{Infinitif} & \textbf{Présent (ich/er)} & \textbf{Particularité} & \textbf{Français} \\
\hline
pflanzen & ich pflanze & régulier & planter \\
\hline
schützen & ich schütze & régulier & protéger \\
\hline
\cle{wachsen} & \cle{er wächst} & \textbf{fort (a $\rightarrow$ ä)} & \textbf{pousser, croître} \\
\hline
säen & ich säe & régulier & semer \\
\hline
\cle{gießen} & \cle{er gießt} & \textbf{fort (ß, pas de changement de voyelle)} & \textbf{arroser} \\
\hline
\cle{abholzen} & \cle{er holzt ab} & \textbf{séparable} & \textbf{abattre, déboiser} \\
\hline
aufforsten & er forstet auf & séparable & reboiser \\
\hline
recyceln & ich recycle & régulier (emprunt) & recycler \\
\hline
sparen & ich spare & régulier & économiser \\
\hline
\end{tabularx}
\end{center}

\begin{attention}[Conjugaison : deux précisions importantes]
\begin{itemize}
    \item \all{gießen} est un verbe \textbf{fort} (Präteritum \all{goss}, Partizip II \all{gegossen}), mais sa voyelle \textbf{ne change pas} au présent : \all{ich gieße, du gießt, er gießt}. Attention à l'orthographe : \textbf{ß} après voyelle longue.
    \item \all{aufforsten} : le radical se termine par \textbf{-t}, donc on ajoute un \textbf{-e-} de liaison : \all{du forstest auf, er forstet auf} (comme \all{arbeiten} $\rightarrow$ \all{er arbeitet}).
\end{itemize}
\end{attention}

\begin{propriete}[Familles de mots verbales : la nominalisation]
En allemand, un verbe peut devenir un nom d'action : il prend la majuscule et, très souvent, le suffixe \all{-ung} (toujours féminin : \all{die ...ung, die ...ungen}).
\begin{itemize}
    \item \all{pflanzen} (v.) $\rightarrow$ \all{die Pflanzung} (la plantation) ; \all{die Pflanze} (la plante).
    \item \all{schützen} (v.) $\rightarrow$ \all{der Schutz} (la protection) $\rightarrow$ \all{der Umweltschutz} (la protection de l'environnement).
    \item \all{abholzen} (v.) $\rightarrow$ \all{die Abholzung} (la déforestation).
    \item \all{aufforsten} (v.) $\rightarrow$ \all{die Aufforstung} (le reboisement).
    \item \all{verschmutzen} (v., salir, polluer) $\rightarrow$ \all{die Verschmutzung} (la pollution).
\end{itemize}
\emph{Astuce :} repérer le suffixe \all{-ung} permet de deviner un nom d'action féminin et de traduire par un nom français en « -tion » ou « -ment ».
\end{propriete}

\begin{exemplebox}[Verbes en action]
\all{„Wir \textbf{pflanzen} heute hundert Sämlinge.“} (« Nous plantons cent plants aujourd'hui. ») \\
\all{„Der Naturpark \textbf{schützt} die Tiere und Pflanzen.“} (« Le parc naturel protège les animaux et les plantes. ») \\
\all{„Die Bäume \textbf{wachsen} hier sehr schnell.“} (« Les arbres poussent très vite ici. ») \\
\all{„Man darf den Wald nicht \textbf{abholzen}.“} (« On ne doit pas déboiser la forêt. ») \\
\all{„Die Gemeinde \textbf{forstet} die kahlen Hügel \textbf{auf}.“} (« La commune reboise les collines nues. »)
\end{exemplebox}

\courssub{Champ 3 : L'environnement et le tourisme (Umwelt und Tourismus)}

Ce champ rassemble les noms abstraits et les acteurs. Attention aux mots composés longs : on les lit de droite à gauche pour en comprendre le sens.

\begin{center}
\begin{tabularx}{\linewidth}{|>{\bfseries}l|X|l|}
\hline
\rowcolor{poporangeL}
\textbf{Allemand (Article + Pluriel)} & \textbf{Français} & \textbf{Analyse / Composition} \\
\hline
die Umwelt, — & l'environnement & um + Welt (= le monde autour) \\
\hline
der Umweltschutz, — & la protection de l'environnement & Umwelt + \cle{Schutz} (der Schutz) \\
\hline
der Naturpark, die Naturparks & le parc naturel & Natur + Park \\
\hline
der Ökotourismus, — & l'écotourisme & Öko- (grec \emph{oikos}) + Tourismus \\
\hline
der Tourist, die Touristen & le touriste & (nom faible : \all{den Touristen}) \\
\hline
die Natur, — & la nature & (pas de pluriel usuel) \\
\hline
die Landschaft, die Landschaften & le paysage & Land + -schaft \\
\hline
\end{tabularx}
\end{center}

\begin{definition}[Le mot composé : der Umweltschutz]
\all{der Umweltschutz} = \all{die Umwelt} (l'environnement) + \all{der Schutz} (la protection).
C'est un nom composé typique. Il s'écrit \textbf{en un seul mot}, sans espace ni trait d'union. Le genre est celui du second élément : \all{der Schutz} $\rightarrow$ \all{der Umweltschutz}.
\end{definition}

\begin{propriete}[Règle des noms composés (Komposita)]
\begin{enumerate}
    \item Le \cle{dernier élément} (le noyau) détermine le \textbf{genre} et le \textbf{pluriel}.
    \item Le(s) premier(s) élément(s) (déterminants) précisent le sens.
    \item Écriture : \textbf{un seul mot}, sans espace, sans trait d'union.
    \item Une lettre de liaison (\cle{Fugenlaut}) peut s'insérer : \textbf{-s-} (\all{die Baumwurzel}), \textbf{-n-/-en-} (\all{der Touristenführer}).
\end{enumerate}
\emph{Exemples du thème :} \all{die Baumwurzel} (die Wurzel), \all{der Obstbaum} (der Baum), \all{der Landschaftsschutz} (der Schutz), \all{der Ökotourist} (der Tourist).
\end{propriete}

\begin{attention}[Faux amis et pièges fréquents]
\begin{itemize}
    \item \all{die Natur} = la nature. \all{natürlich} = \textbf{1. naturel (adjectif) ; 2. naturellement, bien sûr (adverbe)}. Ne pas traduire systématiquement par « naturellement ».
    \item \all{der Akt} = l'acte (juridique, administratif), \textbf{pas} « l'action » au sens général. Pour « une action (écologique) », utilisez \all{die Aktion} ou \all{die Maßnahme}.
    \item \all{die Pflanzung} = la plantation (action de planter ; aussi « plantation » agricole). Ne pas inventer *\all{die Plantage} sur le modèle français.
    \item \all{recyceln} s'écrit avec \textbf{y} et \textbf{c}, pas *recyclen.
    \item \all{der Tourist} est un \textbf{nom faible} : \all{den Touristen, dem Touristen} (accusatif et datif singulier).
\end{itemize}
\end{attention}

\courssub{Champ 4 : Problèmes et solutions (Probleme und Lösungen)}

Vocabulaire pour argumenter, décrire des menaces et proposer des remèdes. Rappel : les adjectifs comme \all{nachhaltig}, \all{sauber}, \all{schmutzig} sont \textbf{invariables} en position d'attribut (après \all{sein} : \all{das Wasser ist sauber}) mais \textbf{s'accordent} en position épithète (après article : \all{das saubere Wasser}).

\begin{center}
\begin{tabularx}{\linewidth}{|>{\bfseries}l|X|l|}
\hline
\rowcolor{poppurpleL}
\textbf{Allemand (Article + Pluriel / Forme)} & \textbf{Français} & \textbf{Collocations utiles} \\
\hline
die Abholzung, die Abholzungen & la déforestation & die illegale Abholzung stoppen \\
\hline
die Verschmutzung, die Verschmutzungen & la pollution & die Luftverschmutzung, die Wasserverschmutzung \\
\hline
das Problem, die Probleme & le problème & ein großes Problem lösen \\
\hline
die Lösung, die Lösungen & la solution & eine nachhaltige Lösung finden \\
\hline
nachhaltig (adj.) & durable, soutenable & nachhaltige Entwicklung \\
\hline
sauber (adj.) & propre & sauberes Wasser, saubere Luft \\
\hline
schmutzig (adj.) & sale & schmutziges Wasser \\
\hline
\end{tabularx}
\end{center}

\begin{exemplebox}[Expressions utiles pour l'oral et l'écrit]
Ces collocations (associations de mots figées) sont à mémoriser par cœur :
\begin{itemize}
    \item \all{die Umwelt schützen} (protéger l'environnement)
    \item \all{Bäume pflanzen} (planter des arbres)
    \item \all{in der Natur wandern} (randonner dans la nature)
    \item \all{Müll trennen} (trier les déchets) — \emph{très important en Allemagne !}
    \item \all{Energie sparen} (économiser l'énergie)
    \item \all{den Wald aufforsten} (reboiser la forêt)
    \item \all{Plastik vermeiden} (éviter le plastique)
\end{itemize}
\end{exemplebox}

\begin{savaistu}[Le tri des déchets en Allemagne : die Mülltrennung]
En Allemagne, la protection de l'environnement commence devant la porte : \all{die Mülltrennung} (le tri des déchets) est obligatoire et très codifiée. Chaque foyer a plusieurs poubelles :
\begin{itemize}
    \item \textbf{Gelbe Tonne / Gelber Sack} (jaune) : emballages en plastique, métal, matériaux composites (\all{Verpackungen}).
    \item \textbf{Blaue Tonne} (bleue) : papier et carton (\all{Papier, Pappe}).
    \item \textbf{Braune Tonne / Biotonne} (brune) : déchets organiques (\all{Bioabfälle}).
    \item \textbf{Schwarze / graue Tonne} (noire/grise) : déchets résiduels (\all{Restmüll}).
    \item \textbf{Glascontainer} (conteneurs à verre) : tri par couleur (blanc, vert, brun).
\end{itemize}
Mal trier peut coûter une amende (\all{das Bußgeld}). C'est une habitude culturelle forte : \all{„Müll trennen ist Pflicht.“} (« Trier les déchets est un devoir. »)
\end{savaistu}

\courssub{Carte mentale récapitulative : « die Umwelt »}

La figure ci-dessous résume visuellement les quatre champs lexicaux. Utilisez-la pour réviser en reliant les mots entre eux (ex. : \all{der Wald} $\rightarrow$ \all{schützen} $\rightarrow$ \all{der Umweltschutz}).

\begin{popfigure}
\begin{center}\tcbox[colback=popdark,colframe=popdark,coltext=white,arc=9pt,boxsep=4pt,left=6pt,right=6pt]{\bfseries\small die Umwelt}\end{center}
\vspace{-2pt}
\begin{tcbraster}[raster columns=2,raster equal height=rows,raster column skip=6pt,raster row skip=6pt,enhanced,blankest]
\begin{tcolorbox}[enhanced,colback=white,colframe=popgreen,boxrule=0.9pt,arc=3pt,title={Wald},fonttitle=\bfseries\scriptsize,coltitle=white,colbacktitle=popgreen,left=4pt,right=4pt,top=2pt,bottom=2pt,boxsep=1pt,toptitle=1.5pt,bottomtitle=1.5pt,halign=left]
\begin{itemize}[nosep,leftmargin=*,itemsep=1pt,font=\scriptsize]
\item der Wald
\item der Baum
\item die Wurzel
\item das Blatt
\item der Sämling
\end{itemize}
\end{tcolorbox}
\begin{tcolorbox}[enhanced,colback=white,colframe=poporange,boxrule=0.9pt,arc=3pt,title={Aktionen},fonttitle=\bfseries\scriptsize,coltitle=white,colbacktitle=poporange,left=4pt,right=4pt,top=2pt,bottom=2pt,boxsep=1pt,toptitle=1.5pt,bottomtitle=1.5pt,halign=left]
\begin{itemize}[nosep,leftmargin=*,itemsep=1pt,font=\scriptsize]
\item pflanzen
\item schützen
\item aufforsten
\item recyceln
\item sparen
\item abholzen
\end{itemize}
\end{tcolorbox}
\begin{tcolorbox}[enhanced,colback=white,colframe=poppurple,boxrule=0.9pt,arc=3pt,title={Tourismus},fonttitle=\bfseries\scriptsize,coltitle=white,colbacktitle=poppurple,left=4pt,right=4pt,top=2pt,bottom=2pt,boxsep=1pt,toptitle=1.5pt,bottomtitle=1.5pt,halign=left]
\begin{itemize}[nosep,leftmargin=*,itemsep=1pt,font=\scriptsize]
\item der Naturpark
\item der Ökotourismus
\item der Tourist
\item wandern
\item die Landschaft
\end{itemize}
\end{tcolorbox}
\begin{tcolorbox}[enhanced,colback=white,colframe=poppink,boxrule=0.9pt,arc=3pt,title={Probleme \& Lösungen},fonttitle=\bfseries\scriptsize,coltitle=white,colbacktitle=poppink,left=4pt,right=4pt,top=2pt,bottom=2pt,boxsep=1pt,toptitle=1.5pt,bottomtitle=1.5pt,halign=left]
\begin{itemize}[nosep,leftmargin=*,itemsep=1pt,font=\scriptsize]
\item die Abholzung
\item die Verschmutzung
\item das Problem
\item die Lösung
\item nachhaltig
\end{itemize}
\end{tcolorbox}
\end{tcbraster}

\legende{Carte mentale du lexique « Umwelt » — Reliez les branches pour créer des phrases (ex. : Tourist + wandern + in + Naturpark).}
\end{popfigure}

\courssub{Texte d'entraînement : « Grüne Zukunft für das Dja-Reservat »}

Le texte suivant réemploie le vocabulaire des quatre champs. Lisez-le à haute voix, soulignez les mots connus, déduisez les inconnus par le contexte et la transparence (mots composés, emprunts).

\begin{textebox}[Texte original — Niveau A2]
„Im Südosten Kameruns liegt das Dja-Reservat, ein großer Naturpark und UNESCO-Weltnaturerbe. Der tropische Wald hier ist einzigartig: riesige Bäume, seltene Pflanzen und viele Tiere wie Gorillas und Waldelefanten finden hier Schutz. Doch die Abholzung durch illegale Holzfäller und die Verschmutzung durch Plastikmüll sind große Probleme für die Umwelt.

Eine lokale Organisation startet jetzt ein Projekt zur Aufforstung. „Wir wollen die kahlen Flächen wieder bewalden“, sagt der Projektleiter. „Jeden Monat pflanzen wir tausend Sämlinge einheimischer Baumarten. Die Dorfbewohner helfen beim Pflanzen und Gießen. So schaffen wir Arbeit und schützen gleichzeitig die Natur.“

Der Ökotourismus spielt auch eine wichtige Rolle. Touristen wandern mit Führern durch den Wald, beobachten Tiere und lernen viel über den Umweltschutz. Das Geld aus dem Tourismus finanziert neue Pflanzungen und Schulen. „Nachhaltigkeit bedeutet: Wir nutzen den Wald, aber wir zerstören ihn nicht“, erklärt ein Guide. „Sauberes Wasser und saubere Luft sind die Basis für unsere Zukunft.“

Die Regierung unterstützt das Projekt. Neue Gesetze verbieten das illegale Abholzen. In den Schulen lernen die Kinder: „Müll trennen, Energie sparen, Bäume pflanzen – das kann jeder.““
\end{textebox}

\begin{center}
\begin{tabularx}{\linewidth}{|>{\bfseries}l|X|}
\hline
\rowcolor{popblueL}
\textbf{Mot difficile (allemand)} & \textbf{Glossaire français} \\
\hline
das Dja-Reservat, — & la réserve du Dja \\
\hline
das Weltnaturerbe, — & le patrimoine naturel mondial \\
\hline
einzigartig & unique \\
\hline
der Waldelefant, -en & l'éléphant de forêt \\
\hline
der Holzfäller, - & le bûcheron \\
\hline
die kahle Fläche, -n & la zone dénudée \\
\hline
bewalden & boiser, couvrir d'arbres \\
\hline
einheimisch & indigène, local \\
\hline
der Dorfbewohner, - & l'habitant du village \\
\hline
gleichzeitig & simultanément \\
\hline
der Führer, - & le guide \\
\hline
beobachten & observer \\
\hline
finanzieren & financer \\
\hline
zerstören & détruire \\
\hline
die Basis, die Basen & la base \\
\hline
verbieten, verbot, hat verboten & interdire \\
\hline
\end{tabularx}
\end{center}

\courssub{Exercice résolu : Construction de mots composés et choix lexical}

\begin{exoresolu}[Former des noms composés et compléter]
\textbf{Énoncé :}
1. Formez le nom composé correspondant à la définition en utilisant les éléments donnés.
   a) Protection de la nature : Natur + Schutz $\rightarrow$ \dots
   b) Racine d'arbre : Baum + Wurzel $\rightarrow$ \dots
   c) Guide touristique : Tourist + Führer $\rightarrow$ \dots
   d) Pollution de l'air : Luft + Verschmutzung $\rightarrow$ \dots

2. Complétez les phrases avec le verbe approprié, conjugué au présent : \textit{pflanzen, schützen, wachsen, gießen, abholzen, aufforsten}.
   a) Der Förster \_\_\_\_\_\_ junge Bäume im Wald.
   b) Die Bäume \_\_\_\_\_\_ im Regenwald sehr schnell.
   c) Wir müssen die Umwelt \_\_\_\_\_\_.
   d) Bitte \_\_\_\_\_\_ Sie die Blumen jeden Tag!
   e) Die Firma \_\_\_\_\_\_ illegal tropisches Holz \_\_\_\_\_\_.
   f) Nach dem Brand \_\_\_\_\_\_ wir den Berg \_\_\_\_\_\_.

\tcblower
\textbf{Solution détaillée :}
1. \textbf{Noms composés (genre = dernier élément) :}
   a) \all{der Naturschutz} (der Schutz) — \emph{ici pas de lettre de liaison.}
   b) \all{die Baumwurzel} (die Wurzel) — \emph{lettre de liaison -s- : Baum + s + Wurzel.}
   c) \all{der Touristenführer} (der Führer) — \emph{liaison -en- : Tourist + en + Führer.}
   d) \all{die Luftverschmutzung} (die Verschmutzung) — \emph{sans liaison.}

2. \textbf{Conjugaison au présent :}
   a) Der Förster \textbf{pflanzt} junge Bäume im Wald. (verbe régulier)
   b) Die Bäume \textbf{wachsen} im Regenwald sehr schnell. (verbe fort : a $\rightarrow$ ä seulement aux 2\textsuperscript{e}/3\textsuperscript{e} pers. du singulier ; au \textbf{pluriel}, forme de l'infinitif : \all{sie wachsen})
   c) Wir müssen die Umwelt \textbf{schützen}. (infinitif après le modal \all{müssen})
   d) Bitte \textbf{gießen} Sie die Blumen jeden Tag! (\textbf{impératif formel} = infinitif + \all{Sie})
   e) Die Firma \textbf{holzt} illegal tropisches Holz \textbf{ab}. (verbe \textbf{séparable} : \all{abholzen} $\rightarrow$ \all{holzt ... ab})
   f) Nach dem Brand \textbf{forsten} wir den Berg \textbf{auf}. (verbe \textbf{séparable} : \all{aufforsten} $\rightarrow$ \all{forsten ... auf} ; la forme \all{wir} = infinitif)
\end{exoresolu}

\courssub{Exercices d'application (à faire en classe ou à la maison)}

\begin{exercice}[Lexique : Mots composés, familles, traduction]
\textbf{1. Familles de mots :} Complétez le tableau en retrouvant le nom d'action (souvent en \textit{-ung}) ou le nom d'agent.
\begin{center}
\begin{tabular}{|l|l|l|}
\hline
\rowcolor{popblueL}
\textbf{Verbe} & \textbf{Nom d'action} & \textbf{Agent (der/die ...er/in)} \\
\hline
pflanzen & \dots & der Pflanzer / die Pflanzerin \\
\hline
schützen & \dots & \dots \\
\hline
abholzen & \dots & \dots \\
\hline
aufforsten & \dots & \dots \\
\hline
wandern & \dots & \dots \\
\hline
\end{tabular}
\end{center}

\textbf{2. Traduisez en allemand (vocabulaire de la leçon) :}
a) La déforestation détruit la forêt. \\
b) Nous plantons des arbres pour l'avenir. \\
c) Le parc naturel protège les animaux. \\
d) L'écotourisme est durable. \\
e) Il faut trier les déchets et économiser l'énergie. \\
f) L'eau sale est un problème pour la santé.

\textbf{3. Expression guidée :} En 5-6 phrases, décrivez une action de reboisement dans votre région (Yaoundé, Douala, village natal). Utilisez au moins 8 mots du lexique (noms, verbes, adjectifs) et deux mots composés.
\end{exercice}

\begin{corrige}[Exercice 3 — Lexique]
\textbf{1. Familles de mots :}
\begin{center}
\begin{tabular}{|l|l|l|}
\hline
\rowcolor{popblueL}
\textbf{Verbe} & \textbf{Nom d'action} & \textbf{Agent} \\
\hline
pflanzen & die Pflanzung & der Pflanzer / die Pflanzerin \\
\hline
schützen & der Schutz & der Schützer / die Schützerin \\
\hline
abholzen & die Abholzung & der Holzfäller / die Holzfällerin \\
\hline
aufforsten & die Aufforstung & der Aufforster / die Aufforsterin \\
\hline
wandern & die Wanderung & der Wanderer / die Wanderin \\
\hline
\end{tabular}
\end{center}

\textbf{2. Traduction :}
a) \all{Die Abholzung zerstört den Wald.} \\
b) \all{Wir pflanzen Bäume für die Zukunft.} \\
c) \all{Der Naturpark schützt die Tiere.} \\
d) \all{Der Ökotourismus ist nachhaltig.} \\
e) \all{Man muss Müll trennen und Energie sparen.} (\textit{ou :} \all{Wir müssen Müll trennen und Energie sparen.}) \\
f) \all{Schmutziges Wasser ist ein Problem für die Gesundheit.}

\textbf{3. Production type (exemple) :}
\all{„In meiner Stadt Yaoundé gibt es ein Projekt zur Aufforstung. Wir pflanzen viele Sämlinge an den Hängen. Die Dorfbewohner schützen die jungen Bäume. Wir gießen sie in der Trockenzeit. Der Naturschutz ist wichtig für saubere Luft. Der Ökotourismus bringt Geld für die Schulen. Das ist eine nachhaltige Lösung.“}
\end{corrige}

\begin{methode}[Comment apprendre ce lexique efficacement ?]
\begin{enumerate}
    \item \textbf{Fiches « Mot + Article + Pluriel » :} écrivez \all{der Wald, die Wälder} au recto, « la forêt » au verso. Révisez 5 minutes par jour.
    \item \textbf{Arbre lexical :} dessinez l'arbre \all{Baum} avec ses branches \all{Obstbaum, Baumstamm, Baumwurzel, Baumsaft}. Faites de même pour \all{Schutz} $\rightarrow$ \all{Umweltschutz, Naturschutz, Landschaftsschutz}.
    \item \textbf{Verbes à afficher :} collez au mur les 9 verbes du Champ 2 avec leur particularité (séparable, fort, régulier). Conjuguez-les mentalement chaque matin.
    \item \textbf{Collocations :} apprenez les 7 expressions de la boîte « Expressions utiles » comme des blocs indissociables (\all{Müll trennen}, \all{Energie sparen}).
    \item \textbf{Contexte camerounais :} reliez chaque mot à une réalité locale (ex. \all{der Sämling} $\rightarrow$ pépinière de Mbalmayo ; \all{der Naturpark} $\rightarrow$ parc de la Lobéké ou de Korup ; \all{das Dja-Reservat} $\rightarrow$ réserve du Dja, patrimoine mondial de l'UNESCO).
\end{enumerate}
\end{methode}

\coursec{Grammaire 1 : l'impératif}

Dans la section précédente, nous avons découvert le vocabulaire de la forêt et de la protection de l'environnement : \all{der Wald}, \all{der Baum}, \all{pflanzen}, \all{schützen}, \all{die Umwelt}. Or, dans le texte \all{„Aufforstung und Ökotourismus in Kamerun“}, ces mots apparaissent souvent dans des phrases qui donnent des ordres ou des conseils : \all{„Pflanzt Bäume!“}, \all{„Geht mit uns!“}, \all{„Schützen Sie die Umwelt!“}. Ces phrases sont à \cle{l'impératif}. C'est le mode que l'on utilise pour agir sur les autres : c'est donc le mode de toutes les campagnes de protection de la nature !

\courssub{Observation : à quoi reconnaît-on un impératif ?}

\begin{textebox}[Extraits du texte support]
„Pflanzt Bäume in euren Dörfern!“\\
„Geht mit uns in den Naturpark!“\\
„Schützen Sie die Umwelt, schützen Sie unsere Zukunft!“\\
„Komm doch mit nach Yaoundé, dort beginnt unsere Tour!“
\end{textebox}

Observons ces quatre phrases :

\begin{itemize}
\item Le \cle{verbe} se trouve en \textbf{première position} de la phrase (comme dans une question en allemand).
\item Il n'y a \textbf{pas de pronom sujet} devant le verbe (sauf \all{Sie} et \all{wir}, placés \emph{après} le verbe).
\item La phrase se termine souvent par un \textbf{point d'exclamation}.
\end{itemize}

\begin{definition}[L'impératif]
L'\cle{impératif} est le mode verbal qui sert à \textbf{donner un ordre}, à \textbf{donner un conseil}, à \textbf{formuler une consigne} ou à \textbf{faire une invitation}. En allemand, il existe quatre formes d'impératif, selon la personne à qui l'on s'adresse : \all{du} (tu), \all{ihr} (vous, pluriel familier), \all{wir} (nous, invitation) et \all{Sie} (vous, politesse).
\end{definition}

Comparaison avec le français : le français possède aussi l'impératif (\emph{plante ! plantons ! plantez !}), mais l'allemand distingue \textbf{quatre} formes là où le français n'en a que trois, car l'allemand sépare le « vous » familier (\all{ihr}) du « vous » de politesse (\all{Sie}).

\courssub{Formation de l'impératif : les quatre formes}

\textbf{1. La forme \all{du} (2\up{e} personne du singulier)}

\begin{propriete}[Impératif singulier : le radical du présent]
À la 2\up{e} personne du singulier, l'impératif se forme avec le \textbf{radical du présent}, c'est-à-dire la forme \all{du} du présent \textbf{sans la terminaison -t} :
\all{du pflanzt} → \all{Pflanz(e)!} ; \all{du kommst} → \all{Komm!} ; \all{du machst} → \all{Mach!}

Le \all{-e} final est facultatif pour la plupart des verbes : \all{Pflanz!} ou \all{Pflanze!} sont tous deux corrects. Il est toutefois \textbf{obligatoire} quand le radical se termine en \all{-d}, \all{-t}, \all{-ig} ou en groupe de consonnes difficile à prononcer : \all{du wartest} → \all{Warte!} ; \all{du öffnest} → \all{Öffne!} ; \all{du entschuldigst} → \all{Entschuldige!}
\end{propriete}

\begin{attention}[Attention : deux types de radicaux au présent]
Ne confonds pas les deux situations : pour un verbe faible comme \all{pflanzen}, la forme \all{du} du présent est \all{du pflanzt} (terminaison \all{-t}, sans \all{-s}), et l'impératif est le radical pur \all{Pflanz(e)!}. Pour un verbe fort comme \all{kommen}, la forme est \all{du kommst} (terminaison \all{-st}), et l'impératif supprime le \all{-st} : \all{Komm!}. Dans les deux cas, on retombe sur le \textbf{radical} du verbe.
\end{attention}

\begin{propriete}[Les verbes à mutation e→i/ie gardent la voyelle changée]
Les verbes forts qui changent leur voyelle \all{e} en \all{i} ou \all{ie} au présent (\all{du hilfst}, \all{du liest}, \all{du nimmst}, \all{du sprichst}) \textbf{conservent cette voyelle changée} à l'impératif, et ne prennent \textbf{jamais} le \all{-e} final :
\all{du hilfst} → \all{Hilf!} ; \all{du liest} → \all{Lies!} ; \all{du nimmst} → \all{Nimm!} ; \all{du sprichst} → \all{Sprich!}
\end{propriete}

\begin{attention}[Piège : les verbes à mutation a→ä ne changent PAS à l'impératif]
Erreur classique des francophones ! Les verbes qui changent \all{a} en \all{ä} au présent (\all{du fährst}, \all{du schläfst}, \all{du trägst}) \textbf{retrouvent le a simple} à l'impératif :
\all{du fährst} → \all{Fahr!} (et non \all{Fähr!}) ; \all{du schläfst} → \all{Schlaf!} ; \all{du trägst} → \all{Trag!}

Astuce : la mutation \all{a→ä} n'existe qu'avec les terminaisons \all{-st}/\all{-t} ; comme l'impératif supprime cette terminaison, la mutation disparaît aussi.
\end{attention}

\textbf{2. La forme \all{ihr} (2\up{e} personne du pluriel)}

\begin{propriete}[Impératif pluriel familier : la forme du présent sans « ihr »]
À la 2\up{e} personne du pluriel, l'impératif est \textbf{identique à la forme \all{ihr} du présent}, sans le pronom :
\all{ihr pflanzt} → \all{Pflanzt!} ; \all{ihr kommt} → \all{Kommt!} ; \all{ihr helft} → \all{Helft!}

Aucune exception (sauf \all{sein} → \all{seid!}). C'est la forme la plus facile !
\end{propriete}

\textbf{3. La forme \all{Sie} (politesse)}

\begin{propriete}[Impératif de politesse : infinitif + Sie]
La forme de politesse reprend l'\textbf{infinitif} (identique à la forme \all{Sie} du présent) suivi du pronom \all{Sie} placé \textbf{après} le verbe :
\all{Pflanzen Sie!} ; \all{Kommen Sie!} ; \all{Schützen Sie die Umwelt!}

Le pronom \all{Sie} est \textbf{obligatoire} et prend toujours la majuscule.
\end{propriete}

\textbf{4. La forme \all{wir} (invitation : « Allons-y ! »)}

\begin{propriete}[Impératif à la 1\up{re} personne du pluriel : infinitif + wir]
Pour inviter quelqu'un à faire quelque chose avec soi (équivalent du français \emph{allons !}, \emph{plantons !}), on utilise l'infinitif suivi de \all{wir} :
\all{Gehen wir in den Wald!} (Allons dans la forêt !) ; \all{Pflanzen wir einen Baum!} (Plantons un arbre !)
\end{propriete}

\begin{exemplebox}[Exemples traduits]
\all{Pflanze einen Baum!} « Plante un arbre ! » (à un ami)\\
\all{Pflanzt Bäume!} « Plantez des arbres ! » (à un groupe d'élèves)\\
\all{Schützen Sie die Umwelt!} « Protégez l'environnement ! » (discours officiel, affiche)\\
\all{Gehen wir in den Wald!} « Allons dans la forêt ! »\\
\all{Hilf mir bitte!} « Aide-moi, s'il te plaît ! »\\
\all{Lies den Text!} « Lis le texte ! »\\
\all{Fahr langsam!} « Roule doucement ! »
\end{exemplebox}

\begin{center}
\begin{tabularx}{\linewidth}{|l|X|X|X|}\hline
\rowcolor{popblueL} Personne & Règle & pflanzen & helfen (e→i) \\ \hline
\all{du} & radical (+ e facultatif) & \all{Pflanz(e)!} & \all{Hilf!} \\ \hline
\all{ihr} & forme \all{ihr} du présent & \all{Pflanzt!} & \all{Helft!} \\ \hline
\all{wir} & infinitif + \all{wir} & \all{Pflanzen wir!} & \all{Helfen wir!} \\ \hline
\all{Sie} & infinitif + \all{Sie} & \all{Pflanzen Sie!} & \all{Helfen Sie!} \\ \hline
\end{tabularx}
\end{center}

\begin{popfigure}
\begin{tikzpicture}[node distance=0.35cm]
\node[draw=popblue, fill=popblueL, rounded corners, align=center, font=\bfseries] (titre) {Der Imperativ : pflanzen};
\node[draw=popgreen, fill=popgreenL, rounded corners, align=center, below left=0.7cm and 1.6cm of titre] (du) {du\\ \all{Pflanz(e)!}\\ \footnotesize radical};
\node[draw=poporange, fill=poporangeL, rounded corners, align=center, below=0.7cm of titre, xshift=-0.4cm] (ihr) {ihr\\ \all{Pflanzt!}\\ \footnotesize présent sans \all{ihr}};
\node[draw=poppurple, fill=poppurpleL, rounded corners, align=center, below=0.7cm of titre, xshift=0.4cm] (wir) {wir\\ \all{Pflanzen wir!}\\ \footnotesize infinitif + \all{wir}};
\node[draw=poppink, fill=poppinkL, rounded corners, align=center, below right=0.7cm and 1.6cm of titre] (sie) {Sie\\ \all{Pflanzen Sie!}\\ \footnotesize infinitif + \all{Sie}};
\draw[-{Stealth}, thick, popblue] (titre) -- (du);
\draw[-{Stealth}, thick, popblue] (titre) -- (ihr);
\draw[-{Stealth}, thick, popblue] (titre) -- (wir);
\draw[-{Stealth}, thick, popblue] (titre) -- (sie);
\end{tikzpicture}
\legende{Les quatre formes de l'impératif du verbe \all{pflanzen} (planter).}
\end{popfigure}

\courssub{Cas particuliers}

\textbf{1. Le verbe \all{sein} (être)} est totalement irrégulier :

\begin{center}
\begin{tabular}{|l|l|l|}\hline
\rowcolor{popblueL} Personne & Impératif & Exemple traduit \\ \hline
\all{du} & \all{Sei!} & \all{Sei ruhig!} « Sois tranquille ! » \\ \hline
\all{ihr} & \all{Seid!} & \all{Seid vorsichtig!} « Soyez prudents ! » \\ \hline
\all{wir} & \all{Seien wir!} & \all{Seien wir freundlich!} « Soyons aimables ! » \\ \hline
\all{Sie} & \all{Seien Sie!} & \all{Seien Sie bitte pünktlich!} « Soyez ponctuel, s.v.p. ! » \\ \hline
\end{tabular}
\end{center}

\textbf{2. Les verbes en \all{-eln} et \all{-ern}} : à l'impératif \all{du}, on garde le \all{-e} du radical : \all{sammeln} → \all{Sammle!} (ramasse !) ; \all{wandern} → \all{Wandere!} (marche, randonne !). Le \all{-e} est ici indispensable à la prononciation.

\textbf{3. Les verbes à particule séparable} : la particule se détache et va \textbf{à la fin}, comme au présent : \all{mitkommen} → \all{Komm mit!} (viens avec nous !) ; \all{aufstehen} → \all{Steht auf!} (levez-vous !) ; \all{zuhören} → \all{Hören Sie zu!} (écoutez !).

\courssub{Ponctuation et atténuation de l'ordre}

\begin{propriete}[Point d'exclamation et adverbes d'atténuation]
L'impératif allemand s'écrit très souvent avec un \textbf{point d'exclamation}. Pour adoucir un ordre et le transformer en conseil ou en invitation amicale, on ajoute les petits mots \cle{bitte} (s'il te plaît), \cle{mal} (un peu, pour voir) ou \cle{doch} (allez, insistance amicale) :
\all{Komm mal her!} « Viens donc voir ! » ; \all{Hilf mir doch!} « Aide-moi donc ! » ; \all{Pflanzen Sie bitte hier!} « Plantez ici, je vous prie ! »
\end{propriete}

\begin{attention}[Erreurs fréquentes des francophones]
\begin{itemize}
\item Ne mets \textbf{jamais le pronom} devant le verbe : on dit \all{Komm!} et non \all{Du komm!}. Seuls \all{Sie} et \all{wir} apparaissent, \emph{après} le verbe.
\item N'oublie pas la mutation \all{e→i/ie} : \all{Hilf!}, \all{Lies!}, et non \all{Helf!}, \all{Les!}.
\item N'ajoute pas de tréma à l'impératif : \all{Fahr!} et non \all{Fähr!}.
\item Attention au faux ami de construction : « Allons-y ! » se dit \all{Gehen wir!} ; le verbe précède toujours le pronom.
\item À l'impératif négatif, la négation \all{nicht} se place normalement : \all{Fahr nicht so schnell!} « Ne roule pas si vite ! » ; \all{Vergiss das nicht!} « N'oublie pas cela ! »
\end{itemize}
\end{attention}

\begin{methode}[Vérifier une forme d'impératif en trois étapes]
\begin{enumerate}
\item \textbf{Identifier la personne} à qui l'on s'adresse : un ami (\all{du}), plusieurs camarades (\all{ihr}), un adulte ou un inconnu (\all{Sie}), une invitation commune (\all{wir}).
\item \textbf{Reconstruire la phrase au présent} avec le pronom : \all{du pflanzt}, \all{ihr pflanzt}, \all{Sie pflanzen}.
\item \textbf{Supprimer le pronom} (sauf \all{Sie}/\all{wir}) et la terminaison pour \all{du} ; vérifier la mutation éventuelle (\all{helfen} → \all{Hilf!} ; \all{fahren} → \all{Fahr!}).
\end{enumerate}
\end{methode}

\courssub{Application au thème : la campagne de reboisement}

\begin{textebox}[Un appel du club écologique du lycée]
Liebe Schülerinnen und Schüler!

Unser Wald ist in Gefahr. Jeden Tag verschwinden Bäume, und mit ihnen die Tiere. Aber wir können handeln! Kommt am Samstag in den Schulgarten und pflanzt mit uns junge Bäume! Nehmt einen Spaten, nehmt einen Freund, und kommt um acht Uhr! Bringt Wasser und gute Laune mit!

Seid vorsichtig mit den jungen Pflanzen: gießt sie jeden Tag und schützt sie vor den Ziegen. Sammelt auch den Müll im Garten — eine saubere Schule ist eine schöne Schule!

Und Sie, liebe Eltern und Lehrer: Unterstützen Sie unsere Aktion! Geben Sie uns Setzlinge und Werkzeuge! Besuchen Sie unseren Naturpark am Wochenende und sprechen Sie mit Ihren Nachbarn über den Ökotourismus in Kamerun!

Gehen wir gemeinsam diesen Weg! Jeder Baum zählt. Pflanzen wir heute, damit unsere Kinder morgen im Schatten sitzen können!

Euer Öko-Club
\end{textebox}

\textbf{Glossaire} : \all{die Gefahr, -en} le danger ; \all{verschwinden} disparaître ; \all{handeln} agir ; \all{der Schulgarten, -gärten} le jardin de l'école ; \all{der Spaten, -} la bêche ; \all{die Laune, -n} l'humeur ; \all{gießen} arroser ; \all{die Ziege, -n} la chèvre ; \all{der Müll} les déchets ; \all{unterstützen} soutenir ; \all{der Setzling, -e} le jeune plant ; \all{das Werkzeug, -e} l'outil ; \all{der Schatten, -} l'ombre.

\textbf{Questions de compréhension} : 1) Que propose le club écologique pour samedi ? 2) Relevez trois impératifs à la forme \all{ihr} et deux à la forme \all{Sie}. 3) Pourquoi faut-il arroser les jeunes plants ? 4) Quelle invitation à la forme \all{wir} trouve-t-on à la fin du texte ?

\begin{exoresolu}[Mets les verbes à l'impératif]
Énoncé : mets les verbes entre parenthèses à l'impératif demandé. a) (helfen, du) \all{... mir bitte im Garten!} b) (mitkommen, ihr) \all{... in den Naturpark!} c) (schützen, Sie) \all{... die Umwelt!} d) (fahren, du) \all{... nicht so schnell!} e) (sein, ihr) \all{... bitte pünktlich!} f) (gehen, wir) \all{... jetzt in den Wald!}
\tcblower
Solution détaillée : a) \all{Hilf mir bitte im Garten!} — verbe fort à mutation \all{e→i} : \all{du hilfst} → \all{Hilf!}, sans \all{-e} final. b) \all{Kommt mit in den Naturpark!} — forme \all{ihr} du présent ; la particule séparable \all{mit} se détache et va à la fin (ou juste après le verbe devant un complément long). c) \all{Schützen Sie die Umwelt!} — infinitif + \all{Sie}. d) \all{Fahr nicht so schnell!} — attention : mutation \all{a→ä} supprimée à l'impératif (et non \all{Fähr!}). e) \all{Seid bitte pünktlich!} — forme irrégulière de \all{sein}. f) \all{Gehen wir jetzt in den Wald!} — infinitif + \all{wir} : « Allons maintenant dans la forêt ! »
\end{exoresolu}

\begin{savaistu}[L'infinitif, la consigne impersonnelle des panneaux allemands]
Dans les pays germanophones, les panneaux officiels n'utilisent souvent pas l'impératif mais l'\textbf{infinitif}, plus neutre et impersonnel : \all{Nicht rauchen!} « Défense de fumer », \all{Bitte nicht auf den Rasen treten!} « Prière de ne pas marcher sur la pelouse », \all{Keinen Müll wegwerfen!} « Ne jetez pas de déchets ». Dans les parcs naturels allemands, autrichiens et suisses, tu liras partout ces consignes à l'infinitif. Retiens cette forme : elle est très appréciée au Bac dans les sujets de compréhension écrite !
\end{savaistu}

\begin{exercice}[À toi de jouer]
1) Écris une affiche pour le club écologique de ton lycée avec quatre impératifs : deux à la forme \all{ihr}, un à la forme \all{du} et un à la forme \all{Sie}. 2) Transforme : \all{du nimmst} → ? ; \all{ihr lest} → ? ; \all{Sie kommen} → ? ; \all{du fährst} → ?
\end{exercice}

\begin{corrige}[Exercice : À toi de jouer]
1) Exemple de réponse : \all{Kommt am Samstag! Pflanzt Bäume! Hilf uns, die Schule sauber zu halten! Unterstützen Sie unseren Öko-Club!} 2) \all{Nimm!} ; \all{Lest!} ; \all{Kommen Sie!} ; \all{Fahr!}
\end{corrige}

\begin{aretenir}[L'essentiel de l'impératif]
Quatre formes : \all{du} = radical du présent (\all{Pflanz!}, avec mutation \all{e→i/ie} conservée : \all{Hilf!}, \all{Lies!}) ; \all{ihr} = forme du présent sans pronom (\all{Pflanzt!}) ; \all{wir} = infinitif + \all{wir} (\all{Gehen wir!}) ; \all{Sie} = infinitif + \all{Sie} (\all{Schützen Sie!}). Le verbe est en tête de phrase, souvent suivi d'un point d'exclamation ; \all{bitte}, \all{mal}, \all{doch} adoucissent l'ordre. Irrégulier : \all{sein} → \all{sei!}, \all{seid!}, \all{seien Sie!}. Jamais de tréma à l'impératif : \all{Fahr!}
\end{aretenir}

\coursec{Grammaire 2 : les prépositions avec l'accusatif}

Dans la section précédente, nous avons appris à former l'impératif pour donner des consignes : \all{Pflanzt Bäume! Schützt die Umwelt!} Mais pour exprimer \emph{où} l'on passe, \emph{pour quoi} l'on agit, \emph{contre quoi} l'on lutte, il faut des prépositions. Le texte sur le reboisement en contient plusieurs qui obéissent toutes à la même loi : elles exigent \cle{l'accusatif}. C'est l'une des règles les plus sûres et les plus rentables de la grammaire allemande : apprends-la une fois, et tu ne te tromperas plus jamais.

\courssub{Observation : partons du texte}

\begin{textebox}[Extraits du texte « Aufforstung und Ökotourismus in Kamerun »]
„Die Jugendlichen pflanzen Bäume \textbf{für die Zukunft} ihrer Kinder.“\\
„Die Touristen wandern \textbf{durch den Wald} und beobachten die Vögel.“\\
„Im Naturpark genießen die Besucher die Natur \textbf{ohne Lärm} und \textbf{ohne Müll}.“\\
„Viele Gruppen kämpfen \textbf{gegen die Abholzung} des Regenwaldes.“
\end{textebox}

Observons les groupes en gras. Après \all{für}, \all{durch}, \all{ohne}, \all{gegen}, l'article du masculin \all{der} est devenu \all{den} (\all{durch \textbf{den} Wald}), tandis que \all{die} reste \all{die} (\all{für \textbf{die} Zukunft}). Ce n'est pas un hasard : ces prépositions gouvernent toujours le même cas, \cle{l'accusatif}.

\courssub{La liste des prépositions + accusatif}

\begin{definition}[Les sept prépositions de l'accusatif]
Sept prépositions sont \textbf{toujours} suivies de l'accusatif, sans exception :
\begin{center}
\begin{tabularx}{\linewidth}{|l|X|X|}
\hline
\rowcolor{popblueL} \textbf{Préposition} & \textbf{Sens principal} & \textbf{Exemple} \\
\hline
\all{durch} & à travers, par & \all{durch den Wald} (à travers la forêt) \\
\hline
\all{für} & pour & \all{für die Umwelt} (pour l'environnement) \\
\hline
\all{gegen} & contre ; vers (heure) & \all{gegen die Abholzung} (contre la déforestation) \\
\hline
\all{ohne} & sans & \all{ohne Geld} (sans argent) \\
\hline
\all{um} & autour de ; à (heure) & \all{um den See} (autour du lac) ; \all{um 8 Uhr} (à 8 heures) \\
\hline
\all{bis} & jusqu'à & \all{bis nächste Woche} (jusqu'à la semaine prochaine) \\
\hline
\all{entlang} & le long de & \all{den Fluss entlang} (le long de la rivière) \\
\hline
\end{tabularx}
\end{center}
\end{definition}

\begin{methode}[Le moyen mnémotechnique DOGFU]
Pour mémoriser les cinq prépositions principales, retiens le mot \cle{DOGFU} :
\begin{enumerate}
\item \textbf{D} — \all{durch}
\item \textbf{O} — \all{ohne}
\item \textbf{G} — \all{gegen}
\item \textbf{F} — \all{für}
\item \textbf{U} — \all{um}
\end{enumerate}
Ajoute ensuite \all{bis} et \all{entlang}. Chaque fois que tu rencontres l'une de ces sept prépositions, applique automatiquement l'accusatif, puis vérifie l'article du masculin : c'est lui qui change (\all{der} $\rightarrow$ \all{den}).
\end{methode}

\begin{popfigure}
\begin{tikzpicture}
\node[circle, draw=popblue, fill=popblueL, thick, minimum size=2.6cm, align=center, font=\bfseries] (c) {Akkusativ};
\node[draw=popgreen, fill=popgreenL, rounded corners, align=center] (d) at (90:3.4cm) {\all{durch}\\à travers};
\node[draw=popgreen, fill=popgreenL, rounded corners, align=center] (o) at (38.6:3.4cm) {\all{ohne}\\sans};
\node[draw=popgreen, fill=popgreenL, rounded corners, align=center] (g) at (-12.8:3.4cm) {\all{gegen}\\contre};
\node[draw=popgreen, fill=popgreenL, rounded corners, align=center] (f) at (-64.3:3.4cm) {\all{für}\\pour};
\node[draw=popgreen, fill=popgreenL, rounded corners, align=center] (u) at (-115.7:3.4cm) {\all{um}\\autour de};
\node[draw=poporange, fill=poporangeL, rounded corners, align=center] (b) at (-167.1:3.4cm) {\all{bis}\\jusqu'à};
\node[draw=poporange, fill=poporangeL, rounded corners, align=center] (e) at (141.4:3.4cm) {\all{entlang}\\le long de};
\draw[-{Stealth}, thick, popblue] (c) -- (d);
\draw[-{Stealth}, thick, popblue] (c) -- (o);
\draw[-{Stealth}, thick, popblue] (c) -- (g);
\draw[-{Stealth}, thick, popblue] (c) -- (f);
\draw[-{Stealth}, thick, popblue] (c) -- (u);
\draw[-{Stealth}, thick, poporange] (c) -- (b);
\draw[-{Stealth}, thick, poporange] (c) -- (e);
\end{tikzpicture}
\legende{Les sept prépositions qui exigent toujours l'accusatif (en vert : DOGFU ; en orange : bis et entlang).}
\end{popfigure}

\courssub{Rappel : les formes de l'accusatif}

Bonne nouvelle : à l'accusatif, \textbf{seul le masculin change}. Le féminin, le neutre et le pluriel restent identiques au nominatif.

\begin{center}
\begin{tabular}{|l|l|l|l|}
\hline
\rowcolor{popblueL} \textbf{Genre} & \textbf{Nominatif} & \textbf{Akkusatif} & \textbf{Exemple} \\
\hline
masculin & \all{der Wald} & \all{den Wald} & \all{durch den Wald} \\
\hline
féminin & \all{die Umwelt} & \all{die Umwelt} & \all{für die Umwelt} \\
\hline
neutre & \all{das Geld} & \all{das Geld} & \all{ohne das Geld} \\
\hline
pluriel & \all{die Bäume} & \all{die Bäume} & \all{für die Bäume} \\
\hline
\end{tabular}
\end{center}

Même logique pour l'article indéfini : \all{ein Baum} $\rightarrow$ \all{einen Baum} ; \all{eine Schule} $\rightarrow$ \all{eine Schule} ; \all{ein Haus} $\rightarrow$ \all{ein Haus}. Les adjectifs possessifs suivent le même modèle : \all{mein} $\rightarrow$ \all{meinen} (masculin), tandis que \all{meine} (féminin et pluriel) et \all{mein} (neutre) restent inchangés : \all{für meinen Bruder}, \all{für meine Schwester}, \all{für mein Dorf}.

\begin{attention}[Seul le masculin change !]
L'erreur classique du francophone consiste à « accorder » tous les articles comme en français. Non : \all{die} et \all{das} ne changent \textbf{jamais} à l'accusatif. On écrit \all{für die Umwelt} et jamais \all{für dien Umwelt} ; \all{ohne das Geld} et jamais \all{ohne den Geld}. Concentre toute ton attention sur le seul cas qui bouge : \all{der} $\rightarrow$ \all{den}, \all{ein} $\rightarrow$ \all{einen}.
\end{attention}

\courssub{Sens et emplois en détail}

\begin{exemplebox}[Exemples traduits]
\begin{itemize}
\item \all{Wir wandern durch den Naturpark.} « Nous randonnons à travers le parc naturel. »
\item \all{Sie arbeiten für die Umwelt.} « Ils travaillent pour l'environnement. »
\item \all{Er kämpft gegen die Abholzung.} « Il lutte contre la déforestation. »
\item \all{Ohne Bäume kein Leben.} « Sans arbres, pas de vie. »
\item \all{Der Weg führt um den See.} « Le chemin fait le tour du lac. »
\item \all{Der Kurs beginnt um neun Uhr.} « Le cours commence à neuf heures. »
\item \all{Wir pflanzen Bäume bis nächsten Montag.} « Nous plantons des arbres jusqu'à lundi prochain. »
\item \all{Wir gehen den Fluss entlang.} « Nous marchons le long de la rivière. »
\end{itemize}
\end{exemplebox}

Deux emplois temporels à retenir : \all{um} + heure exacte (\all{um 14 Uhr} « à 14 heures ») et \all{gegen} + heure approximative (\all{gegen 14 Uhr} « vers 14 heures »). \all{bis} s'utilise souvent sans article : \all{bis Montag}, \all{bis morgen}.

\begin{attention}[« entlang » se place souvent APRÈS le nom]
Contrairement aux autres prépositions, \all{entlang} se met le plus souvent \emph{après} le nom à l'accusatif : \all{Wir gehen den Fluss entlang.} (« Nous marchons le long de la rivière. ») ; \all{Sie fahren die Straße entlang.} (« Ils roulent le long de la route. »). Ne dis jamais \all{entlang den Fluss} dans la langue courante : place \all{entlang} derrière.
\end{attention}

\courssub{Les contractions usuelles}

À l'oral comme à l'écrit, certaines prépositions fusionnent avec l'article \all{das} :

\begin{center}
\begin{tabularx}{\linewidth}{|X|X|X|}
\hline
\rowcolor{popblueL} \textbf{Contraction} & \textbf{Forme complète} & \textbf{Exemple} \\
\hline
\all{fürs} & \all{für das} & \all{ein Geschenk fürs Leben} (un cadeau pour la vie) \\
\hline
\all{durchs} & \all{durch das} & \all{durchs Dorf} (à travers le village) \\
\hline
\all{ums} & \all{um das} & \all{ums Haus} (autour de la maison) \\
\hline
\end{tabularx}
\end{center}

Attention : on ne contracte qu'avec le neutre \all{das}. Avec le masculin, on garde les deux mots : \all{für den Wald}, jamais de contraction.

\courssub{Comparaison français / allemand}

\begin{center}
\begin{tabularx}{\linewidth}{|X|X|X|}
\hline
\rowcolor{popblueL} \textbf{Français} & \textbf{Deutsch} & \textbf{Remarque} \\
\hline
sans un arbre & \all{ohne einen Baum} & masculin : \all{ein} $\rightarrow$ \all{einen} \\
\hline
pour l'école & \all{für die Schule} & féminin : inchangé \\
\hline
à travers le parc & \all{durch den Park} & masculin : \all{der} $\rightarrow$ \all{den} \\
\hline
le long de la route & \all{die Straße entlang} & préposition postposée \\
\hline
jusqu'à demain & \all{bis morgen} & souvent sans article \\
\hline
\end{tabularx}
\end{center}

\begin{attention}[Ne confonds pas avec les prépositions à deux cas]
Les prépositions \all{in}, \all{an}, \all{auf}, \all{über}, \all{unter} (que tu étudieras plus tard) hésitent entre accusatif et datif selon le sens (mouvement ou position). Ici, aucune hésitation : après \all{durch}, \all{für}, \all{gegen}, \all{ohne}, \all{um}, \all{bis}, \all{entlang}, c'est \textbf{toujours} l'accusatif, dans tous les contextes. C'est une règle sans exception.
\end{attention}

\courssub{Mise en pratique}

\begin{exoresolu}[Complète avec la bonne préposition et le bon article]
Complète : 1. Wir wandern \all{durch} \dots\dots\ Wald. (der) — 2. Die Schüler sammeln Müll \all{für} \dots\dots\ Umwelt. (die) — 3. \all{Ohne} \dots\dots\ Baum (ein) gibt es keinen Schatten. — 4. Wir gehen \dots\dots\ Fluss (der) \all{entlang}. — 5. Der Bus kommt \all{um} \dots\dots\ Uhr. (8)
\tcblower
\textbf{Solution détaillée :}
\begin{enumerate}
\item \all{durch \textbf{den} Wald} — masculin, accusatif : \all{der} $\rightarrow$ \all{den}.
\item \all{für \textbf{die} Umwelt} — féminin, inchangé.
\item \all{Ohne \textbf{einen} Baum} — masculin avec article indéfini : \all{ein} $\rightarrow$ \all{einen}.
\item \all{\textbf{den} Fluss entlang} — masculin à l'accusatif, et \all{entlang} se place après le nom.
\item \all{um \textbf{8} Uhr} — heure exacte, sans article.
\end{enumerate}
\end{exoresolu}

\begin{exercice}[À toi de jouer]
Complète avec la préposition DOGFU qui convient et mets l'article à l'accusatif :
\begin{enumerate}
\item Die Touristen fahren \dots\dots\ \dots\dots\ Regenwald. (der ; à travers)
\item Viele Menschen kämpfen \dots\dots\ \dots\dots\ Abholzung. (die ; contre)
\item Er reist \dots\dots\ \dots\dots\ Pass. (der ; sans)
\item Wir sammeln Geld \dots\dots\ \dots\dots\ Naturpark. (der ; pour)
\item Sie laufen \dots\dots\ See (der) entlang.
\end{enumerate}
\end{exercice}

\begin{corrige}[Exercice « À toi de jouer »]
1. \all{durch den Regenwald} — 2. \all{gegen die Abholzung} — 3. \all{ohne den Pass} — 4. \all{für den Naturpark} — 5. \all{den See entlang}. Vérifie à chaque fois : la préposition appartient-elle à la liste DOGFU + \all{bis}/\all{entlang} ? Si oui, accusatif automatique, et seul le masculin change.
\end{corrige}

\begin{aretenir}[L'essentiel]
Après \all{durch}, \all{ohne}, \all{gegen}, \all{für}, \all{um} (DOGFU), ainsi que \all{bis} et \all{entlang}, l'allemand exige \textbf{toujours} l'accusatif. À l'accusatif, seul le masculin change : \all{der} $\rightarrow$ \all{den}, \all{ein} $\rightarrow$ \all{einen}. La préposition \all{entlang} se place généralement après le nom : \all{den Fluss entlang}. Contractions fréquentes avec le neutre : \all{fürs}, \all{durchs}, \all{ums}.
\end{aretenir}

\coursec{Méthode du commentaire et production guidée}

Après avoir maîtrisé les prépositions à régime accusatif (section \textbf{Grammaire 2}), nous disposons maintenant de tous les outils linguistiques — vocabulaire thématique, impératif, prépositions \cle{durch, für, gegen, ohne, um} — pour analyser le texte support et produire nos propres énoncés. Cette section vous guide pas à pas vers l'autonomie en compréhension et en expression, compétences centrales de l'examen de classe et du Baccalauréat.

\courssub{La méthode du commentaire de texte en 4 étapes}

\begin{methode}[Comment réussir un commentaire de texte en allemand]
Un bon commentaire suit une structure logique en quatre temps. Respectez cet ordre pour organiser vos idées et obtenir la note maximale.
\begin{enumerate}
    \item \textbf{Présenter} (\all{einleiten}) : donnez le type de document (article, rapport, interview), le thème principal (\all{Aufforstung, Ökotourismus}), la source supposée (\all{Zeitung, Schulzeitung, Internet}) et l'intention de l'auteur (\all{informieren, warnen, überzeugen}).
    \item \textbf{Résumer} (\all{zusammenfassen}) : reprenez les idées principales \emph{dans l'ordre du texte} sans copier-coller. Utilisez vos propres mots (paraphrase). Limitez-vous à l'essentiel (3 à 4 phrases maximum).
    \item \textbf{Analyser} (\all{analysieren}) : montrez \emph{comment} le texte est construit. Relevez le champ lexical (\all{Wortfelder} : \all{der Wald, der Schutz}), les structures grammaticales (impératif, prépositions + accusatif), les connecteurs logiques (\all{zuerst, deshalb}). Expliquez l'effet produit sur le lecteur.
    \item \textbf{Donner son avis argumenté} (\all{Stellung nehmen}) : exprimez votre opinion personnelle \emph{en vous appuyant sur le texte}. Utilisez les expressions de l'opinion (\all{Meiner Meinung nach...}). Liez votre avis aux arguments du texte (\all{deshalb, weil}).
\end{enumerate}
\cle{Règle d'or} : citez toujours un exemple précis du texte (numéro de ligne ou citation courte entre guillemets „…“) pour justifier votre analyse ou votre avis.
\end{methode}

\begin{popfigure}
\begin{tikzpicture}[
    node distance=1.2cm and 1.5cm,
    every node/.style={align=center, font=\small},
    box/.style={draw=popdark, thick, rounded corners=6pt, fill=popcream, minimum width=3.2cm, minimum height=1.8cm},
    arrow/.style={-Stealth, thick, popblue}
]
\node[box, fill=popblueL, align=center] (p1){\textbf{1. Présenter}\\\textit{(einleiten)}\\Type, thème,\\source, intention};
\node[box, fill=popgreenL, right=of p1, align=center] (p2){\textbf{2. Résumer}\\\textit{(zusammenfassen)}\\Idées principales,\\ordre, propres mots};
\node[box, fill=poporangeL, right=of p2, align=center] (p3){\textbf{3. Analyser}\\\textit{(analysieren)}\\Lexique, grammaire,\\connecteurs, effet};
\node[box, fill=poppinkL, right=of p3, align=center] (p4){\textbf{4. Avis personnel}\\\textit{(Stellung nehmen)}\\Opinion, arguments\\du texte, connecteurs};

\draw[arrow] (p1.east) -- (p2.west);
\draw[arrow] (p2.east) -- (p3.west);
\draw[arrow] (p3.east) -- (p4.west);
\end{tikzpicture}
\legende{Organigramme : les 4 étapes du commentaire de texte structuré}
\end{popfigure}

\courssub{Commentaire modèle sur le texte support}

Le texte support de la section 1 („Aufforstung und Ökotourismus in Kamerun“) est un article d'information d'environ 180 mots. Voici un commentaire modèle respectant les 4 étapes et la longueur attendue (8 à 10 lignes).

\begin{exemplebox}[Commentaire modèle : „Aufforstung und Ökotourismus in Kamerun“]
\textbf{Présentation :} Il s'agit d'un article de presse en ligne traitant du reboisement et de l'écotourisme au Cameroun. L'auteur vise à informer le grand public sur les initiatives locales et à convaincre de l'utilité de la protection des forêts.

\textbf{Résumé :} D'abord, le texte décrit la dégradation des forêts camerounaises due à l'exploitation illégale. Ensuite, il présente le projet „Green Cameroon“, qui plante des arbres indigènes avec les villageois. De plus, il explique que l'écotourisme crée des emplois durables pour les jeunes. Enfin, il appelle à la responsabilité de chacun.

\textbf{Analyse :} Le champ lexical de la nature (\all{der Wald, der Baum, die Artenvielfalt}) domine. L'auteur utilise l'impératif (\all{Pflanzen Sie Bäume!}) pour mobiliser le lecteur. Les prépositions accusatives (\all{für die Umwelt}, \all{durch den Tourismus}) précisent les buts et les causes. Le connecteur \cle{deshalb} marque la conséquence logique.

\textbf{Avis :} \all{Meiner Meinung nach ist das Projekt vorbildlich, weil es Umweltschutz mit wirtschaftlicher Entwicklung verbindet. Deshalb sollten wir alle solche Initiativen unterstützen.} « À mon avis, le projet est exemplaire parce qu'il associe la protection de l'environnement au développement économique. C'est pourquoi nous devrions tous soutenir de telles initiatives. »
\end{exemplebox}

\begin{attention}[Connecteurs et variété syntaxique]
Ne commencez \textbf{pas} toutes vos phrases par \all{Ich} ou \all{Der Text}. Variez les débuts de phrase avec des connecteurs (topicalisation) : le verbe conjugué reste en \cle{deuxième position}.
\begin{center}
\begin{tabularx}{\linewidth}{|X|X|X|}\hline
\rowcolor{popblueL} \textbf{Au lieu de...} & \textbf{Utilisez...} & \textbf{Exemple} \\ \hline
\all{Ich denke, das ist gut.} & \cle{Meiner Meinung nach} + V2 & \all{Meiner Meinung nach ist der Wald wichtig.} \\ \hline
\all{Der Text sagt: „...“} & \cle{Zuerst / Dann / Außerdem} + V2 & \all{Zuerst beschreibt der Text die Probleme.} \\ \hline
\all{Es ist wichtig, weil...} & \cle{Deshalb / Daher} + V2 & \all{Der Wald schützt das Klima. Deshalb ist er wichtig.} \\ \hline
\end{tabularx}
\end{center}
\end{attention}

\courssub{Expressions pour donner son avis (Stellungnahme)}

\begin{definition}[Formules types pour l'argumentation (niveau A2)]
Ces structures permettent d'exprimer une opinion nuancée. Attention : après \cle{dass} (que), le verbe conjugué va \textbf{à la fin} de la proposition subordonnée.
\end{definition}

\begin{center}
\begin{tabularx}{\linewidth}{|X|X|X|}\hline
\rowcolor{popblueL} \textbf{Structure allemande} & \textbf{Traduction française} & \textbf{Exemple contextualisé} \\ \hline
\all{Meiner Meinung nach} + V2 & À mon avis / Selon moi & \all{Meiner Meinung nach ist Ökotourismus die Zukunft.} \\ \hline
\all{Ich denke, dass}... (verbe à la fin) & Je pense que... & \all{Ich denke, dass wir mehr Bäume pflanzen müssen.} \\ \hline
\all{Ich finde es wichtig, dass}... & Je trouve important que... & \all{Ich finde es wichtig, dass die Jugend mitmacht.} \\ \hline
\all{Meines Erachtens} + V2 & À mon sens / D'après moi & \all{Meines Erachtens schützt der Wald das Wasser.} \\ \hline
\all{Ich bin der Meinung, dass}... & Je suis d'avis que... & \all{Ich bin der Meinung, dass jeder helfen kann.} \\ \hline
\all{Dafür bin ich. / Dagegen bin ich.} & Je suis pour. / Je suis contre. & \all{Schutz der Natur? Dafür bin ich!} \\ \hline
\end{tabularx}
\end{center}

\begin{attention}[Piège majeur : l'ordre des mots après « dass »]
En français : « Je pense \textbf{que la forêt est importante}. » (sujet + verbe).
En allemand : \all{Ich denke, dass der Wald \textbf{wichtig ist}.} (sujet ... verbe conjugué \textbf{à la fin}).

\textbf{Faux :} \all{Ich denke, dass der Wald ist wichtig.}

\textbf{Correct :} \all{Ich denke, dass der Wald wichtig \textbf{ist}.}

\emph{Règle :} \cle{dass} (comme \cle{weil}, \cle{wenn}) renvoie le verbe conjugué en position finale dans la subordonnée.
\end{attention}

\courssub{Production écrite guidée : „Was können wir für die Umwelt tun?“}

\begin{methode}[Rédiger un paragraphe engagé (6 à 8 phrases)]
Suivez ce plan pour intégrer naturellement les points de grammaire vus dans la leçon.
\begin{enumerate}
    \item \textbf{Phrase d'accroche} (\all{Themensatz}) : constat général + \cle{Meiner Meinung nach}.
    \item \textbf{3 impératifs (appel à l'action) :} utilisez l'impératif formel (forme \cle{Sie}) ou familier (forme \cle{ihr}) pour donner des conseils concrets. Variez les verbes (\all{pflanzen, schützen, vermeiden, sammeln, nutzen}).
    \item \textbf{2 prépositions + accusatif (justification) :} expliquez le but (\cle{für}), la cause ou le moyen (\cle{durch}) ou l'opposition (\cle{gegen}).
    \item \textbf{Phrase de conclusion} (\all{Schluss}) : appel collectif avec \cle{wir} + modal (\all{sollen / müssen / wollen}) + connecteur \cle{deshalb} ou \cle{zusammen}.
\end{enumerate}
\end{methode}

\begin{exoresolu}[Production écrite modèle]
\textbf{Consigne :} rédigez un paragraphe de 7 phrases sur „Was können wir für die Umwelt tun?“. Incluez \textbf{3 impératifs} et \textbf{2 prépositions + accusatif}. Soulignez-les.

\textbf{Proposition corrigée :}

\all{Meiner Meinung nach ist der Schutz unserer Umwelt in Kamerun sehr wichtig.} (1. avis + V2)

\all{\textbf{Pflanzen Sie} jeden Monat einen Baum!} (2. \textbf{impératif 1})

\all{\textbf{Schützen Sie} die Wälder \textbf{gegen} illegale Abholzung!} (3. \textbf{impératif 2} + \textbf{préposition 1 : gegen + accusatif})

\all{\textbf{Vermeiden Sie} Plastikmüll im Alltag!} (4. \textbf{impératif 3})

\all{Wir können Fahrräder \textbf{für} kurze Wege \textbf{nutzen}.} (5. \textbf{préposition 2 : für + accusatif})

\all{Die Schulen sollten Umweltklubs gründen.} (6. proposition concrète, modal \cle{sollten})

\all{\textbf{Deshalb} handeln wir jetzt zusammen für eine grüne Zukunft!} (7. conclusion + \cle{deshalb} + V2)
\tcblower
\textbf{Analyse grammaticale :}
\begin{itemize}
    \item \textbf{3 impératifs (forme Sie) :} \all{Pflanzen Sie}, \all{Schützen Sie}, \all{Vermeiden Sie} (verbe en position 1, \cle{Sie} en position 2).
    \item \textbf{2 prépositions + accusatif :} \all{gegen illegale Abholzung} (accusatif de \cle{die Abholzung}), \all{für kurze Wege} (accusatif pluriel de \cle{der Weg}).
    \item \textbf{Connecteurs :} \cle{Meiner Meinung nach} (début), \cle{deshalb} (conséquence).
    \item \textbf{Position V2 respectée :} \all{Wir können ... nutzen}, \all{Deshalb handeln wir ...}.
\end{itemize}
\end{exoresolu}

\begin{exercice}[À vous de jouer : production écrite]
Rédigez votre propre paragraphe (6 à 8 phrases) sur le thème : \all{„Unser Schulhof soll grüner werden.“} (Notre cour d'école doit devenir plus verte).

\textbf{Contraintes obligatoires :}
\begin{enumerate}
    \item Utilisez \cle{Meiner Meinung nach} ou \cle{Ich finde es wichtig, dass} en phrase 1.
    \item Insérez \textbf{exactement 3 impératifs} (forme \cle{Sie} ou \cle{ihr}) aux phrases 2, 3 et 4.
    \item Insérez \textbf{exactement 2 prépositions accusatives} (\cle{für, gegen, durch, ohne, um}) dans le texte.
    \item Terminez par une phrase avec \cle{deshalb} ou \cle{zusammen}.
\end{enumerate}
\textit{Rédigez d'abord un brouillon, puis recopiez au propre en vérifiant la place du verbe et les majuscules des noms.}
\vspace{4cm}
\end{exercice}

\begin{corrige}[Production écrite : éléments de correction attendus]
Un paragraphe réussi pourra ressembler à ceci (modèle, d'autres formulations sont possibles) :

\all{Meiner Meinung nach soll unser Schulhof grüner und schöner werden. \textbf{Pflanzt} Blumen und Bäume im Schulhof! \textbf{Sammelt} den Müll \textbf{um} das Schulgebäude! \textbf{Schützt} die jungen Setzlinge \textbf{gegen} die Tiere! Wir können auch einen Schulgarten \textbf{für} frisches Gemüse anlegen. Die Lehrer sollen uns dabei helfen. \textbf{Deshalb} arbeiten wir jeden Freitag zusammen im Garten!}

Points de vérification : 3 impératifs corrects (\all{Pflanzt}, \all{Sammelt}, \all{Schützt} — forme \cle{ihr} sans \textbf{-e} final facultatif), 2 à 3 prépositions + accusatif (\all{um das Schulgebäude}, \all{gegen die Tiere}, \all{für frisches Gemüse}), verbe en position 2 après \all{Deshalb}, majuscules à tous les noms.
\end{corrige}

\courssub{Production orale guidée : présenter un projet de reboisement}

\begin{experience}[Présentation orale : 1 minute — „Ein Projekt für mein Viertel / meine Schule“]
\textbf{Situation :} vous représentez le club environnement de votre lycée (par exemple le Lycée de Nkolbisson à Yaoundé ou un collège de Bonamoussadi à Douala). Vous présentez le projet „Un arbre par élève“ lors de la Journée mondiale de l'environnement (5 juin).

\textbf{Structure guidée (4 étapes = 4 phrases types) :}
\begin{enumerate}
    \item \textbf{Accroche + sujet :} \all{Guten Tag! Ich stelle unser Projekt „Ein Baum für jeden Schüler“ vor.} (Bonjour ! Je présente notre projet « Un arbre pour chaque élève ».)
    \item \textbf{But (préposition + accusatif) :} \all{Wir pflanzen Bäume \textbf{für} ein besseres Klima und \textbf{gegen} die Erosion.} (Nous plantons des arbres pour un meilleur climat et contre l'érosion.)
    \item \textbf{Action (impératif ou forme wir) :} \all{Helft mit! \textbf{Pflanzt} einen Baum im Schulhof oder zu Hause.} (Aidez-nous ! Plantez un arbre dans la cour ou à la maison.)
    \item \textbf{Appel final :} \all{Meiner Meinung nach machen wir unsere Schule grüner. \textbf{Deshalb}: Machen Sie mit!} (À mon avis, nous rendons notre école plus verte. C'est pourquoi : participez !)
\end{enumerate}

\textbf{Critères de réussite (grille d'évaluation orale) :}
\begin{itemize}
    \item \cle{Prononciation} claire : Umlaute \all{ä, ö, ü}, \all{z} prononcé « ts », \all{sch} prononcé « ch », \all{eu} prononcé « oï ».
    \item \cle{Grammaire} : au moins 1 impératif correct + 1 préposition + accusatif correcte.
    \item \cle{Vocabulaire} : au moins 5 mots du thème (\all{der Setzling, die Aktion, der Schatten, die Luft, die Gemeinschaft}).
    \item \cle{Fluidité} : pas de longs silences, utilisation de \cle{zuerst, dann, außerdem}.
\end{itemize}
\end{experience}

\courssub{Grille d'auto-évaluation : suis-je prêt pour le contrôle ?}

Remplissez cette grille \textbf{honnêtement} après avoir fait l'exercice écrit et la production orale. Cochez la colonne qui correspond à votre niveau de maîtrise.

\begin{aretenir}[Check-list des compétences — Leçon ALA10]
\begin{center}
\begin{tabular}{|p{4.5cm}|c|c|p{4cm}|}\hline
\rowcolor{popblueL} \textbf{Compétence ciblée} & \textbf{Maîtrisé} & \textbf{À revoir} & \textbf{Action de révision} \\ \hline
Vocabulaire du thème (\all{der Wald, pflanzen, schützen, der Naturpark, die Umwelt, die Artenvielfalt, der Setzling}) & & & Refaire les fiches de mots avec articles et pluriels \\ \hline
\textbf{Impératif} : formes \cle{du / ihr / Sie} régulières et irrégulières (\all{sein} : \all{Sei! / Seien Sie!}) & & & Conjuguer 5 verbes forts à l'impératif \\ \hline
\textbf{Prépositions + accusatif} : \cle{durch, für, gegen, ohne, um} + article correct (\all{den, die, das, die}) & & & Refaire le tableau de déclinaison + écrire 5 phrases \\ \hline
\textbf{Ordre des mots (V2)} : verbe en position 2 après un connecteur (\all{Deshalb, Meiner Meinung nach, Zuerst}) & & & Transformer 5 phrases en commençant par un connecteur \\ \hline
\textbf{Subordonnée avec « dass »} : verbe à la fin (\all{Ich denke, dass der Wald wichtig ist.}) & & & Écrire 3 opinions avec \cle{dass} \\ \hline
\textbf{Commentaire} : structure en 4 étapes + citations du texte + connecteurs logiques & & & Refaire le commentaire modèle à l'oral \\ \hline
\textbf{Production écrite} : paragraphe de 6 à 8 phrases, 3 impératifs, 2 prépositions + accusatif, connecteurs & & & Refaire l'exercice écrit en temps limité (15 min) \\ \hline
\end{tabular}
\end{center}
\end{aretenir}

\begin{savaistu}[Le commentaire de texte au Baccalauréat camerounais]
Au Baccalauréat série A (LV2 allemand), l'épreuve écrite comporte souvent un \textbf{commentaire de texte} dont la notation valorise trois grandes compétences :
\begin{itemize}
    \item \textbf{Compréhension :} restitution fidèle des idées (résumé), réponses précises aux questions de détail.
    \item \textbf{Langue et grammaire :} correction morphosyntaxique (cas, temps, ordre des mots), richesse lexicale (champ thématique), absence de calques du français (par exemple : ne jamais traduire « il y a » par \all{es hat} au lieu de \all{es gibt}).
    \item \textbf{Organisation et méthode :} respect de la structure en 4 étapes, utilisation de connecteurs logiques (\cle{zuerst, dann, außerdem, deshalb, schließlich}), cohérence de l'avis personnel par rapport au texte.
\end{itemize}
\cle{Astuce de correcteur} : un paragraphe « avis personnel » sans lien explicite avec le texte (pas de \cle{deshalb}, pas de référence du type \all{Wie im Text gesagt...}) est toujours pénalisé. \textbf{Ancrez toujours votre opinion dans le texte !}
\end{savaistu}

\newpage

\coursec{Activité d'intégration}

\courssub{Objectifs de l'activité}

À la fin de cette activité d'intégration, tu seras capable de :
\begin{itemize}
\item réinvestir le \cle{vocabulaire thématique} de la forêt et de la protection de l'environnement (\all{der Wald}, \all{der Baum}, \all{pflanzen}, \all{schützen}, \all{die Umwelt}, \all{der Naturpark}) dans une production authentique ;
\item former et employer correctement l'\cle{impératif} (2\up{e} pers. du pluriel et 1\up{re} pers. du pluriel) pour lancer un appel à l'action ;
\item utiliser les \cle{prépositions suivies de l'accusatif} (\all{durch}, \all{für}, \all{ohne}, \all{um}, \all{gegen}) dans des slogans et des phrases de sensibilisation ;
\item coopérer en groupe, présenter oralement un projet en deux minutes et évaluer la production des autres selon des critères précis.
\end{itemize}

\courssub{Situation}

Ton lycée de Yaoundé participe au programme d'échange scolaire avec le \emph{Gymnasium am Stadtpark} de Vienne (Autriche). Chaque année, les deux établissements réalisent un projet commun. Cette année, le thème choisi est : \emph{« Un arbre pour l'avenir »}. Le lycée autrichien a envoyé une lettre à ta classe pour proposer une campagne commune de reboisement.

\begin{textebox}[Lettre du partenaire autrichien]
Liebe Freunde in Yaoundé,\\
wir planen eine Aktion für die Umwelt: „Ein Baum für die Zukunft!“ In Österreich pflanzen viele Schulen Bäume in den Naturparks. Und ihr? Habt ihr einen Garten oder einen Wald in der Nähe von eurer Schule?\\
Macht bitte ein Plakat für eure Kampagne! Das Plakat braucht einen Slogan, schöne Wörter zum Thema Wald und eine Einladung für alle Schüler. Schickt uns Fotos von euren Plakaten! Wir zeigen sie dann in unserer Schule in Wien.\\
Viele Grüße\\
Eure Partnerklasse aus Wien
\end{textebox}

\textbf{Glossaire :} \all{die Aktion, -en} : l'action ; \all{in der Nähe von} : à proximité de ; \all{das Plakat, -e} : l'affiche ; \all{die Einladung, -en} : l'invitation ; \all{zeigen} : montrer.

Ta classe relève le défi : chaque groupe de 3 ou 4 élèves crée une affiche de campagne (ou un court message radio de 30 secondes) pour sensibiliser les élèves du lycée au reboisement de la cour de l'école et du quartier. Les meilleures affiches seront photographiées et envoyées aux partenaires de Vienne.

\courssub{Consignes}

Travaillez par groupes de 3 ou 4. Suivez les étapes dans l'ordre.

\begin{enumerate}
\item \textbf{Compréhension.} Lisez la lettre de la classe partenaire. Soulignez les informations importantes : Quel est le nom de la campagne ? Que demandent les partenaires autrichiens ? Que vont-ils faire des photos ?
\item \textbf{Lexique.} En groupe, choisissez \cle{cinq mots} du vocabulaire thématique de la leçon (par exemple \all{der Wald}, \all{der Baum}, \all{pflanzen}, \all{die Umwelt}, \all{schützen}, \all{der Naturpark}, \all{der Müll}, \all{die Luft}, \all{das Wasser}). Écrivez chaque mot avec son article et son pluriel sur un brouillon.
\item \textbf{Grammaire 1 — l'impératif.} Rédigez un \cle{slogan} à l'impératif, 2\up{e} personne du pluriel. Modèle : \all{Pflanzt Bäume für die Zukunft!} « Plantez des arbres pour l'avenir ! » Vérifiez la forme : radical du verbe + \all{-t} (attention aux verbes à particule : \all{Macht mit!}).
\item \textbf{Grammaire 2 — prépositions + accusatif.} Écrivez \cle{deux phrases} avec des prépositions suivies de l'accusatif : \all{durch den Wald}, \all{für die Umwelt}, \all{ohne Müll}, \all{um die Schule}, \all{gegen die Umweltverschmutzung}. Vérifiez la déclinaison : article accusatif (\all{den}, \all{die}, \all{das}).
\item \textbf{Invitation.} Ajoutez une invitation à la 1\up{re} personne du pluriel : \all{Kommt mit uns!} ou \all{Helfen wir der Natur!} « Venez avec nous ! / Aidons la nature ! »
\item \textbf{Mise en page.} Réalisez l'affiche sur une grande feuille : le slogan en grand, les cinq mots du lexique bien visibles, les phrases lisibles, un dessin simple (arbre, forêt, école).
\item \textbf{Présentation orale.} Préparez une présentation de \cle{deux minutes} : chaque membre du groupe parle. Présentez le slogan, lisez les phrases, expliquez votre choix de mots.
\item \textbf{Vote.} La classe vote pour la meilleure campagne selon trois critères : correction de la langue, richesse du vocabulaire, force du message.
\end{enumerate}

\begin{attention}[Pièges à éviter]
\begin{itemize}
\item N'oublie pas la majuscule à tous les noms : \all{der Baum}, \all{die Umwelt}.
\item À l'impératif pluriel, le verbe se termine en \all{-t} : \all{pflanzt}, \all{schützt}, \all{kommt} — pas de \all{-en} !
\item Après \all{für}, \all{durch}, \all{ohne}, \all{um}, \all{gegen} : toujours l'accusatif. On écrit \all{für die Umwelt} et \all{durch den Wald}, jamais \all{für der Umwelt}.
\item Verbe à particule séparable : la particule va à la fin : \all{Macht mit!} « Participez ! »
\end{itemize}
\end{attention}

\courssub{Ressources à mobiliser}

\begin{aretenir}[Boîte à outils de la campagne]
\textbf{Vocabulaire :} \all{der Wald, die Wälder} (la forêt) ; \all{der Baum, die Bäume} (l'arbre) ; \all{pflanzen} (planter) ; \all{schützen} (protéger) ; \all{die Umwelt} (l'environnement) ; \all{der Naturpark, -s} (le parc naturel) ; \all{die Luft} (l'air) ; \all{das Wasser} (l'eau) ; \all{der Müll} (les déchets) ; \all{die Zukunft} (l'avenir).

\textbf{Impératif (2\up{e} pers. pluriel) :} radical + \all{-t} : \all{pflanzen} $\rightarrow$ \all{Pflanzt!} ; \all{helfen} $\rightarrow$ \all{Helft!} ; \all{sein} $\rightarrow$ \all{Seid!} « Soyez ! »

\textbf{Impératif (1\up{re} pers. pluriel) :} verbe à l'infinitif + \all{wir} : \all{Helfen wir!} « Aidons ! » ; \all{Gehen wir!} « Allons-y ! »

\textbf{Prépositions + accusatif :} \all{durch den Wald} (à travers la forêt) ; \all{für die Umwelt} (pour l'environnement) ; \all{ohne Müll} (sans déchets) ; \all{um die Schule} (autour de l'école) ; \all{gegen die Umweltverschmutzung} (contre la pollution).
\end{aretenir}

\begin{center}
\begin{tabularx}{\linewidth}{|X|X|X|}
\hline
\rowcolor{popblueL} Français & Deutsch & Remarque \\
\hline
Plantez des arbres ! & Pflanzt Bäume! & impératif pluriel en \all{-t} \\
\hline
Protégez l'environnement ! & Schützt die Umwelt! & \all{die Umwelt} : accusatif féminin, article inchangé \\
\hline
Marchez à travers la forêt ! & Geht durch den Wald! & \all{durch} + accusatif : \all{der} $\rightarrow$ \all{den} \\
\hline
Aidons la nature ! & Helfen wir der Natur! & \all{helfen} + datif : attention, exception ! \\
\hline
Participez ! & Macht mit! & particule \all{mit} à la fin \\
\hline
\end{tabularx}
\end{center}

\courssub{Proposition de production et corrigé commenté}

\begin{exoresolu}[Affiche modèle du groupe 3]
Énoncé : rédiger l'affiche complète d'une campagne de reboisement pour le lycée de Yaoundé, avec slogan, cinq mots du lexique, deux phrases avec prépositions + accusatif et une invitation.

\tcblower

\textbf{Texte de l'affiche :}

\all{PFLANZT BÄUME – SCHÜTZT DIE UMWELT!}

\all{Unser Wald – unsere Zukunft!}

\all{der Baum — der Wald — die Umwelt — pflanzen — schützen}

\all{Wir pflanzen zwanzig Bäume um die Schule. Ohne Bäume gibt es keine saubere Luft! Geht mit uns durch den Garten und helft beim Pflanzen!}

\all{Kommt mit uns! Helfen wir der Natur!}

\all{Eure Klasse 1\up{re} A, Lycée de Yaoundé}

\textbf{Traduction :} « Plantez des arbres — protégez l'environnement ! Notre forêt — notre avenir ! l'arbre — la forêt — l'environnement — planter — protéger. Nous plantons vingt arbres autour de l'école. Sans arbres, il n'y a pas d'air propre ! Venez avec nous dans le jardin et aidez à planter ! Venez avec nous ! Aidons la nature ! »

\textbf{Commentaire des choix de langue :}
\begin{itemize}
\item \all{Pflanzt} et \all{Schützt} : impératif pluriel correct (radical + \all{-t}), en majuscules sur l'affiche pour l'effet visuel.
\item \all{um die Schule} : préposition \all{um} + accusatif, article féminin inchangé.
\item \all{Ohne Bäume} : \all{ohne} + accusatif pluriel (pas d'article devant le pluriel indéfini).
\item \all{durch den Garten} : \all{durch} + accusatif masculin : \all{der} devient \all{den}.
\item \all{Kommt mit uns!} : invitation à l'impératif ; \all{Helfen wir der Natur!} : 1\up{re} pers. du pluriel, avec \all{helfen} + datif (\all{der Natur}), structure vue comme exception.
\item Cinq mots du lexique sont présents, avec articles pour les noms.
\end{itemize}
\end{exoresolu}

\courssub{Bilan de l'activité}

Cette activité t'a permis de mobiliser ensemble toutes les compétences de la leçon :
\begin{itemize}
\item \textbf{Compréhension de l'écrit :} tu as lu et exploité une lettre authentique de partenaires germanophones ;
\item \textbf{Lexique :} tu as réemployé au moins cinq mots du thème « forêt et environnement » avec leurs articles ;
\item \textbf{Grammaire :} tu as formé des slogans à l'impératif et construit des phrases avec les prépositions suivies de l'accusatif ;
\item \textbf{Expression orale :} tu as présenté ton affiche en deux minutes, avec une prononciation claire ;
\item \textbf{Coopération :} tu as travaillé en équipe et évalué les productions des autres groupes selon des critères objectifs.
\end{itemize}

\begin{experience}[Pour aller plus loin]
Enregistre avec ton groupe un message radio de 30 secondes pour la radio du lycée : commence par le slogan, donne deux arguments (\all{Die Bäume geben uns Schatten und saubere Luft.}) et termine par l'invitation \all{Kommt mit uns!}. Les enregistrements seront envoyés à la classe partenaire de Vienne avec les photos des affiches.
\end{experience}

\newpage

\coursec{Méthodes et exercices résolus}

\courssub{Méthodes et exercices résolus}

\begin{methode}[Compréhension globale : identifier le sujet, les acteurs et l'intention du texte]
    \textbf{Objectif :} saisir le sens général sans traduire mot à mot.
    \begin{enumerate}
        \item \textbf{Lire le titre et les images :} \all{Aufforstung und Ökotourismus in Kamerun} → sujet : reboisement + écotourisme ; lieu : le Cameroun.
        \item \textbf{Repérer les mots-clés transparents (mots amis) :} \all{das Projekt, die Natur, das Klima, der Tourismus, der Partner, die Region, aktiv}.
        \item \textbf{Identifier la structure :} introduction (constat) → actions (reboisement) → conséquences (tourisme, revenus) → perspectives.
        \item \textbf{Formuler l'idée principale en une phrase en français.}
    \end{enumerate}
    \textbf{Réflexe :} ne pas s'arrêter sur un mot inconnu ; le deviner par le contexte (famille de mots, prépositions, mots amis).
\end{methode}

\begin{exoresolu}[Compréhension globale : questions type examen sur le texte « Aufforstung und Ökotourismus in Kamerun »]
    \textbf{Énoncé :} après lecture du texte, répondez en allemand par des phrases complètes :
    \begin{enumerate}
        \item \all{Was ist das Hauptthema des Textes?}
        \item \all{Welche zwei Hauptaktivitäten werden in Kamerun beschrieben?}
        \item \all{Wer profitiert von den Projekten? (Nennen Sie zwei Gruppen.)}
        \item \all{Warum ist der Wald für das Klima wichtig? (Ein Satz.)}
    \end{enumerate}
    \tcblower
    \textbf{Solution détaillée :}
    \begin{enumerate}
        \item \textbf{Analyse de la question :} \all{Was ist das Hauptthema?} (Quel est le sujet principal ?). On cherche la réponse dans le premier paragraphe.
        \item \textbf{Réponse 1 :} \all{Das Hauptthema ist die Aufforstung und der Ökotourismus in Kamerun.} « Le sujet principal est le reboisement et l'écotourisme au Cameroun. » \textit{Justification :} reprise du titre et du premier paragraphe.
        \item \textbf{Réponse 2 :} \all{Die zwei Hauptaktivitäten sind das Pflanzen von Bäumen (die Aufforstung) und der sanfte Tourismus (der Ökotourismus).} « Les deux activités principales sont la plantation d'arbres (le reboisement) et le tourisme doux (l'écotourisme). » \textit{Justification :} le verbe \all{pflanzen} et le nom \all{der Ökotourismus} reviennent souvent. On utilise la nominalisation \all{das Pflanzen} (verbe nominalisé → neutre, majuscule).
        \item \textbf{Réponse 3 :} \all{Die lokale Bevölkerung und die Umwelt profitieren von den Projekten. Auch die Touristen erleben die Natur.} « La population locale et l'environnement profitent des projets. Les touristes aussi vivent la nature. » \textit{Justification :} repérer \all{profitieren von + Datif} : \all{von den Projekten}. \all{die Bevölkerung} (féminin singulier), \all{die Umwelt} (féminin singulier, sans pluriel).
        \item \textbf{Réponse 4 :} \all{Der Wald speichert CO2 und schützt das Klima.} « La forêt stocke le CO2 et protège le climat. » \textit{Justification :} le verbe \all{speichern} (stocker) + accusatif ; \all{schützen} + accusatif \all{das Klima} (neutre, inchangé à l'accusatif).
    \end{enumerate}
    \textbf{Conclusion :} des réponses courtes, syntaxiquement correctes (verbe conjugué en 2\up{e} position), qui réemploient le vocabulaire du texte.
\end{exoresolu}

\begin{methode}[Compréhension détaillée et lexique : exploiter les familles de mots et le contexte]
    \textbf{Objectif :} déduire le sens précis des mots inconnus et structurer le lexique thématique.
    \begin{enumerate}
        \item \textbf{Surligner les verbes d'action :} \all{pflanzen, schützen, erhalten, fördern, besuchen, unterstützen}.
        \item \textbf{Construire le tableau des familles de mots (\all{Wortfamilien}) :} verbe → nom (article + pluriel) → adjectif/participe.
        \item \textbf{Repérer les connecteurs logiques :} \all{deshalb} (c'est pourquoi), \all{außerdem} (de plus), \all{aber} (mais), \all{damit} (afin que), \all{weil} (parce que) → relations de cause, conséquence, but.
        \item \textbf{Traduire les segments clés (version) :} ne traduire que les phrases porteuses de l'argumentation.
    \end{enumerate}
    \textbf{Astuce :} pour chaque nom nouveau, apprendre \textbf{article + pluriel + sens} : \all{der Wald, die Wälder} ; \all{der Naturpark, die Naturparks} ; \all{der Baum, die Bäume}.
\end{methode}

\begin{exoresolu}[Travail lexical : familles de mots et version ciblée]
    \textbf{Énoncé :}
    \begin{enumerate}
        \item Complétez le tableau des familles de mots à partir du texte.
        \item Traduisez en français le passage suivant :\\
        \all{„Durch die Aufforstung entstehen neue Lebensräume für Tiere. Die Einheimischen finden Arbeit als Ranger oder Führer. Damit wird die Armut in der Region verringert.“}
    \end{enumerate}
    \begin{center}
    \begin{tabular}{|l|l|l|l|}\hline
    \rowcolor{popblueL} \textbf{Verbe} & \textbf{Nom (Art. + Pl.)} & \textbf{Participe II} & \textbf{Sens (FR)} \\ \hline
    \all{pflanzen} & \all{die Pflanzung, -en} & \all{gepflanzt} & planter / plantation \\ \hline
    \all{schützen} & \all{der Schutz} (sans Pl.) & \all{geschützt} & protéger / protection \\ \hline
    \all{erhalten} & \all{der Erhalt} (sans Pl.) & \all{erhalten} & préserver / préservé \\ \hline
    \all{fördern} & \all{die Förderung, -en} & \all{gefördert} & encourager / soutien \\ \hline
    \all{besuchen} & \all{der Besuch, -e} & \all{besucht} & visiter / visite \\ \hline
    \end{tabular}
    \end{center}
    \tcblower
    \textbf{Solution détaillée :}
    \begin{enumerate}
        \item \textbf{Tableau complété :} voir ci-dessus. \textit{Notes :} \all{der Schutz} et \all{der Erhalt} sont des noms d'action sans pluriel usuel. \all{die Aufforstung} (le reboisement) vient du verbe \all{aufforsten} (particule inséparable : \all{er forstet auf, er hat aufgeforstet}). Attention : \all{erhalten} est un verbe fort : \all{er erhält, er hat erhalten}.
        \item \textbf{Version pas à pas :}
        \all{Durch die Aufforstung} (Grâce au reboisement — \all{durch} + accusatif ; \all{die Aufforstung} est féminin, donc inchangé)
        \all{entstehen neue Lebensräume für Tiere.} (naissent de nouveaux habitats pour les animaux — verbe intransitif, sujet : \all{neue Lebensräume}, d'où le verbe au pluriel ; \all{der Lebensraum, die Lebensräume}).
        \all{Die Einheimischen} (Les habitants du pays / la population locale — adjectif nominalisé : majuscule, pluriel)
        \all{finden Arbeit als Ranger oder Führer.} (trouvent du travail comme gardes forestiers ou guides.)
        \all{Damit wird die Armut in der Region verringert.} (Ainsi, la pauvreté dans la région est réduite — \textbf{passif au présent} : \all{werden} + participe II \all{verringert} ; sujet : \all{die Armut}, sans pluriel).
        \item \textbf{Points de vigilance :} \all{die Einheimischen} (adjectif nominalisé, majuscule) ; \all{der Ranger, die Ranger} (emprunt à l'anglais) ; \all{der Führer, die Führer} (le guide) ; \all{die Arbeit} (le travail, sans article ici car sens général).
    \end{enumerate}
\end{exoresolu}

\begin{methode}[Grammaire : former et employer l'impératif (der Imperativ)]
    \textbf{Objectif :} donner un ordre, un conseil, une invitation ou exprimer une interdiction.
    \begin{enumerate}
        \item \textbf{Identifier le destinataire :} \all{du} (tu), \all{ihr} (vous, familier), \all{Sie} (vous, formel), \all{wir} (nous, invitation).
        \item \textbf{Appliquer la formation selon la personne :}
        \begin{itemize}
            \item \all{du} : radical du présent (sans \all{-st}), en général \textbf{sans} \all{-e} final : \all{pflanz!}, \all{komm!}. Le \all{-e} est obligatoire seulement après un radical en \all{-d, -t, -ig} ou en consonne + \all{-m/-n} : \all{arbeite!}, \all{öffne!}, \all{rechne!}, \all{entschuldige!}. Les verbes forts à voyelle modifiée \all{e→i/ie} gardent la modification, sans \all{-e} : \all{lies!} (< \all{lesen}), \all{hilf!} (< \all{helfen}), \all{nimm!} (< \all{nehmen}). En revanche, les verbes à Umlaut (\all{fahren} → \all{du fährst}) perdent l'Umlaut : \all{fahr!}.
            \item \all{ihr} : forme du présent à \all{ihr}, sans pronom : \all{pflanzt!}, \all{seid!} (irrégulier).
            \item \all{Sie} : forme du présent à \all{Sie} + pronom \all{Sie} placé après le verbe : \all{pflanzen Sie!}, \all{seien Sie!} (irrégulier).
            \item \all{wir} : forme du présent à \all{wir} + pronom \all{wir} après le verbe : \all{pflanzen wir!} (Plantons !).
        \end{itemize}
        \item \textbf{Position du verbe :} en \textbf{position 1} (début de phrase), suivi du point d'exclamation.
        \item \textbf{Négation :} \all{nicht} ou \all{kein-} : \all{Wirf keinen Müll weg!} (Ne jette pas de déchets !).
        \item \textbf{Verbes à particule séparable :} la particule va \textbf{en fin de phrase} : \all{Pass auf!} (Fais attention !), \all{Werfen Sie den Müll weg!}.
    \end{enumerate}
    \textbf{Emplois :} consignes (\all{Lest den Text!}), conseils (\all{Schütze den Wald!}), interdictions (\all{Werft keinen Müll weg!}).
\end{methode}

\begin{exoresolu}[Transformation : mettre les phrases à l'impératif (toutes les personnes)]
    \textbf{Énoncé :} transformez les phrases à l'indicatif en impératif pour la personne indiquée. Attention aux verbes irréguliers et aux verbes à particule séparable.
    \begin{enumerate}
        \item \all{Du pflanzt einen Baum.} → (\all{du})
        \item \all{Sie schützen die Umwelt.} → (\all{Sie})
        \item \all{Ihr besucht den Naturpark.} → (\all{ihr})
        \item \all{Wir sammeln Müll im Wald.} → (\all{wir})
        \item \all{Du liest das Schild.} → (\all{du}) \textit{(verbe fort : voyelle modifiée)}
        \item \all{Sie werfen den Müll weg.} → (\all{Sie}) \textit{(verbe à particule séparable : \all{wegwerfen})}
        \item \all{Du hilfst dem Ranger.} → (\all{du}) \textit{(verbe construit avec le datif)}
    \end{enumerate}
    \tcblower
    \textbf{Solution détaillée :}
    \begin{enumerate}
        \item \all{Pflanz einen Baum!} (Radical \all{pflanz-}, sans \all{-e} final.) « Plante un arbre ! »
        \item \all{Schützen Sie die Umwelt!} (Forme \all{Sie} + pronom après le verbe.) « Protégez l'environnement ! »
        \item \all{Besucht den Naturpark!} (Forme \all{ihr} sans pronom ; \all{den Naturpark} : accusatif masculin.) « Visitez le parc naturel ! »
        \item \all{Sammeln wir Müll im Wald!} (Forme \all{wir} + pronom.) « Ramassons des déchets dans la forêt ! »
        \item \all{Lies das Schild!} (Verbe fort \all{lesen} → \all{du liest} → impératif \all{lies} : on garde \all{ie}, sans \all{-e}.) « Lis le panneau ! »
        \item \all{Werfen Sie den Müll weg!} (Verbe \all{wegwerfen} : verbe conjugué en position 1, particule \all{weg} \textbf{en fin de phrase}.) « Jetez les déchets ! »
        \item \all{Hilf dem Ranger!} (\all{helfen} + \textbf{datif} : \all{dem Ranger} ; radical \all{hilf-} avec voyelle modifiée, sans \all{-e}.) « Aide le garde forestier ! »
    \end{enumerate}
    \textbf{Conclusion :} il faut maîtriser les 4 formes, la place de la particule séparable (à la fin) et la conservation de la voyelle modifiée \all{e→i/ie} à la forme \all{du}.
\end{exoresolu}

\begin{methode}[Grammaire : maîtriser les prépositions suivies de l'accusatif (Präpositionen mit Akkusativ)]
    \textbf{Objectif :} utiliser correctement \all{durch, für, gegen, ohne, um} et \all{entlang}.
    \begin{enumerate}
        \item \textbf{Mémoriser la liste (moyen mnémotechnique) :} \cle{FUDOG} : \all{Für, Um, Durch, Ohne, Gegen} (+ \all{entlang}, postposée).
        \item \textbf{Appliquer la déclinaison à l'accusatif :} seul le \textbf{masculin} change : \all{der → den}, \all{ein → einen}, \all{mein → meinen}, \all{kein → keinen}. Féminin, neutre et pluriel : \textbf{inchangés} (\all{die, das, die}).
        \item \textbf{Cas particuliers :}
        \begin{itemize}
            \item \all{entlang} se place \textbf{après} le nom : \all{den Fluss entlang} (le long du fleuve).
            \item \all{bis} s'emploie souvent sans article : \all{bis morgen}, \all{bis Yaoundé}.
            \item \all{ohne} + nom sans article quand le sens est général : \all{ohne Guide} (sans guide), \all{ohne Geld}.
        \end{itemize}
        \item \textbf{Ne pas confondre} avec les prépositions à deux cas (\all{in, auf, an, über, unter, vor, hinter, neben, zwischen}) : accusatif pour la direction (\all{wohin?}), datif pour la position (\all{wo?}) : \all{Wir gehen in den Wald.} (nous allons dans la forêt) / \all{Wir wandern im Wald.} (nous marchons dans la forêt).
    \end{enumerate}
    \textbf{Réflexe :} avant d'écrire le nom, se demander : « Quelle préposition ? » → si c'est \cle{FUDOG}, l'accusatif est obligatoire.
\end{methode}

\begin{exoresolu}[Application : prépositions accusatives dans un texte à trous]
    \textbf{Énoncé :} complétez avec la préposition convenable (\all{durch, für, gegen, ohne, um, entlang}) et mettez le nom entre parenthèses au bon cas. Traduisez chaque phrase.
    \begin{enumerate}
        \item \all{Wir wandern \_\_\_\_\_\_ (der Fluss) \_\_\_\_\_\_.} (entlang)
        \item \all{\_\_\_\_\_\_ (der Schutz) der Wälder pflanzen wir Bäume.} (für)
        \item \all{\_\_\_\_\_\_ (die Aufforstung) entsteht neuer Wald.} (durch)
        \item \all{Man kann den Park nicht \_\_\_\_\_\_ (der Guide) besuchen.} (ohne)
        \item \all{Wir kämpfen \_\_\_\_\_\_ (die Abholzung).} (gegen)
        \item \all{Viele Touristen reisen \_\_\_\_\_\_ (das Land), um die Natur zu sehen.} (durch)
    \end{enumerate}
    \tcblower
    \textbf{Solution détaillée :}
    \begin{enumerate}
        \item \all{Wir wandern den Fluss entlang.} (\all{entlang} se place \textbf{après} le nom ; \all{der Fluss} masculin → \all{den Fluss}.) « Nous marchons le long du fleuve. »
        \item \all{Für den Schutz der Wälder pflanzen wir Bäume.} (\all{für} + accusatif ; \all{der Schutz} masculin → \all{den Schutz} ; \all{der Wälder} : génitif pluriel, complément du nom.) « Pour la protection des forêts, nous plantons des arbres. »
        \item \all{Durch die Aufforstung entsteht neuer Wald.} (\all{durch} + accusatif ; \all{die Aufforstung} féminin → inchangé ; \all{neuer Wald} : nominatif masculin sans article, adjectif en \all{-er}.) « Grâce au reboisement, une nouvelle forêt naît. »
        \item \all{Man kann den Park nicht ohne Guide besuchen.} (\all{ohne} + accusatif ; ici sens général, donc sans article. Avec article : \all{ohne einen Guide}.) « On ne peut pas visiter le parc sans guide. »
        \item \all{Wir kämpfen gegen die Abholzung.} (\all{gegen} + accusatif ; \all{die Abholzung} féminin → inchangé.) « Nous luttons contre la déforestation. »
        \item \all{Viele Touristen reisen durch das Land, um die Natur zu sehen.} (\all{durch} + accusatif ; \all{das Land} neutre → inchangé ; \all{um ... zu} + infinitif exprime le but.) « Beaucoup de touristes voyagent à travers le pays pour voir la nature. »
    \end{enumerate}
    \textbf{Conclusion :} vérifier systématiquement le genre du nom : seul le masculin change (\all{der → den}, \all{ein → einen}). Attention à \all{entlang}, placée après le nom.
\end{exoresolu}

\begin{attention}[Piège hors programme : \all{um ... willen}]
    La tournure \all{um ... willen} (pour l'amour de, pour le bien de) gouverne le \textbf{génitif} et non l'accusatif : \all{um der Natur willen} (pour le bien de la nature). Elle est signalée ici à titre culturel, mais à l'examen de Première, préférez les tournures du programme : \all{für die Natur} (préposition + accusatif) ou \all{um die Natur zu schützen} (\all{um ... zu} + infinitif pour exprimer le but).
\end{attention}

\begin{methode}[Expression écrite guidée : rédiger un paragraphe cohérent (résumé ou avis personnel)]
    \textbf{Objectif :} produire un texte de 60 à 80 mots, structuré, en réemployant le lexique et la grammaire de la leçon.
    \begin{enumerate}
        \item \textbf{Planifier (3 idées) :} 1. Constat ou problème. 2. Action ou solution (reboisement, parc naturel). 3. Résultat ou bénéfice (climat, emplois, tourisme).
        \item \textbf{Choisir les connecteurs :} \all{zuerst / zunächst} (d'abord), \all{dann / danach} (ensuite), \all{außerdem} (de plus), \all{deshalb} (c'est pourquoi), \all{schließlich} (finalement), \all{meiner Meinung nach} (à mon avis).
        \item \textbf{Contrôler la syntaxe (ordre des mots) :}
        \begin{itemize}
            \item Phrase principale : verbe conjugué en position 2 (si un complément ouvre la phrase, le sujet passe après le verbe : \all{Deshalb kommen viele Touristen.}).
            \item Subordonnée (\all{weil, dass, wenn}) : verbe conjugué à la fin : \all{..., weil der Wald das Klima schützt.}
            \item But : \all{um ... zu} + infinitif à la fin : \all{..., um die Natur zu schützen.}
        \end{itemize}
        \item \textbf{Vérifier :} majuscules à tous les noms, accusatif après \cle{FUDOG}, conjugaison des verbes forts, particules séparables en fin de phrase.
    \end{enumerate}
    \textbf{Modèle de début :} \all{In Kamerun gibt es viele Projekte für die Aufforstung. Durch das Pflanzen neuer Bäume wird der Wald geschützt. Deshalb kommen Touristen in die Naturparks. Das bringt Geld für die Region.} « Au Cameroun, il existe de nombreux projets de reboisement. Grâce à la plantation de nouveaux arbres, la forêt est protégée. C'est pourquoi des touristes viennent dans les parcs naturels. Cela rapporte de l'argent à la région. »
\end{methode}

\begin{exoresolu}[Rédaction : « L'écotourisme, une chance pour le Cameroun ? » (production écrite, environ 80 mots)]
    \textbf{Énoncé :} rédigez un court paragraphe en allemand (environ 80 mots) sur le thème : \all{„Ökotourismus in Kamerun – Chance oder Risiko?“}. Utilisez au moins : 3 verbes au présent, 2 prépositions suivies de l'accusatif (\all{für, durch, gegen}), 1 impératif (conseil), 2 connecteurs logiques.
    \tcblower
    \textbf{Proposition de rédaction (corrigé type) :}\\
    \all{„Der Ökotourismus in Kamerun ist eine große Chance. Durch den Schutz der Wälder bleibt die Natur erhalten. Für die lokale Bevölkerung entstehen neue Jobs als Ranger oder Führer. Außerdem lernen die Touristen die Artenvielfalt kennen. Man muss aber aufpassen: Besuchen Sie die Parks mit Respekt! Werfen Sie keinen Müll weg! Nur so hilft der Tourismus gegen die Armut und für das Klima.“}\\
    « L'écotourisme au Cameroun est une grande chance. Grâce à la protection des forêts, la nature est préservée. Pour la population locale, de nouveaux emplois apparaissent, comme gardes forestiers ou guides. De plus, les touristes découvrent la biodiversité. Mais il faut faire attention : visitez les parcs avec respect ! Ne jetez pas de déchets ! C'est seulement ainsi que le tourisme aide contre la pauvreté et pour le climat. »
    \textbf{Analyse grammaticale du corrigé :}
    \begin{itemize}
        \item \textbf{Verbes au présent :} \all{ist, bleibt, entstehen, lernen ... kennen, muss, hilft}. (\all{entstehen} au pluriel car le sujet est \all{neue Jobs} ; \all{kennenlernen} est séparable : \all{lernen ... kennen}.)
        \item \textbf{Prépositions + accusatif :} \all{durch den Schutz} (masculin : \all{der Schutz → den Schutz}), \all{für die lokale Bevölkerung} (féminin, inchangé), \all{gegen die Armut} (féminin, inchangé), \all{für das Klima} (neutre, inchangé).
        \item \textbf{Impératifs :} \all{Besuchen Sie die Parks mit Respekt!} (forme \all{Sie}) ; \all{Werfen Sie keinen Müll weg!} (particule \all{weg} en fin de phrase ; \all{keinen Müll} : accusatif masculin de \all{kein-}).
        \item \textbf{Connecteurs :} \all{außerdem} (en position 1, verbe en position 2 : \all{Außerdem lernen die Touristen ...}) ; \all{nur so} (position 1, inversion : \all{Nur so hilft der Tourismus ...}).
        \item \textbf{Majuscules :} \all{Ökotourismus, Chance, Schutz, Wälder, Natur, Bevölkerung, Jobs, Ranger, Führer, Touristen, Artenvielfalt, Parks, Respekt, Müll, Tourismus, Armut, Klima}.
    \end{itemize}
    \textbf{Note :} ce texte obtient la note maximale car il respecte toutes les contraintes, emploie le vocabulaire du thème, applique correctement la grammaire ciblée et reste cohérent du début à la fin.
\end{exoresolu}

\begin{attention}[Les 5 erreurs qui coûtent des points à l'examen]
    \begin{enumerate}
        \item \textbf{Majuscules oubliées sur les noms :} \all{die umwelt, der baum, der naturpark} → faute majeure. \textit{Règle :} \cle{tout nom prend une majuscule en allemand} : \all{die Umwelt, der Baum, der Naturpark}.
        \item \textbf{Accusatif masculin non marqué après une préposition FUDOG :} \all{für der Schutz, durch der Wald} → faute de cas. \textit{Correct :} \all{für den Schutz, durch den Wald}. (Féminin, neutre, pluriel : inchangés : \all{für die Umwelt, für das Klima}.)
        \item \textbf{Particule séparable mal placée à l'impératif :} \all{Wegwerfen Sie den Müll!} → faute de structure. \textit{Correct :} \all{Werfen Sie den Müll weg!} (particule en \textbf{fin de phrase}).
        \item \textbf{Impératif \all{du} mal formé :} \all{Pflanze einen Baum!} (avec \all{-e} superflu) ou \all{Lese das Schild!} (perte de la voyelle modifiée). \textit{Correct :} \all{Pflanz einen Baum!}, \all{Lies das Schild!}. Et attention à l'Umlaut qui disparaît : \all{Fahr langsam!} (et non \all{Fähr!}).
        \item \textbf{Confusion entre \all{für} + accusatif et \all{um ... zu} + infinitif :} \all{Um den Schutz pflanzen wir Bäume.} → faute. \textit{Correct :} \all{Für den Schutz pflanzen wir Bäume.} (préposition + nom à l'accusatif) ou \all{Um den Wald zu schützen, pflanzen wir Bäume.} (\all{um ... zu} + infinitif en fin de subordonnée).
    \end{enumerate}
    \textbf{Conseil final :} relisez votre copie en traquant dans l'ordre : \textbf{1.} majuscules des noms ; \textbf{2.} accusatif masculin (\all{den / einen / keinen}) ; \textbf{3.} verbe conjugué en position 2 (ou à la fin en subordonnée) ; \textbf{4.} particules séparables en fin de phrase.
\end{attention}

\newpage

\coursec{Exercices}

\courssub{Je vérifie mes connaissances}

\begin{exercice}[QCM : le texte support]
Entoure la bonne réponse.
\begin{enumerate}
\item Le texte \all{Aufforstung und Ökotourismus in Kamerun} parle surtout :
\begin{enumerate}
\item du tourisme en Europe \item de la protection des forêts au Cameroun \item de la chasse
\end{enumerate}
\item \all{die Aufforstung} signifie :
\begin{enumerate}
\item la déforestation \item le reboisement \item la pollution
\end{enumerate}
\item \all{der Ökotourismus} est un tourisme qui :
\begin{enumerate}
\item détruit la nature \item respecte la nature et aide la population \item coûte très cher
\end{enumerate}
\item Dans le texte, les habitants des villages :
\begin{enumerate}
\item coupent tous les arbres \item plantent des arbres et guident les touristes \item vendent la forêt
\end{enumerate}
\end{enumerate}
\end{exercice}

\begin{exercice}[Vrai ou faux ? Justifie]
Réponds par \all{Richtig} ou \all{Falsch} et justifie ta réponse par une courte phrase en allemand tirée du texte support.
\begin{enumerate}
\item Die Wälder in Kamerun sind nicht in Gefahr.
\item Viele Dorfbewohner pflanzen neue Bäume.
\item Der Ökotourismus bringt den Dörfern kein Geld.
\item Die Touristen können im Naturpark Tiere beobachten.
\end{enumerate}
\end{exercice}

\begin{exercice}[Vocabulaire : trouve le mot]
Complète chaque phrase avec le mot qui convient : \all{der Wald, der Baum, pflanzen, der Naturpark, schützen, die Umwelt}.
\begin{enumerate}
\item Im \_\_\_\_\_\_\_\_\_\_\_\_ wachsen viele Bäume.
\item Wir müssen die \_\_\_\_\_\_\_\_\_\_\_\_ \_\_\_\_\_\_\_\_\_\_\_\_ : sie ist unser Leben.
\item Die Schüler \_\_\_\_\_\_\_\_\_\_\_\_ jedes Jahr zwanzig Bäume.
\item Ein \_\_\_\_\_\_\_\_\_\_\_\_ hat grüne Blätter und eine lange Wurzel.
\item Im \_\_\_\_\_\_\_\_\_\_\_\_ können Touristen Elefanten und Vögel sehen.
\end{enumerate}
\end{exercice}

\begin{exercice}[Conjugaison : l'impératif]
Complète le tableau avec la bonne forme de l'impératif.

\begin{center}
\begin{tabular}{|l|l|l|l|}
\hline
\rowcolor{popblueL} Infinitif & du & ihr & Sie \\ \hline
pflanzen & Pflanze \ldots ! & \ldots ! & \ldots ! \\ \hline
helfen & \ldots ! & Helft \ldots ! & \ldots ! \\ \hline
lesen & \ldots ! & \ldots ! & Lesen Sie \ldots ! \\ \hline
fahren & \ldots ! & \ldots ! & \ldots ! \\ \hline
sein & Sei \ldots ! & \ldots ! & \ldots ! \\ \hline
\end{tabular}
\end{center}
\end{exercice}

\courssub{Je m'entraîne}

\begin{exercice}[Mets les phrases à l'impératif]
Transforme chaque phrase à l'impératif à la personne indiquée. Attention aux mutations vocaliques !

Exemple : \all{Du pflanzt einen Baum.} $\rightarrow$ \all{Pflanze einen Baum!}
\begin{enumerate}
\item Du hilfst den Dorfbewohnern. (du)
\item Ihr lest den Text über den Naturpark. (ihr)
\item Sie fahren langsam durch den Wald. (Sie)
\item Du nimmst eine Tasche für den Müll. (du)
\item Ihr schützt die Umwelt. (ihr)
\item Sie bleiben auf dem Weg. (Sie)
\end{enumerate}
\end{exercice}

\begin{exercice}[Prépositions + accusatif]
Complète avec la bonne préposition (choisie parmi \all{durch, für, gegen, ohne, um}) et le bon article à l'accusatif.
\begin{enumerate}
\item Wir wandern \_\_\_\_\_\_ (durch) \_\_\_\_\_\_ Wald.
\item Sie kämpft \_\_\_\_\_\_ (gegen) \_\_\_\_\_\_ Abholzung.
\item Die Schüler sammeln Geld \_\_\_\_\_\_ (für) \_\_\_\_\_\_ Naturpark.
\item Er geht \_\_\_\_\_\_ (ohne) \_\_\_\_\_\_ Freund in den Wald.
\item Die Gruppe sitzt \_\_\_\_\_\_ (um) \_\_\_\_\_\_ Baum.
\end{enumerate}
\end{exercice}

\begin{textebox}[Unser Projekt]
\all{\_\_\_\_\_\_ (1) Bäume für die Zukunft! In unserer Schule in Yaoundé machen wir ein Projekt \_\_\_\_\_\_ (2) die Natur. Wir gehen \_\_\_\_\_\_ (3) den Schulgarten und pflanzen zehn \_\_\_\_\_\_ (4). So \_\_\_\_\_\_ (5) wir die Erde und die Luft. Die \_\_\_\_\_\_ (6) ist wichtig für alle Menschen in Kamerun.}
\end{textebox}

\begin{exercice}[Texte à trous]
Complète le texte « Unser Projekt » ci-dessus avec les mots suivants : \all{Bäume, Umwelt, durch, für, Pflanzt, schützen}.
\end{exercice}

\begin{exercice}[Appariements]
Relie chaque mot allemand à sa traduction française.

\begin{center}
\begin{tabular}{|l|l|}
\hline
\rowcolor{popblueL} Deutsch & Français \\ \hline
1. der Wald, die Wälder & a. planter \\ \hline
2. der Baum, die Bäume & b. protéger \\ \hline
3. pflanzen & c. l'environnement \\ \hline
4. schützen & d. la forêt \\ \hline
5. die Umwelt & e. le parc naturel \\ \hline
6. der Naturpark, die Naturparks & f. l'arbre \\ \hline
\end{tabular}
\end{center}
\end{exercice}

\courssub{Je me prépare au Bac}

\begin{textebox}[Ein Naturpark in Deutschland]
\all{Der Naturpark Schwarzwald liegt im Südwesten Deutschlands. Viele Touristen besuchen ihn jedes Jahr. Sie wandern durch den Wald, beobachten Vögel und schlafen in kleinen Hotels. Die Mitarbeiter des Parks sagen zu den Besuchern: „Bleiben Sie auf den Wegen! Nehmen Sie Ihren Müll mit! Pflanzen Sie einen Baum für die Zukunft!“ Die Touristen lernen viel über die Natur, und die Dörfer in der Region verdienen Geld mit dem Ökotourismus. So schützt der Park die Umwelt und hilft den Menschen.}

\emph{Glossaire : der Weg, die Wege : le chemin ; der Müll : les déchets ; verdienen : gagner.}
\end{textebox}

\begin{exercice}[Compréhension écrite type Bac]
Lis le texte « Ein Naturpark in Deutschland » ci-dessus, puis réponds aux questions \textbf{en allemand, par des phrases complètes}.

\textbf{Fragen zum Text :}
\begin{enumerate}
\item Wo liegt der Naturpark Schwarzwald?
\item Was machen die Touristen im Park? (Nenne zwei Aktivitäten.)
\item Was sagen die Mitarbeiter zu den Besuchern? (Nenne zwei Regeln.)
\item Warum ist der Ökotourismus gut für die Dörfer?
\item Was lernen die Touristen im Park?
\end{enumerate}

\textbf{Richtig oder Falsch ?} Justifie ta réponse par une citation du texte.
\begin{enumerate}
\item Die Touristen schlafen in großen Hotels.
\item Die Besucher müssen ihren Müll mitnehmen.
\item Der Ökotourismus bringt der Region kein Geld.
\end{enumerate}

\textbf{Fragen zur Sprache :}
\begin{enumerate}
\item Relève dans le texte deux impératifs et précise la personne (du, ihr ou Sie).
\item Relève deux prépositions suivies de l'accusatif et recopie le groupe prépositionnel complet.
\end{enumerate}
\end{exercice}

\begin{exercice}[Thème et version]
\textbf{A. Thème.} Traduis en allemand :
\begin{enumerate}
\item Plantez des arbres pour l'avenir ! (vous = \all{Sie})
\item Nous marchons à travers la forêt sans bruit.
\item Aide les villageois, mon ami ! (tu)
\item Elle lutte contre la déforestation au Cameroun.
\end{enumerate}

\textbf{B. Version.} Traduis en français le texte « Besuch im Naturpark » ci-dessous.
\end{exercice}

\begin{textebox}[Besuch im Naturpark]
\all{Kommen Sie in den Naturpark und entdecken Sie die Natur! Wandern Sie durch den grünen Wald und beobachten Sie die Tiere. Viele Familien fahren am Wochenende in den Park. Die Kinder lernen dort viel über die Bäume und die Vögel. Der Park kämpft gegen die Abholzung und pflanzt jedes Jahr neue Bäume. Besuchen Sie uns — die Natur wartet auf Sie!}

\emph{Glossaire : entdecken : découvrir ; die Abholzung : la déforestation ; warten auf : attendre.}
\end{textebox}

\begin{exercice}[Production écrite et jeu de rôle]
\textbf{A. Production écrite (80 à 100 mots).}

Sujet : \all{Was kannst du für die Umwelt in deiner Stadt tun?}

Rédige un court texte en allemand. Contraintes obligatoires :
\begin{itemize}
\item au moins \textbf{trois impératifs} (par exemple des conseils à tes camarades) ;
\item au moins \textbf{deux prépositions suivies de l'accusatif} (\all{durch, für, gegen, ohne, um}) ;
\item au moins cinq mots du vocabulaire du chapitre (\all{der Wald, der Baum, pflanzen, schützen, die Umwelt, der Naturpark}).
\end{itemize}

\textbf{B. Jeu de rôle oral (par binômes).}

Un élève joue le guide d'un parc naturel camerounais (par exemple le parc de Waza ou la réserve du Dja) ; l'autre joue un touriste allemand. Le guide donne au moins cinq consignes à l'impératif (\all{Bleiben Sie\ldots, Fotografieren Sie\ldots, Werfen Sie keinen Müll\ldots}) ; le touriste pose trois questions. Puis échangez les rôles.
\end{exercice}

\begin{methode}[Réussir la production écrite]
\begin{enumerate}
\item \textbf{Prépare} : liste cinq mots du chapitre et trois conseils à donner (planter, protéger, ne pas jeter de déchets).
\item \textbf{Rédige} : commence par une phrase d'introduction (\all{Die Umwelt ist wichtig für uns alle.}), puis donne tes conseils à l'impératif (\all{Pflanzt Bäume! Schützt die Tiere!}).
\item \textbf{Intègre} les prépositions + accusatif : \all{Wir kämpfen gegen die Abholzung. Wir sammeln Geld für den Naturpark.}
\item \textbf{Vérifie} : verbe en deuxième position, majuscule à tous les noms, forme correcte de l'impératif, article à l'accusatif (\all{der} $\rightarrow$ \all{den}).
\item \textbf{Compte} tes mots : entre 80 et 100 mots.
\end{enumerate}
\end{methode}

\begin{attention}[Erreurs fréquentes des francophones]
\begin{itemize}
\item Ne dis pas \all{Pflanzen Bäume!} pour t'adresser à des camarades : à \all{ihr}, la forme correcte est \all{Pflanzt Bäume!} (sans \all{Sie}).
\item Après \all{durch, für, gegen, ohne, um}, le masculin passe à l'accusatif : \all{durch den Wald}, et non \all{durch der Wald}.
\item Faux ami : \all{die Abholzung} = la déforestation ; \all{die Aufforstung} = le reboisement. Ne les confonds pas !
\item \all{helfen} se construit avec le datif : \all{Hilf den Dorfbewohnern!} (et non l'accusatif).
\item N'oublie pas la majuscule à tous les noms : \all{der Baum, die Umwelt, der Naturpark}.
\end{itemize}
\end{attention}

\newpage

\coursec{Corrigés des exercices}

\begin{corrige}[Exercice 1 — QCM : le texte support]
1. b — le texte traite de la protection des forêts au Cameroun. \quad 2. b — \all{die Aufforstung} = le reboisement (attention au faux ami : la déforestation se dit \all{die Abholzung}). \quad 3. b — l'écotourisme respecte la nature et aide la population locale. \quad 4. b — les villageois plantent des arbres et guident les touristes.
\end{corrige}

\begin{corrige}[Exercice 2 — Vrai ou faux ? Justifie]
\begin{enumerate}
\item \all{Falsch. Die Wälder in Kamerun sind in Gefahr.} (Les forêts du Cameroun sont en danger.)
\item \all{Richtig. Viele Dorfbewohner pflanzen neue Bäume.}
\item \all{Falsch. Der Ökotourismus bringt den Dörfern Geld.} (L'écotourisme rapporte de l'argent aux villages.)
\item \all{Richtig. Die Touristen können im Naturpark Tiere beobachten.}
\end{enumerate}
Points de vigilance : la justification doit être une phrase complète avec le verbe en deuxième position ; \all{in Gefahr sein} = être en danger ; \all{bringen} + datif (\all{den Dörfern}) + accusatif (\all{Geld}).
\end{corrige}

\begin{corrige}[Exercice 3 — Vocabulaire : trouve le mot]
1. \all{Im Wald wachsen viele Bäume.} (\all{im} = \all{in dem}, datif de lieu.) \quad 2. \all{Wir müssen die Umwelt schützen.} \quad 3. \all{Die Schüler pflanzen jedes Jahr zwanzig Bäume.} \quad 4. \all{Ein Baum hat grüne Blätter und eine lange Wurzel.} \quad 5. \all{Im Naturpark können Touristen Elefanten und Vögel sehen.}

Points de vigilance : après \all{müssen}, l'infinitif \all{schützen} va en fin de phrase ; \all{die Umwelt} est féminin singulier, l'article ne change pas à l'accusatif ; majuscule à tous les noms.
\end{corrige}

\begin{corrige}[Exercice 4 — Conjugaison : l'impératif]
\begin{center}
\begin{tabular}{|l|l|l|l|}
\hline
\rowcolor{popblueL} Infinitif & du & ihr & Sie \\ \hline
pflanzen & Pflanze \ldots ! & Pflanzt \ldots ! & Pflanzen Sie \ldots ! \\ \hline
helfen & Hilf \ldots ! & Helft \ldots ! & Helfen Sie \ldots ! \\ \hline
lesen & Lies \ldots ! & Lest \ldots ! & Lesen Sie \ldots ! \\ \hline
fahren & Fahr(e) \ldots ! & Fahrt \ldots ! & Fahren Sie \ldots ! \\ \hline
sein & Sei \ldots ! & Seid \ldots ! & Seien Sie \ldots ! \\ \hline
\end{tabular}
\end{center}
Rappels : à \all{du}, on prend le radical de la 2\up{e} personne du présent sans \all{-st}, avec la mutation vocalique éventuelle (\all{helfen} $\rightarrow$ \all{Hilf!}, \all{lesen} $\rightarrow$ \all{Lies!}, \all{fahren} $\rightarrow$ \all{Fahr(e)!}). À \all{ihr}, on reprend la forme du présent sans le pronom, jamais de mutation (\all{Pflanzt!}, \all{Helft!}, \all{Lest!}). À \all{Sie}, infinitif + \all{Sie} (\all{Fahren Sie!}). Le verbe \all{sein} est irrégulier : \all{Sei! Seid! Seien Sie!}
\end{corrige}

\begin{corrige}[Exercice 5 — Mets les phrases à l'impératif]
\begin{enumerate}
\item \all{Hilf den Dorfbewohnern!} (mutation e $\rightarrow$ i ; \all{helfen} + datif : \all{den Dorfbewohnern}.)
\item \all{Lest den Text über den Naturpark!} (pas de mutation à \all{ihr} ; \all{den Text} reste à l'accusatif.)
\item \all{Fahren Sie langsam durch den Wald!} (\all{durch} + accusatif : \all{den Wald}.)
\item \all{Nimm eine Tasche für den Müll!} (\all{nehmen} $\rightarrow$ \all{du nimmst} $\rightarrow$ \all{Nimm!} : chute du -e- et gémination du m ; \all{für} + accusatif.)
\item \all{Schützt die Umwelt!} (verbe régulier, pas de mutation.)
\item \all{Bleiben Sie auf dem Weg!} (\all{auf dem Weg} : datif de lieu, pas de changement.)
\end{enumerate}
\end{corrige}

\begin{corrige}[Exercice 6 — Prépositions + accusatif]
\begin{enumerate}
\item \all{Wir wandern durch den Wald.} (\all{der Wald} $\rightarrow$ accusatif masculin \all{den}.)
\item \all{Sie kämpft gegen die Abholzung.} (féminin : \all{die} ne change pas à l'accusatif.)
\item \all{Die Schüler sammeln Geld für den Naturpark.} (masculin : \all{der} $\rightarrow$ \all{den}.)
\item \all{Er geht ohne einen Freund in den Wald.} (\all{ein} $\rightarrow$ \all{einen} ; on accepte aussi \all{ohne seinen Freund}.)
\item \all{Die Gruppe sitzt um den Baum.} (masculin : \all{den Baum}.)
\end{enumerate}
À retenir : \all{durch, für, gegen, ohne, um} sont toujours suivies de l'accusatif. Seul le masculin change de forme : \all{der} $\rightarrow$ \all{den}, \all{ein} $\rightarrow$ \all{einen}.
\end{corrige}

\begin{corrige}[Exercice 7 — Texte à trous]
1. \all{Pflanzt} (impératif \all{ihr}, majuscule en début de phrase) ; 2. \all{für} (un projet \emph{pour} la nature : \all{für} + accusatif, \all{die Natur} ne change pas) ; 3. \all{durch} (\all{durch den Schulgarten} : à travers le jardin de l'école) ; 4. \all{Bäume} (accusatif pluriel : \all{zehn Bäume}) ; 5. \all{schützen} (sujet \all{wir} : \all{schützen wir}, verbe en deuxième position après \all{So}) ; 6. \all{Umwelt} (\all{die Umwelt}, sujet de \all{ist}).

Texte complet : \all{Pflanzt Bäume für die Zukunft! In unserer Schule in Yaoundé machen wir ein Projekt für die Natur. Wir gehen durch den Schulgarten und pflanzen zehn Bäume. So schützen wir die Erde und die Luft. Die Umwelt ist wichtig für alle Menschen in Kamerun.}
\end{corrige}

\begin{corrige}[Exercice 8 — Appariements]
1 -- d (\all{der Wald, die Wälder} = la forêt) ; 2 -- f (\all{der Baum, die Bäume} = l'arbre) ; 3 -- a (\all{pflanzen} = planter) ; 4 -- b (\all{schützen} = protéger) ; 5 -- c (\all{die Umwelt} = l'environnement) ; 6 -- e (\all{der Naturpark, die Naturparks} = le parc naturel).
\end{corrige}

\begin{corrige}[Exercice 9 — Compréhension écrite type Bac]
\textbf{Fragen zum Text :}
\begin{enumerate}
\item \all{Der Naturpark Schwarzwald liegt im Südwesten Deutschlands.}
\item \all{Die Touristen wandern durch den Wald und beobachten Vögel.} (ou : \all{Sie schlafen in kleinen Hotels.})
\item \all{Die Mitarbeiter sagen: „Bleiben Sie auf den Wegen!“ und „Nehmen Sie Ihren Müll mit!“}
\item \all{Der Ökotourismus ist gut für die Dörfer, weil sie Geld verdienen.} (attention : après \all{weil}, le verbe va en fin de phrase.)
\item \all{Die Touristen lernen viel über die Natur.}
\end{enumerate}
\textbf{Richtig oder Falsch ?}
\begin{enumerate}
\item \all{Falsch.} Le texte dit : \all{„Sie schlafen in kleinen Hotels.“}
\item \all{Richtig.} Le texte dit : \all{„Nehmen Sie Ihren Müll mit!“}
\item \all{Falsch.} Le texte dit : \all{„Die Dörfer in der Region verdienen Geld mit dem Ökotourismus.“}
\end{enumerate}
\textbf{Fragen zur Sprache :}
\begin{enumerate}
\item \all{Bleiben Sie!}, \all{Nehmen Sie Ihren Müll mit!}, \all{Pflanzen Sie einen Baum!} : impératifs à la forme de politesse \all{Sie} (le pronom \all{Sie} est placé après le verbe).
\item \all{durch den Wald} (\all{durch} + accusatif) ; \all{für die Zukunft} (\all{für} + accusatif).
\end{enumerate}
Barème indicatif (sur 10) : Fragen zum Text 5 pts (1 pt par réponse complète et correcte) ; Richtig/Falsch 3 pts (1 pt par réponse justifiée) ; Fragen zur Sprache 2 pts. Points de vigilance : répondre par des phrases complètes, verbe en deuxième position, ne pas recopier le texte sans l'adapter à la question.
\end{corrige}

\begin{corrige}[Exercice 10 — Thème et version]
\textbf{A. Thème.}
\begin{enumerate}
\item \all{Pflanzen Sie Bäume für die Zukunft!} (impératif \all{Sie} ; \all{für} + accusatif.)
\item \all{Wir gehen durch den Wald ohne Lärm.} (\all{durch} + accusatif : \all{den Wald} ; \all{ohne} + accusatif : \all{ohne Lärm}.)
\item \all{Hilf den Dorfbewohnern, mein Freund!} (impératif \all{du} avec mutation e $\rightarrow$ i ; \all{helfen} + datif : \all{den Dorfbewohnern}.)
\item \all{Sie kämpft gegen die Abholzung in Kamerun.} (\all{gegen} + accusatif : \all{die Abholzung}.)
\end{enumerate}
\textbf{B. Version.} « Venez au parc naturel et découvrez la nature ! Randonnez à travers la forêt verte et observez les animaux. Beaucoup de familles vont au parc le week-end. Les enfants y apprennent beaucoup de choses sur les arbres et les oiseaux. Le parc lutte contre la déforestation et plante chaque année de nouveaux arbres. Visitez-nous — la nature vous attend ! »

Points de vigilance : \all{wandern} = randonner (et non « errer ») ; \all{am Wochenende} = le week-end ; \all{warten auf} + accusatif = attendre quelqu'un.
\end{corrige}

\begin{corrige}[Exercice 11 — Production écrite et jeu de rôle]
\textbf{A. Proposition de rédaction modèle (environ 90 mots) :}

\all{Die Umwelt ist wichtig für uns alle. In meiner Stadt Yaoundé können wir viel tun. Pflanzt Bäume in eurem Viertel! Schützt die Tiere und den Wald! Werft keinen Müll auf die Straße! Wir kämpfen gegen die Abholzung, denn ohne Bäume gibt es keine saubere Luft. Unsere Schule sammelt Geld für einen Naturpark. Wir gehen oft durch den Schulgarten und pflanzen junge Bäume. Helft den Dorfbewohnern und lernt viel über die Natur! So schützen wir gemeinsam die Umwelt für die Zukunft.}

Vérification des contraintes : impératifs \all{Pflanzt!, Schützt!, Werft!, Helft!, lernt!} ; prépositions + accusatif : \all{für uns alle, gegen die Abholzung, ohne Bäume, für einen Naturpark, durch den Schulgarten, für die Zukunft} ; vocabulaire du chapitre : \all{der Wald, der Baum, pflanzen, schützen, die Umwelt, der Naturpark}.

Barème indicatif (sur 10) : respect des contraintes (3 impératifs, 2 prépositions + accusatif, 5 mots du thème) 4 pts ; correction grammaticale 3 pts ; vocabulaire 2 pts ; cohérence et longueur (80--100 mots) 1 pt.

Points de vigilance : à \all{ihr}, pas de \all{Sie} après le verbe (\all{Pflanzt!}, non \all{Pflanzen!}) ; \all{helfen} + datif ; majuscule à tous les noms ; verbe en deuxième position.

\textbf{B. Jeu de rôle — exemple de dialogue :}

\all{Guide: Guten Tag! Willkommen im Naturpark. Bleiben Sie bitte auf den Wegen! Fotografieren Sie die Tiere, aber kommen Sie nicht zu nah! Werfen Sie keinen Müll in den Wald! Nehmen Sie Wasser mit! Und pflanzen Sie am Ende einen Baum für die Zukunft!}

\all{Tourist: Danke! Können wir Elefanten sehen? Wann beginnt die Tour? Und was kostet der Besuch?}

\all{Guide: Ja, wir sehen oft Elefanten. Die Tour beginnt um neun Uhr und kostet zehn Euro.}
\end{corrige}

\newpage

\coursec{Fiche bilan}

\begin{popfigure}
\begin{center}\tcbox[colback=popgreen,colframe=popgreen,coltext=white,arc=9pt,boxsep=4pt,left=6pt,right=6pt]{\bfseries\small Reboisement et écotourisme}\end{center}
\vspace{-2pt}
\begin{tcbraster}[raster columns=3,raster equal height=rows,raster column skip=6pt,raster row skip=6pt,enhanced,blankest]
\begin{tcolorbox}[enhanced,colback=white,colframe=popblue,boxrule=0.9pt,arc=3pt,title={Texte support},fonttitle=\bfseries\scriptsize,coltitle=white,colbacktitle=popblue,left=4pt,right=4pt,top=2pt,bottom=2pt,boxsep=1pt,toptitle=1.5pt,bottomtitle=1.5pt,halign=left]
\begin{itemize}[nosep,leftmargin=*,itemsep=1pt,font=\scriptsize]
\item Aufforstung in Kamerun
\item compréhension globale + détaillée
\end{itemize}
\end{tcolorbox}
\begin{tcolorbox}[enhanced,colback=white,colframe=poporange,boxrule=0.9pt,arc=3pt,title={Lexique},fonttitle=\bfseries\scriptsize,coltitle=white,colbacktitle=poporange,left=4pt,right=4pt,top=2pt,bottom=2pt,boxsep=1pt,toptitle=1.5pt,bottomtitle=1.5pt,halign=left]
\begin{itemize}[nosep,leftmargin=*,itemsep=1pt,font=\scriptsize]
\item der Wald, der Baum
\item pflanzen, schützen
\end{itemize}
\end{tcolorbox}
\begin{tcolorbox}[enhanced,colback=white,colframe=poppurple,boxrule=0.9pt,arc=3pt,title={Impératif},fonttitle=\bfseries\scriptsize,coltitle=white,colbacktitle=poppurple,left=4pt,right=4pt,top=2pt,bottom=2pt,boxsep=1pt,toptitle=1.5pt,bottomtitle=1.5pt,halign=left]
\begin{itemize}[nosep,leftmargin=*,itemsep=1pt,font=\scriptsize]
\item du / ihr / Sie
\end{itemize}
\end{tcolorbox}
\begin{tcolorbox}[enhanced,colback=white,colframe=popgold,boxrule=0.9pt,arc=3pt,title={Prépositions + Akkusativ},fonttitle=\bfseries\scriptsize,coltitle=white,colbacktitle=popgold,left=4pt,right=4pt,top=2pt,bottom=2pt,boxsep=1pt,toptitle=1.5pt,bottomtitle=1.5pt,halign=left]
\begin{itemize}[nosep,leftmargin=*,itemsep=1pt,font=\scriptsize]
\item durch, für, gegen, ohne, um
\end{itemize}
\end{tcolorbox}
\begin{tcolorbox}[enhanced,colback=white,colframe=poppink,boxrule=0.9pt,arc=3pt,title={Méthode},fonttitle=\bfseries\scriptsize,coltitle=white,colbacktitle=poppink,left=4pt,right=4pt,top=2pt,bottom=2pt,boxsep=1pt,toptitle=1.5pt,bottomtitle=1.5pt,halign=left]
\begin{itemize}[nosep,leftmargin=*,itemsep=1pt,font=\scriptsize]
\item questions, résumé
\item court paragraphe
\end{itemize}
\end{tcolorbox}
\begin{tcolorbox}[enhanced,colback=white,colframe=popgreen!70,boxrule=0.9pt,arc=3pt,title={Culture},fonttitle=\bfseries\scriptsize,coltitle=white,colbacktitle=popgreen!70,left=4pt,right=4pt,top=2pt,bottom=2pt,boxsep=1pt,toptitle=1.5pt,bottomtitle=1.5pt,halign=left]
\begin{itemize}[nosep,leftmargin=*,itemsep=1pt,font=\scriptsize]
\item Naturparks, Ökotourismus
\end{itemize}
\end{tcolorbox}
\end{tcbraster}

\legende{Carte mentale de la leçon : texte, lexique, grammaire et méthode}
\end{popfigure}

\begin{aretenir}[L'essentiel de la leçon]
\textbf{1. Lexique clé (à connaître avec article et pluriel)}
\begin{itemize}
  \item \all{der Wald, die Wälder} : la forêt ; \all{der Baum, die Bäume} : l'arbre ; \all{die Pflanze, -n} : la plante.
  \item \all{pflanzen} : planter ; \all{schützen} : protéger ; \all{retten} : sauver ; \all{abholzen} (particule séparable) : abattre, déboiser.
  \item \all{die Umwelt} (sans pluriel) : l'environnement ; \all{der Umweltschutz} (sans pluriel) : la protection de l'environnement.
  \item \all{der Naturpark, -s} : le parc naturel ; \all{der Ökotourismus} (sans pluriel) : l'écotourisme ; \all{die Aufforstung, -en} : le reboisement.
  \item \all{der Tourist, -en / die Touristin, -nen} : le touriste / la touriste ; \all{die Natur, -en} : la nature ; \all{sauber} : propre.
\end{itemize}

\textbf{2. Grammaire à connaître par cœur}
\begin{center}
\begin{tabularx}{\linewidth}{|l|X|X|}
\hline
\rowcolor{popblueL} \textbf{Point} & \textbf{Règle} & \textbf{Exemple} \\
\hline
Impératif (du) & radical du verbe (2\textsuperscript{e} pers. sg. sans \emph{-st}), \emph{-e} facultatif & \all{Pflanz(e) einen Baum!} \\
\hline
Impératif (ihr) & forme \emph{ihr} sans le pronom & \all{Schützt den Wald!} \\
\hline
Impératif (Sie) & forme \emph{Sie} + pronom placé après le verbe & \all{Besuchen Sie den Naturpark!} \\
\hline
Impératif de \all{sein} & \all{Sei! Seid! Seien Sie!} & \all{Seid vorsichtig!} \\
\hline
Prépositions + Akkusativ & \all{durch, für, gegen, ohne, um} (+ \all{entlang}, postposée) & \all{Wir gehen durch den Wald.} \\
\hline
\end{tabularx}
\end{center}

\textbf{3. Points de vigilance}
\begin{itemize}
  \item À l'impératif, le verbe est en \cle{première position} : \all{Pflanze Bäume!} (jamais \all{Du pflanze}).
  \item Après \all{durch, für, gegen, ohne, um}, toujours l'\cle{accusatif} : \all{für die Umwelt}, \all{ohne den Wald}.
  \item Verbes à particule séparable : \all{abholzen} se sépare au Präsens : \all{Man holzt die Bäume ab.}
  \item Majuscule à tous les noms : \all{der Wald, die Umwelt, der Naturpark}.
\end{itemize}

\textbf{4. Méthode express}
\begin{enumerate}
  \item Lire le texte une première fois pour l'idée générale (qui ? où ? quoi ?).
  \item Souligner les mots du thème (forêt, protection, tourisme) et repérer les mots transparents (\all{Naturpark, Tourist, Ökotourismus}).
  \item Répondre aux questions par des phrases complètes en reprenant les mots du texte.
  \item Pour le résumé : 3 à 5 phrases au Präsens, sans recopier le texte.
  \item Pour la production : 5 à 8 phrases avec des impératifs et des prépositions + accusatif.
\end{enumerate}
\end{aretenir}

\begin{exemplebox}[Modèles de phrases à réutiliser à l'écrit comme à l'oral]
\all{Pflanzt Bäume für die Zukunft!} « Plantez des arbres pour l'avenir ! »

\all{Wir müssen die Umwelt schützen.} « Nous devons protéger l'environnement. »

\all{Die Touristen wandern durch den Naturpark.} « Les touristes font une randonnée à travers le parc naturel. »

\all{Ohne den Wald gibt es kein Leben.} « Sans la forêt, il n'y a pas de vie. »

\all{Seien Sie freundlich zur Natur!} « Soyez bienveillant envers la nature ! »
\end{exemplebox}

\begin{attention}[Erreurs fréquentes des francophones]
\begin{itemize}
  \item Ne pas confondre \all{der Baum} (l'arbre) et \all{der Wald} (la forêt) : le genre change, l'article aussi.
  \item \all{die Umwelt} ne prend jamais de pluriel et s'emploie presque toujours avec l'article défini.
  \item À l'impératif \emph{Sie}, le pronom \all{Sie} est obligatoire et se place après le verbe : \all{Kommen Sie!} (et non \all{Kommen!} pour vouvoyer).
  \item Attention aux verbes qui changent de voyelle : \all{lesen} donne \all{Lies!}, \all{nehmen} donne \all{Nimm!} ; mais les verbes réguliers comme \all{pflanzen} gardent leur radical : \all{Pflanze!}.
  \item \all{entlang} se place \emph{après} le nom : \all{Wir gehen den Fluss entlang.} « Nous marchons le long de la rivière. »
\end{itemize}
\end{attention}

\begin{aretenir}[Checklist avant le Bac]
\begin{itemize}
  \item $\square$ Je connais l'article et le pluriel de \all{der Wald}, \all{der Baum}, \all{der Naturpark}, \all{die Umwelt}.
  \item $\square$ Je sais conjuguer \all{pflanzen}, \all{schützen}, \all{schneiden} au Präsens.
  \item $\square$ Je sais former les trois impératifs : \all{Pflanz(e)!}, \all{Pflanzt!}, \all{Pflanzen Sie!}.
  \item $\square$ Je connais l'impératif irrégulier de \all{sein} : \all{sei / seid / seien Sie}.
  \item $\square$ Je cite par cœur les prépositions + accusatif : \all{durch, für, gegen, ohne, um}.
  \item $\square$ Je mets toujours l'accusatif après ces prépositions : \all{für den Wald}.
  \item $\square$ Je place le verbe en première position à l'impératif.
  \item $\square$ Je mets la majuscule à tous les noms allemands.
  \item $\square$ Je sais séparer les verbes à particule : \all{Man holzt Bäume ab.}
  \item $\square$ Je sais résumer un texte en 3 à 5 phrases au Präsens.
  \item $\square$ Je sais écrire un court paragraphe sur la protection de l'environnement au Cameroun.
  \item $\square$ Je relis ma production : place du verbe, cas, orthographe des Umlaute.
\end{itemize}
\end{aretenir}

\findecours

\end{document}
